\chapter{SLAM 与稠密建图}

\paragraph*{本章阅读锚点}
把 3DGS 放进 SLAM 后，问题不再只是离线优化一个可渲染场景，而是相机位姿、地图几何、渲染外观和历史关键帧要在同一个在线系统里相互牵制。原始 3DGS 的锚点是显式高斯、可微 splatting 和按残差驱动的密度控制；SLAM 的锚点则是跟踪、建图、回环、局部或全局 bundle adjustment。二者结合后，核心问题可以拆成几条路线：RGB-only 方法要用单目深度、光流或直接法补足尺度和几何；RGB-D 方法把深度、轮廓和不确定性接入高斯优化；多传感器方法用 LiDAR、IMU、事件或雷达稳定位姿和稀疏几何；定位方法把高斯地图从“渲染资产”改造成可匹配的 3D 地标；语义和动态方法则要求地图知道哪些高斯属于物体、类别或动态干扰。阅读本章时，重点看每篇方法把原始 3DGS 或已有 SLAM 系统的哪一个环节替换掉：初始化、跟踪损失、地图更新、密度控制、回环约束、传感器融合，还是最终的地图维护策略。

\section{单目与 RGB SLAM}

RGB-only 3DGS SLAM 的难点在于，原始 3DGS 默认相机位姿和稀疏点云已经较可靠，而单目 SLAM 同时缺少绝对尺度和稠密深度。早期方法常把 ORB-SLAM、DROID-SLAM 或直接法视觉里程计作为前端，再让 3DGS 负责高保真建图；后续方法开始让高斯地图反过来参与跟踪、尺度校正、关键帧选择和历史记忆压缩。这一节的主线不是“是否用了 3DGS”，而是看高斯在系统中承担的是地图、优化目标、几何先验，还是历史状态压缩器。



% METHOD: 2311.16728
\paragraph{Photo-SLAM（特征定位与照片级高斯地图分工）（2023.11）}
Photo-SLAM~\cite{2311_16728} 的锚点是传统特征 SLAM 与 3DGS 的分工。它没有把高斯强行承担所有定位职责，而是保留 ORB 特征、因子图和多传感器前端来稳定相机轨迹，再用 3D Gaussian hyper primitives 生成照片级地图。这样做的原因是，原始 3DGS 渲染质量高，但在在线 SLAM 中直接从光度残差恢复鲁棒位姿并不总是可靠；传统特征前端在快速运动、立体或 RGB-D 场景里仍有工程优势。

系统支持 monocular、stereo 和 RGB-D。跟踪侧使用 ORB 特征进行定位和回环，建图侧把关键帧和稀疏几何转成高斯，并通过 Gaussian pyramid progressive training 逐步提高渲染质量。它还用主动 densification 根据几何特征决定哪里应增加高斯，而不是完全依赖图像残差。相对 Gaussian Splatting SLAM 的直接法追踪，Photo-SLAM 的改变是把定位可靠性留给成熟 SLAM 前端，把高斯作为照片级地图层；因此它更像一个工程锚点，说明 3DGS 可以嵌入现有 SLAM 栈而不必替代所有模块。
\textbf{实现细节：} 与 SplaTAM 将跟踪与建图统一在同一可微渲染管线中的范式不同，Photo-SLAM（2023.11） 采用了跟踪与建图解耦的架构。因此 Photo-SLAM 将定位任务完全交给 ORB 特征前端的因子图优化（用 Levenberg--Marquardt 算法最小化重投影误差），将照片级渲染交给后端 3DGS 优化，两者通过超基元地图（hyper primitives map）桥接。超基元地图与 SplaTAM 简化高斯的差异。与 SplaTAM 将高斯参数压缩为 8 维各向同性球状高斯不同，Photo-SLAM 的超基元保留了完整 3DGS 表达能力。每个超基元存储位置 $\mathbf{P} \in \mathbb{R}^{}3$、旋转 $\mathbf{r} \in SO(3)$、缩放 $\mathbf{s} \in \mathbb{R}^{}3$、密度 $\alpha \in \mathbb{R}^{}1$、球谐系数 $\mathrm{SH} \in \mathbb{R}^{16}$（对应 0--1 阶球谐，可建模视角相关颜色），以及 ORB 描述子 $\mathbf{O} \in \mathbb{R}^{256}$。换言之，Photo-SLAM 不需要简化高斯表示，因为定位任务已由前端 ORB 模块承担，高斯参数只需优化渲染质量。渲染公式与标准 3DGS 的 $\alpha$-混合一致：$C(R,t) = \sum_{i} c_i \alpha_i \prod_{j=1}^{i-1}(1-\alpha_j)$。几何致密化：利用 ORB 特征引导高斯增生。SplaTAM 用渲染轮廓和深度不一致性驱动在线致密化，而 Photo-SLAM 利用前端已有的 ORB 特征实现一种互补的致密化策略。Photo-SLAM 在每个关键帧创建时，主动将这些非活跃 ORB 特征点转化为临时超基元：在 RGB-D 模式下直接用深度值投影；在单目模式下用最近邻活跃特征点的深度估计；在双目模式下用立体匹配估计深度。高斯金字塔训练。Photo-SLAM 进一步引入高斯金字塔（Gaussian-Pyramid）渐进训练方法：默认使用 3 级金字塔（$n=2$），训练从最高层级（最粗分辨率）开始，随迭代增加逐级降至原始分辨率。
\textbf{方法信息：} 引用很多，Semantic Scholar 约 228 次。
\textbf{设计动机与边界：} 相对原始 3DGS 以独立高斯和图像重建误差为主的设定，Photo-SLAM 把问题落在以下可实现的改动上：Photo-SLAM的核心是超原语图，包含点云存储的ORB特征（用于定位）和旋转、缩放、密度、球谐系数（用于渲染）。系统分两个线程：跟踪线程利用ORB特征通过因子图求解器优化相机位姿；建图线程通过3D-GS的可微分渲染反向传播损失优化超原语。关键创新是高斯金字塔训练方法：渐进式地从低分辨率到高分辨率学习多级特征，先学习场景的大尺度结构再细化细节，避免单目设置中因缺乏深度信息导致的优化困难。主动致密化策略基于几何特征在信息不足区域增加超原语。其设计目的不是单独增加一个模块，而是用上述表示、监督或渲染约束消除原流程在目标任务中的具体失配。提出Photo-SLAM框架，创新性地将3D Gaussian Splatting集成到SLAM系统中。维护超原语图同时利用ORB特征进行定位和3D-GS进行逼真建图。引入基于高斯金字塔的渐进式训练方法学习多级特征。支持单目、双目和RGB-D输入，在Replica数据集上PSNR比SOTA高30%，渲染速度快数百倍，可在Jetson AGX Orin上实时运行。适用时还应保留论文给出的边界：单目模式下深度初始化不如RGB-D精确；对快速运动可能丢失跟踪；大场景的内存管理待优化。
% METHOD: 2312.06741
\paragraph{Gaussian Splatting SLAM（解析位姿雅可比的早期锚点）（2023.12）}
Gaussian Splatting SLAM~\cite{2312_06741} 是把 3DGS 直接用作在线视觉 SLAM 表示的早期工作之一。它的关键问题是：如果高斯既是地图又是渲染器，如何在没有神经场长时间反向传播的情况下直接优化相机位姿。方法的改变不是引入复杂前端，而是为 EWA splatting 推导关于 $SE(3)$ 位姿的解析雅可比，使帧到模型追踪能直接在高斯渲染图上求解。

为了在线运行，它还简化了原始 3DGS 表示：颜色使用 view-independent 形式，并省略球谐函数的视角相关外观；跟踪时通过相机位姿梯度直接最小化渲染误差。地图更新不是盲目 densification，而是结合几何验证和尺度正则，检查新高斯是否被已有视角充分约束，并抑制退化成极端扁平形状的高斯。系统支持 monocular、stereo 和 RGB-D 输入，在线输出相机轨迹和持续更新的高斯地图；单目模式需要处理尺度对齐，RGB-D 模式则可直接利用深度稳定几何。代码 README 也特别说明，仓库名 MonoGS 只是简称，论文引用时应称为 Gaussian Splatting SLAM，并提供带相机位姿梯度的 differential Gaussian rasterizer。相对后续 Splat-SLAM 或 DROID-Splat，它没有复杂全局 BA 和学习式前端，但它提供了一个重要局部锚点：高斯渲染本身可以成为直接法位姿优化目标，而不是只在建图结束后用于可视化。
\textbf{实现细节：} Gaussian Splatting SLAM（MonoGS）（2023.12） 是首个将 3DGS 应用于单目 SLAM 的工作。动机。在此之前的 NeRF-SLAM 方法需要逐射线采样和体渲染，单帧渲染耗时数百毫秒，无法实时运行；而标准 3DGS 虽然渲染速度快，但其自适应致密化机制和位姿假设是为离线新视角合成设计的，不适合帧序到达的在线 SLAM 场景。本方法修改了 3DGS 的相机位姿优化和建图环节——用解析雅可比矩阵替代自动微分进行位姿梯度计算，用各向同性正则化约束高斯形状，用基于共视性的关键帧管理替代固定间隔策略——3DGS 的光栅化渲染管线和高斯基元表示保持不变。其背景挑战在于：没有深度传感器时，高斯基元无法从深度图直接初始化位置，且位姿优化的收敛域大幅缩小。跟踪与解析雅可比。标准 3DGS 的 CUDA 光栅化仅对高斯参数（位置、颜色、不透明度、协方差）计算了显式梯度，但并未推导相机位姿 $\mathbf{T} \in \mathrm{SE}(3)$ 关于渲染结果的梯度。在基线做法中，若要端到端地通过可微渲染优化相机位姿，只能依赖自动微分（autodiff）逐层回传梯度，这在每帧需要多次迭代的 SLAM 场景下引入了显著的计算开销。MonoGS 的核心创新是推导了 3DGS 中相机位姿的解析雅可比矩阵——这是首次在 EWA（椭圆加权平均）Splatting 框架下给出 $\mathrm{SE}(3)$ 位姿关于高斯参数的完整解析梯度。直观地理解，3D 高斯的世界坐标 $\boldsymbol{\mu}^w$ 通过相机位姿 $\mathbf{T}$ 投影到图像平面得到相机坐标 $\boldsymbol{\mu}^c$，而这个投影过程关于位姿的导数决定了渲染误差如何反向传播到位姿参数。
\textbf{锚点细化：} MonoGS：是首个单目 3DGS-SLAM 工作，针对单目无深度传感器、高斯无法直接初始化、位姿优化收敛域大幅缩小的局限（单目还有一个固有难题——\textbf{尺度歧义}：单目相机无法区分"近处的小物体"和"远处的大物体"，因为两者在图像上投影完全一样），用\textbf{解析雅可比 + 单目深度方差初始化}替代自动微分位姿计算，核心思想是\textbf{在 EWA splatting 框架下推导 SE(3) 位姿完整解析梯度}——解析雅可比直接给出位姿关于渲染损失的闭式梯度，相比自动微分逐层回传不仅省去每帧多次迭代的额外开销，还能用李代数最小参数化保证优化始终在 SE(3) 流形上进行。基于高斯共视性（IOU 与重叠系数）管理关键帧，对有渲染深度的像素低方差初始化、无深度像素高方差初始化使多数错误高斯自然消失，各向同性正则化抑制稀疏视角下针状高斯退化。
\textbf{方法信息：} 代码：\url{https://github.com/muskie82/MonoGS}；GitHub 关注度很高，约 2100 stars；引用很多，Semantic Scholar 约 597 次。



% METHOD: 2405.11129
\paragraph{MotionGS（运动过滤驱动的紧凑 SLAM）（2024.05）}
MotionGS~\cite{2405_11129} 继承的是 MonoGS 一类在线 Gaussian SLAM 的路线，但它不满足于把每个关键帧都转成新高斯并持续 densification。它观察到 SLAM 输入中大量帧只带来很小运动，若照搬原始 3DGS 的残差驱动增密，会产生许多对位姿和渲染贡献有限的小高斯。方法的改变是把“是否需要记住这一帧”交给 motion filter：跟踪阶段先提取深度特征并根据相机运动选择运动关键帧，再用直接位姿优化和信息滤波窗口更新局部状态。

建图阶段，MotionGS 在关键帧位姿和高斯参数上联合优化，同时给每个高斯引入一个可学习二值 mask 倾向。mask 由尺度、不透明度和辅助变量共同决定，通过 sigmoid 和 stop-gradient 训练，让小尺度、低贡献、重复解释的 Mini-Gaussians 在密度控制时被抑制或移除。项目仓库材料与 MotionGS 项目一致，说明代码围绕 motion filtering 和 compact Gaussian map 展开。相对 SplaTAM 的 RGB-D 稠密增量建图，MotionGS 更像在 RGB-only 场景中把关键帧管理和高斯剪枝合到一起，目标是让在线地图更小、更稳定。
\textbf{实现细节：} MotionGS（2024.05） 同样基于可微渲染位姿优化，但不同于 S3PO-GS 依赖预训练点图建立 2D-3D 对应，MotionGS 从跟踪管线的帧选择机制切入，提出双关键帧策略以减少冗余帧参与优化。动机。标准 3DGS-SLAM 对所有帧统一处理，但大量帧之间的视角变化极小，参与联合优化时仅带来冗余梯度而无新增信息，不仅浪费计算还可能导致过拟合。双关键帧策略分为两层：运动滤波器提取每帧的密集视觉特征（6 层残差模块 + 3 层下采样，输出 256 维 1/8 分辨率特征图），通过 4 层相关性金字塔计算帧间运动向量，当运动向量范数超过阈值或帧间距超过上限时选为运动关键帧；信息滤波器则进一步从运动关键帧中筛选——当前帧相对前一信息关键帧的场景几何相对补集 $\mathrm{RC}(i,j) = (G_i - G_i \cap G_j)/(G_i \cup G_j)$ 超过阈值时，才加入滑动窗口成为信息关键帧，窗口满时按重叠系数 $\mathrm{OC}_{i,j}$ 淘汰最冗余帧。这种两级筛选使参与联合优化的帧数大幅减少，同时保持几何覆盖的完整性。在场景表示方面，MotionGS 借鉴 Compact 3DGS 的 mask 思路，为每个高斯的尺度 $s_i$ 和不透明度 $o_i$ 引入可学习二值掩码 $B_i = \sigma(b_i)$，在优化过程中自动抑制对渲染贡献极小的"微型高斯"，配合 mask 损失 $\mathcal{L}_M = \frac{1}{N}\sum_{i}\sigma(b_i)$ 鼓励稀疏化，优化结束后丢弃掩码。
\textbf{设计动机与边界：} 相对原始 3DGS 以独立高斯和图像重建误差为主的设定，MotionGS 把问题落在以下可实现的改动上：MotionGS包含两个核心线程：(1) 跟踪线程：对每帧图像提取深度特征，通过运动滤波器判断是否为运动关键帧（减少冗余帧处理），使用专为3DGS设计的直接位姿优化方法进行位姿估计，通过信息滤波器持续更新和维护信息关键帧的滑动窗口；(2) 建图线程：在diff-gaussian-rasterization框架下同时优化关键帧位姿和3D高斯参数。紧凑表示通过在优化过程中引入基于高斯尺度和不透明度的二值掩码实现：为每个高斯设置辅助参数b，通过sigmoid和停止梯度操作生成掩码，消除对渲染影响极小的Mini-Gaussian，减少参与优化和存储的高斯数量。设计新型掩码损失L\_M来平衡高保真渲染与减少场景几何。在每次densification中除split/clone/prune外，还基于掩码移除部分场景几何。其设计目的不是单独增加一个模块，而是用上述表示、监督或渲染约束消除原流程在目标任务中的具体失配。提出基于运动滤波的3DGS-SLAM方法，通过深度特征提取和运动滤波减少处理帧数，双重关键帧选择实现由粗到细的位姿估计，新型损失函数实现紧凑的高斯场景表示，在跟踪和建图性能上超越现有方法且内存使用更少适用时还应保留论文给出的边界：目前仅支持RGB-D输入，未扩展到单目或立体相机设置；深度特征提取增加了计算开销；在纹理缺失或重复纹理区域的表现有待改善
\textbf{方法信息：} 代码：\url{https://github.com/ruijiezhu94/motiongs}；GitHub 关注度较高，约 155 stars。



% METHOD: 2405.16544
\paragraph{Splat-SLAM（RGB-only 的全局优化锚点）（2024.05）}
Splat-SLAM~\cite{2405_16544} 是 RGB-only 3DGS SLAM 的重要锚点。它继承 DROID-SLAM 的 dense flow 和 bundle adjustment 思路，同时把 3DGS 地图作为可变形的全局表示。原始 3DGS 依赖给定相机和稀疏点云，普通 RGB-only SLAM 又难以产生足够高质量的可渲染稠密地图；Splat-SLAM 的改变是让 tracking、depth proxy、local BA、global BA 和 Gaussian deformation 共同工作。

具体流程中，前端用稠密光流估计帧间对应，并通过 DSPO 层同时优化 disparity、尺度和相机位姿。局部与全局 BA 更新关键帧位姿和深度后，系统不是重新建图，而是主动变形已有 3DGS 地图，使高斯中心、尺度和颜色跟随新的关键帧几何对齐。单目深度来自 DPT 之类的 proxy，并通过尺度和平移对齐进入优化。相对 Gaussian Splatting SLAM 的解析 pose Jacobian 路线，Splat-SLAM 更强调全局 BA 与地图形变；相对 DROID-Splat，它不是简单把 DROID 前端接到 renderer 后面，而是把高斯地图作为随全局优化不断修正的状态。
\textbf{实现细节：} Splat-SLAM（2024.05） 是首个实现全局优化的 RGB-only 高斯 SLAM 系统。动机。MonoGS 采用 frame-to-model 跟踪，没有全局位姿和地图优化，长期运行会产生不可逆的累积漂移；同时 MonoGS 仅依赖渲染深度进行高斯初始化，在多视角约束不足的区域容易产生几何退化。Splat-SLAM 的核心思路是将frame-to-frame 光流跟踪与可变形高斯地图结合，通过回环检测和全局光束法平差（BA）消除漂移，并通过在线地图变形将全局优化结果传播到高斯场景表示中。在建图环节，Splat-SLAM 沿用了 MonoGS 的高斯共视性关键帧选择、基于可见性的剪枝和增密策略以及各向同性正则化，但在跟踪和全局一致性上做出了根本性的架构改变。
\textbf{实现补充：} 实现上还要核对输入、输出、参数、特征、初始化、训练损失、正则、渲染或推理接口，以及优化中的剪枝、分裂、采样、编码或解码步骤；这些抓手决定该方法到底改了 3DGS 管线的哪一环。
\textbf{设计动机与边界：} 相对原始 3DGS 以独立高斯和图像重建误差为主的设定，Splat-SLAM 把问题落在以下可实现的改动上：Splat-SLAM是首个结合全局优化跟踪与3DGS稠密建图的仅RGB SLAM系统。核心组件：1）跟踪模块：使用预训练循环密集光流模型（RAFT）结合DSPO（视差、尺度、位姿优化）层联合优化相机位姿和逐像素视差。局部BA在滑动窗口内优化，全局BA在回环检测触发时执行扩展图优化。DSPO层包含两个交替优化的目标：密集BA（基于光流）和单目深度约束。2）3DGS地图变形：当关键帧位姿或深度被全局BA更新时，3D高斯地图主动变形以适应新参数。这是通过将高斯中心位置和深度与关键帧位姿绑定实现的——位姿更新时高斯自动跟随变形。3）代理深度图：引入单目深度估计器（DPT模型）提供相对深度先验，通过尺度和平移参数对齐到优化后的视差。代理深度图辅助高斯初始化和正则化。4）高斯管理：采用MonoGS的初始化和剪枝策略，基于共可见性进行关键帧选择和窗口管理。其设计目的不是单独增加一个模块，而是用上述表示、监督或渲染约束消除原流程在目标任务中的具体失配。提出首个利用全局优化跟踪的仅RGB 3DGS SLAM系统Splat-SLAM。通过密集光流进行帧到帧跟踪，实现回环检测和全局BA。关键创新是3D高斯地图能动态适应关键帧位姿和深度更新，主动变形以保持一致性。引入代理深度图结合单目深度估计器改进深度精度。在Replica、TUM-RGBD和ScanNet上达到SOTA或持平性能。适用时还应保留论文给出的边界：依赖光流模型的质量；单目深度估计精度有限；实时性受全局BA计算影响；对快速运动和纹理缺失区域可能跟踪失败。
\textbf{方法信息：} 代码：\url{https://github.com/eriksandstroem/splat-slam}；GitHub 关注度较高，约 107 stars；引用很多，Semantic Scholar 约 90 次。



% METHOD: 2405.17187
\paragraph{Memorize What Matters（多次 traverses 中只记住静态结构）（2024.05）}
Memorize What Matters~\cite{2405_17187} 更接近长期建图而非单次 SLAM，但它对 RGB 高斯地图很重要：同一路段被多次 traverse 采集时，车辆、行人、天气和光照都会改变，普通 3DGS 会把这些 transient 内容混入静态地图。方法的局部锚点是鲁棒特征蒸馏和场景分解，而不是前端位姿优化。

具体做法是把 DINOv2 特征蒸馏进 3DGS，并跨 traverses 比较特征残差，找出哪些 2D 区域在几何和语义上反复不一致。由这些 residual mining 得到 transient segmentation 后，训练阶段降低动态或短时内容的权重，使静态环境被更稳定地写入高斯地图。它相对 MotionGS 的关键帧压缩不同：MotionGS 选择“哪些帧值得记”，该方法选择“帧里哪些内容不应被记”。在自动驾驶或长期地图更新中，这类机制会影响高斯地图是否能作为可复用的环境记忆。
\textbf{实现细节：} 3DGM（2024.05） 面向自动驾驶场景中的多次遍历（multitraverse）建图问题：同一区域被车辆反复经过，每次遇到不同的行人和车辆。传统 3DGS 在此场景下的困境是——短暂物体（行人、车辆）破坏多视角一致性，而遮挡、非朗伯表面等噪声进一步干扰重建。3DGM 的核心洞察是：环境结构在遍历间保持一致，短暂物体则频繁变化，因此可用重复遍历的共识信号作为自监督，自动区分环境与物体。技术流程分三个阶段。第一阶段（初始化）将多次遍历的图像联合输入 SfM 获取初始高斯点和相机位姿，多遍历的丰富视角天然提升三角化质量。第二阶段（EmerSeg 短暂物体分割）通过鲁棒特征蒸馏实现：在每个高斯上附加 $d$ 维 DINOv2 语义特征 $f_i$，联合优化 RGB 和特征渲染损失 $L = \sum_t(L_{\mathrm{rgb}}(I_t(\xi_t; G), I_t) + L_{\mathrm{feat}}(F_t(\xi_t; G), F_t))$。由于环境像素（内点）在数量上远超短暂物体像素（外点），模型主要被内点梯度驱动去学习环境结构，导致短暂物体区域在特征空间产生高残差。接着利用特征残差挖掘：归一化残差图后，低于阈值 $\tau_1 = 0.3$ 的像素被抑制，再经轮廓检测、小轮廓过滤、天空轮廓去除和凸包提取，最终生成短暂物体掩码 $M_t$——整个过程无需额外训练。第三阶段（EnvGS 鲁棒优化）用掩码屏蔽外点：$L = \sum_t L_{\mathrm{rgb}}(M_t \odot I_t(\xi_t; G), M_t \odot I_t)$，仅从环境像素优化高斯属性，输出干净的环境地图。
\textbf{方法信息：} NVIDIA 参与了这篇论文。



% METHOD: 2409.09295
\paragraph{GEVO（用高斯地图压缩单目里程计历史）（2024.09）}
GEVO~\cite{2409_09295} 从视觉里程计的内存瓶颈切入。传统单目 VO 为了做局部优化和回环，常需要保存大量历史图像或特征；GEVO 的改变是用 Gaussian map 渲染历史观测，从而不再长期保存所有原图。它不是把 3DGS 作为最终美观地图，而是把高斯视为可查询的历史状态压缩器。

方法先用 MVS cost volume 和 confidence propagation 做 occupancy-preserving initialization，使初始高斯更能覆盖真实可见结构。更新时，它只优化当前视角下不一致且足够可见的高斯，并处理 recursive occlusion 和 incomplete ray occlusion，避免被遮挡区域反复写错。相对 Photo-SLAM 的 hyper primitive map，GEVO 更偏里程计系统设计：高斯的价值在于用几十 MB 级别状态替代历史图像缓存，同时保留足够外观来支持后续 tracking。
\textbf{实现细节：} GEVO（2024.09） 同样利用 3DGS 地图渲染替代原始图像存储，但其关注点与 Splat-SLAM 的全局一致性不同——GEVO 聚焦于内存效率与灾难性遗忘问题。动机。GEVO 的核心策略是仅保留当前滑动窗口内的原始图像（8--10 帧），用已有地图渲染的图像替代历史关键帧——但由于渲染图像的保真度会随地图退化而逐渐下降，单纯替换并不足够，因此 GEVO 提出了三个互补机制来维持渲染图像的可用性。占用保持初始化。与 Splat-SLAM 的代理深度图（融合多视点深度与单目先验）不同，GEVO 采用基于多视立体（MVS）代价体的初始化策略：在滑动窗口的 8--10 个关键帧上构建光度一致性代价体 $V(u,d)$（64 个深度假设，图像 4 倍下采样至仅需 4.7 MB），经信念传播提取深度图后上采样至全分辨率，再通过 SPGF$^{*}$ 算法生成高斯。关键创新在于高斯同时表示障碍物和自由区域——自由区域高斯编码了无障碍空间，用于后续识别和剪枝错误遮挡已有高斯的新高斯（即缓解 RO）。
\textbf{方法信息：} Amazon 参与了这篇论文。
\textbf{设计动机与边界：} 相对原始 3DGS 以独立高斯和图像重建误差为主的设定，GEVO 把问题落在以下可实现的改动上：GEVO包含三个核心创新：(1)渲染替代存储：不再存储大量历史图像，而是从现有3DGS地图渲染图像来指导高斯优化，从根本上降低内存开销；(2)占用保持初始化：基于多视图立体（MVS）构建代价体积，通过置信度传播估计深度，紧凑编码障碍物自由区域，用于识别和修剪导致回溯遮挡的新高斯；(3)两阶段优化：首先更新与当前滑动窗口不一致且充分可见的少量高斯，避免过去视图的高斯过拟合到当前窗口。该方法有效解决了回溯遮挡（RO）和不完全射线遮蔽（IRO）两种灾难性遗忘模式。其设计目的不是单独增加一个模块，而是用上述表示、监督或渲染约束消除原流程在目标任务中的具体失配。GEVO是首个内存高效的高斯SLAM框架，通过渲染替代存储历史图像将内存开销降至约58MB，比先前方法减少最多94倍，同时保持可比的地图保真度。适用时还应保留论文给出的边界：渲染图像的质量随时间缓慢退化；仅支持单目RGB输入；在纹理贫乏场景中初始化效果受限。
% METHOD: 2409.13055
\paragraph{MGSO（直接法 SLAM 点云到高斯的在线桥接）（2024.09）}
MGSO~\cite{2409_13055} 继承 DSO-like photometric SLAM 的密集像素跟踪路线。它的增量是把直接法 SLAM 前端产生的 dense structured point cloud 及时转成 3DGS，而不是等离线 SfM 完成后再训练。SLAM 模块负责跟踪稠密像素、估计相机运动和结构点，高斯模块从这些点云初始化，并在未映射区域补充新的高斯。

具体流程是并行维护直接法 SLAM 和 3DGS 两个线程。前端持续提供局部点云和位姿，后端把已有点云转成高斯位置、尺度和颜色初值，再根据渲染误差优化高斯属性。项目实现面向 monocular real-time photometric SLAM。相对 GSO-SLAM 的双向耦合，MGSO 更像单向桥接：直接法系统先保证跟踪，高斯层负责高保真地图；但它比纯离线 3DGS 更接近实时 RGB-only 建图。
\textbf{方法信息：} 代码：\url{https://github.com/8bit-nyk/mgso_dev}。
\textbf{设计动机与边界：} 相对原始 3DGS 以独立高斯和图像重建误差为主的设定，MGSO 把问题落在以下可实现的改动上：双模块并行架构：(1) SLAM模块：基于DSO的光度SLAM，跟踪密集像素集（非稀疏特征点），通过像素选择策略（分块选最高梯度像素）生成密集结构化点云，窗式关键帧管理保证分布均匀；(2) 3DGS模块：用SLAM输出的密集点云初始化高斯，在未映射区域持续添加高斯，高质量初始化减少所需优化步骤；(3) 关键观察：原始3DGS最终高斯位置和分布与DSO点云输出相似，因此用DSO初始化可大幅减少优化需求。系统仅使用RGB输入，不依赖深度传感器。其设计目的不是单独增加一个模块，而是用上述表示、监督或渲染约束消除原流程在目标任务中的具体失配。提出首个单目光度SLAM-3DGS集成系统MGSO，通过密集结构化点云初始化加速3DGS优化，在质量、内存效率和速度上达到最佳平衡，仅需RGB输入即可在笔记本硬件运行适用时还应保留论文给出的边界：单目尺度模糊性；快速运动时跟踪可能失败；光度SLAM对光照变化敏感
% METHOD: 2411.17660
\paragraph{DROID-Splat（端到端 SLAM 前端与高斯渲染后端的耦合）（2024.11）}
DROID-Splat~\cite{2411_17660} 的锚点很明确：使用 DROID-SLAM 负责鲁棒跟踪和全局优化，用 3DGS 负责实时可渲染地图。许多早期 3DGS SLAM 系统在前端位姿漂移后很难修正地图；DROID-Splat 的改变是把 recurrent update、局部前端、全局后端和 loop closure 保留下来，再把高置信 disparity 转成高斯初始化。

系统由并行模块组成。局部前端根据 dense flow 和 GRU 更新估计相邻关键帧约束，后端做全局 BA 和回环，渲染线程从 disparity 初始化高斯，并通过共视和置信度过滤避免把错误深度写入地图。它同时支持 monocular、RGB-D 和未知内参输入，未知内参时还可近似 DROID-Calib 那样优化相机参数。项目实现围绕 DROID-style tracking 与 Gaussian rendering 管线展开。相对 Splat-SLAM，DROID-Splat 更像把成熟端到端 SLAM 前端作为主干，高斯地图是其稠密可视化和新视角合成后端。
\textbf{实现细节：} 同一技术路线中更早的探索是 DROID-Splat（2024.11） ，同样以 DROID-SLAM 的端到端稠密光流网络为跟踪核心，但采用了更简洁的"直耦合"架构——跟踪器与渲染器通过共享位姿状态直接耦合，无需 HI-SLAM2 那样的四模块分离与离线精化流程。动机。纯高斯渲染优化的跟踪精度始终落后于传统方法和端到端 SLAM，而端到端跟踪器虽鲁棒但无法产生照片级场景表示；DROID-Splat 旨在弥合这一差距，通过"稠密跟踪 + 稠密渲染"的统一框架同时实现高精度定位和高质量重建。
\textbf{实现补充：} 实现上还要核对输入、输出、参数、特征、初始化、训练损失、正则、渲染或推理接口，以及优化中的剪枝、分裂、采样、编码或解码步骤；这些抓手决定该方法到底改了 3DGS 管线的哪一环。
\textbf{设计动机与边界：} 相对原始 3DGS 以独立高斯和图像重建误差为主的设定，DROID-Splat 把问题落在以下可实现的改动上：系统架构包含四个并行模块：(1) 局部前端：基于DROID-SLAM的密集光流估计和Bundle Adjustment，通过帧图管理关键帧，使用循环GRU更新视差和位姿。(2) 全局后端：对关键帧图进行全局BA优化。(3) 回环检测：检测回环并通过位姿图BA实现全局一致性。(4) 密集渲染器：基于3DGS，从前端获取密集视差图，经共可见性过滤和高置信度筛选后初始化和优化高斯原语。支持集成Metric3D/ZoeDepth/DepthAnything等单目深度先验。还支持相机内参联合优化（类似DROID-calib），处理未知内参的wild视频数据。其设计目的不是单独增加一个模块，而是用上述表示、监督或渲染约束消除原流程在目标任务中的具体失配。将端到端SLAM的鲁棒跟踪与3DGS的实时渲染结合，在常见SLAM基准上同时达到SOTA跟踪和渲染结果。支持单目和RGBD输入，通过单目深度预测和相机标定目标实现在未知内参wild数据上的鲁棒重建。组件级消融分析各模块实际贡献。适用时还应保留论文给出的边界：依赖DROID-SLAM的预训练权重；实时性受限于GPU性能；极端动态场景仍有挑战
\textbf{方法信息：} 代码：\url{https://github.com/chenhoy/droid-splat}；GitHub 关注度较高，约 414 stars。



% METHOD: 2509.14191
\paragraph{MCGS-SLAM（纯 RGB 多相机高斯 SLAM）（2025.09）}
MCGS-SLAM~\cite{2509_14191} 把 RGB-only 高斯 SLAM 推到多相机设置。它面对的问题不是单个单目相机缺尺度，而是多相机系统中不同视角、不同曝光和不同局部尺度估计如何组成统一 Gaussian map。方法以 Metric3Dv2 这类单目深度与法向估计器提供几何先验，再通过 RFT keyframes 选择关键帧，避免每路相机都向地图写入冗余高斯。

核心改变在两个优化模块。JDSA 负责联合深度和尺度对齐，让不同相机和不同时间的单目深度进入同一尺度；MCBA 则做多相机 bundle adjustment，同时约束相机位姿、高斯地图和尺度一致性。方法还用 low-rank scale consistency 防止各相机尺度漂移，并在离线 refinement 阶段进一步提高渲染和几何质量。相对 Splat-SLAM 的单相机 RGB 路线，MCGS-SLAM 的增量是把尺度对齐从单序列问题变成跨相机系统问题。
\textbf{实现细节：} MCGS-SLAM（2025.09） 将 S3PO-GS 的单目设定推广至多相机 rig（3--5 个已标定相机），核心挑战从"单目尺度漂移"转为"多相机间的跨视角深度与尺度一致性"。动机。多相机 rig 中每个相机独立运行单目深度估计，各相机的深度预测存在独立的空间变化尺度偏移，直接融合会导致跨视角深度不一致，进而引起几何退化。，MCGS-SLAM 设计了多相机光束调整（MCBA）模块，将时序对（同相机跨时间）和跨视角对（同时间不同相机）的投影残差统一到一个光束调整框架中。时序对通过已知外参 $\mathbf{T}_C^B$ 将体位姿变换映射到相机坐标系，跨视角对则直接利用预标定的相机间外参 $\mathbf{T}_{C_i C_j}$，使位姿更新 $\Delta\boldsymbol{\xi}$ 和深度更新 $\Delta\mathbf{d}$ 通过 Schur complement 联合求解。在尺度对齐方面，S3PO-GS 的 patch 级局部对齐仅在单帧内操作，而多相机设定中各相机的单目深度预测（Metric3Dv2）存在独立的空间变化尺度偏移。MCGS-SLAM 为此引入联合深度-尺度对齐（JDSA），为每个关键帧分配一个可学习的 $m \times n$ 尺度网格，通过双线性插值生成逐像素尺度因子，将预测深度重缩放为优化深度；JDSA 与 MCBA 交替运行，既利用多视角光度-几何一致性约束尺度，又用尺度对齐损失 $\|\tilde{d}_i \cdot B_i - d_i\|^{}2$ 保证对齐稳定，避免了直接联合优化尺度与位姿导致的不收敛问题。建图阶段，当 MCBA 更新关键帧位姿 $\Delta\mathbf{T}_t$ 时，所有锚定于该帧的高斯的均值和协方差同步传播（$\boldsymbol{\mu}_i \leftarrow \Delta\mathbf{T}_t \cdot \boldsymbol{\mu}_i$，$\Sigma_i \leftarrow \mathbf{R}\Sigma_i\mathbf{R}^{}\top$），避免重新初始化或重新渲染。
\textbf{设计动机与边界：} 相对原始 3DGS 以独立高斯和图像重建误差为主的设定，MCGS-SLAM 把问题落在以下可实现的改动上：MCGS-SLAM系统包含在线跟踪和离线精化两个阶段：(1) 在线跟踪：使用Metric3Dv2估计每帧深度和法线图，基于RFT流进行关键帧选择；联合深度和尺度对齐（JDSA）解决单目深度预测的尺度歧义，使用可学习m×n尺度网格；多相机光束调整（MCBA）通过密集光度残差和几何残差联合精化位姿和深度，JDSA损失函数强制多视角光度一致性；(2) 3DGS映射：增量将精化后的关键帧融合到全局3DGS地图；(3) 尺度一致性模块：使用低秩先验强制跨视角度量对齐；(4) 离线精化：联合优化所有相机位姿和高斯参数，提升全局一致性。系统仅使用RGB输入，支持实时大规模运行。为了消除原始 3DGS 在该目标任务中的具体失配，方法把上述表示、监督或渲染约束作为可学习和可执行的改动，而不是单独堆叠模块。提出MCGS-SLAM，首个纯RGB多相机3DGS-SLAM系统，通过多相机光束调整和尺度一致性模块实现更鲁棒的位姿估计和更完整的场景重建，多相机广视野可重建单目遗漏的侧视区域。适用时还应保留论文给出的边界：需要多相机标定；多相机融合增加计算复杂度；尺度一致性模块的收敛可能不稳定；对相机配置的泛化性有待验证。
\textbf{方法信息：} Microsoft 参与了这篇论文。



% METHOD: 2602.11714
\paragraph{GSO-SLAM（DSO 与 2DGS 的双向耦合）（2026.02）}
GSO-SLAM~\cite{2602_11714} 继承直接稀疏里程计 DSO 和 2DGS 表面化表示。普通做法往往是前端给位姿，后端只负责渲染地图；GSO-SLAM 的改变是让 DSO 与 2DGS 互相提供观测。它把系统写成近似 EM 框架：E-step 中，DSO 给出半稠密深度、关键帧位姿和像素关联，2DGS 渲染深度又作为伪观测反过来补充直接法；M-step 中，系统更新高斯地图，使其和最新轨迹、深度一致。

具体实现中，高斯从 DSO 的半稠密点和关键帧初始化，随后通过 ray-splat intersection 获取多视图一致深度。2DGS 的优势在于高斯更贴近表面，深度比体积式 splatting 更适合约束直接法前端。项目实现围绕 DSO 和 Gaussian Splatting 双向耦合展开。相对 Photo-SLAM 的“特征前端加渲染后端”，GSO-SLAM 更紧密：高斯地图不是被动输出，而是直接参与视觉里程计的深度和位姿估计。
\textbf{设计动机与边界：} 相对原始 3DGS 以独立高斯和图像重建误差为主的设定，GSO-SLAM 把问题落在以下可实现的改动上：GSO-SLAM的核心创新：1)EM框架下的双向耦合：E步利用DSO的半稠密深度和2DGS的渲染深度作为伪观测联合优化，M步更新高斯场景表示；2)高斯初始化策略：利用DSO提供的图像信息、关键帧位姿和像素关联直接生成接近最终状态的高斯初始化，消除启发式初始化方法；3)2DGS表示：采用2DGS而非3DGS，利用射线-splat交点实现多视角一致深度和更好的几何精度；4)实时性能：紧密耦合优化避免了松耦合方法中的信息冗余，同时保留了耦合方法的深度利用重建场景的能力。为了消除原始 3DGS 在该目标任务中的具体失配，方法把上述表示、监督或渲染约束作为可学习和可执行的改动，而不是单独堆叠模块。提出在EM框架下双向耦合DSO与2DGS的SLAM系统，通过半稠密深度估计和高斯渲染深度的伪观测联合优化，消除冗余计算开销。设计高斯初始化策略利用VO信息直接近似最终场景。适用时还应保留论文给出的边界：依赖单目相机的光度一致性假设；对快速运动和纹理缺乏场景可能退化；2DGS在复杂3D结构表示上可能受限。
\textbf{方法信息：} 代码：\url{https://github.com/lab-of-ai-and-robotics/gso-slam}。



\section{RGB-D 与稠密融合 SLAM}

RGB-D 方法的基础优势是深度给了尺度和可见表面，劣势是深度噪声、遮挡边界、空洞和设备范围限制会直接污染高斯地图。本节方法大多围绕 SplaTAM 这一锚点展开：先用 RGB-D 观测初始化或更新高斯，再通过颜色、深度和轮廓损失优化相机与地图。后续变化主要体现在三处：用 TSDF 或结构点云稳定几何，用不确定性或频率分析控制哪些观测可信，用系统剪枝和硬件调度让在线 RGB-D SLAM 真正实时。


% METHOD: 2312.02126
\paragraph{SplaTAM（RGB-D 高斯 SLAM 的代表性锚点）（2023.12）}
SplaTAM~\cite{2312_02126} 是本章最重要的 RGB-D 锚点之一。它把 3DGS 从离线重建改造成 dense RGB-D SLAM：每个高斯带有位置、半径、不透明度和 RGB 颜色，系统同时渲染颜色、深度和 silhouette。原始 3DGS 依赖已知相机和离线训练，而传统 RGB-D SLAM 的地图通常不是照片级可渲染；SplaTAM 的改变是让同一组高斯既支持跟踪又支持建图。

跟踪阶段，系统固定地图并优化当前相机位姿，使渲染颜色和深度对齐输入帧。建图阶段，系统优化高斯参数和关键帧位姿，并用 silhouette mask 判断哪些区域已经被地图解释、哪些区域仍是 unmapped，从而只在未覆盖区域添加新高斯。这个机制比单纯残差 densification 更适合在线 SLAM，因为它把“新增高斯”绑定到可见覆盖，而不是让所有高残差像素都触发膨胀。公开代码和高星标说明其实现已成为后续 RGB-D Gaussian SLAM 的常用参照。相对 Gaussian Splatting SLAM 的早期解析追踪，SplaTAM 的锚点意义在于完整定义了 track-and-map 的 RGB-D 高斯闭环。
\textbf{锚点细化：} SplaTAM：较早将 3D 高斯作为 RGB-D SLAM 的唯一三维表示，针对标准 3DGS 自适应致密化依赖离线全局视角统计、不适合在线帧序到达的局限，用\textbf{轮廓感知在线致密化}替代基于梯度统计的自适应致密化。核心思想是\textbf{渲染轮廓图像量化逐像素认知不确定性}——轮廓图像是渲染过程的副产品，不记录颜色或深度，只记录每个像素被高斯"覆盖了多少"（接近 1 表示已充分观测、接近 0 表示尚未建图），由此系统较早能显式判断"哪里没建好"。基于这一不确定性图，SplaTAM 用各向同性球状高斯（8 维参数）替代完整协方差椭球，跟踪阶段仅对轮廓超过 0.99 的像素施加约束，致密化阶段在轮廓不足或深度不一致的区域添加新高斯。
\textbf{锚点细化：} SplaTAM：首次将 3D 高斯作为 RGB-D SLAM 的唯一三维表示，针对标准 3DGS 自适应致密化依赖离线全局视角统计、不适合在线帧序到达的局限，用\textbf{轮廓感知在线致密化}替代基于梯度统计的自适应致密化，核心思想是\textbf{渲染轮廓图像量化逐像素认知不确定性}——轮廓图像是渲染过程的副产品，不记录颜色或深度，只记录每个像素被高斯"覆盖了多少"（接近 1 表示已充分观测、接近 0 表示尚未建图），由此系统首次能显式判断"哪里没建好"。基于这一不确定性图，SplaTAM 用各向同性球状高斯（8 维参数）替代完整协方差椭球，跟踪阶段仅对轮廓超过 0.99 的像素施加约束，致密化阶段在轮廓不足或深度不一致的区域添加新高斯。
\textbf{方法信息：} NVIDIA 参与了这篇论文；代码：\url{https://github.com/spla-tam/splatam}；GitHub 关注度高，约 2146 stars；引用很多，Semantic Scholar 约 632 次。



% METHOD: 2403.16095
\paragraph{CG-SLAM（不确定性感知的早期 RGB-D 高斯场）（2024.03）}
CG-SLAM~\cite{2403_16095} 是 RGB-D 高斯 SLAM 中较早强调 uncertainty-aware field 的工作。它继承 Gaussian Splatting SLAM 的解析 EWA 位姿导数，但加入深度不确定性模型来判断哪些高斯对跟踪更有价值。普通 RGB-D 损失会把噪声深度和真实表面同等看待，CG-SLAM 的改变是让低不确定 primitive 在 frame-to-model tracking 中承担更多权重。

具体上，它推导 3D Gaussian field 关于相机位姿的解析雅可比，并实现专门 CUDA kernel 支持在线优化。地图更新时使用尺度正则避免高斯退化，跟踪时最小化低不确定基元的重渲染误差。相对 SplaTAM 的 silhouette-driven map expansion，CG-SLAM 的重点在位姿优化和高斯选择：不是所有可见高斯都同样适合追踪，系统应优先使用几何和深度更可靠的部分。
\textbf{实现细节：} CG-SLAM（2024.03） 修改了标准 3DGS 的场表示与优化策略——用一致性几何稳定的不确定性感知高斯场替代原始无约束的高斯场，$\alpha$-混合 渲染管线保持不变。CG-SLAM 系统性地解决了标准 3DGS 直接用于 SLAM 时的三重缺陷。在基线做法中，原始 3D 高斯场作为纯视角合成技术，其强各向异性使其倾向于过拟合输入图像而非准确建模 3D 表面。具体表现为高斯基元在优化后呈现出"箭形伪影"（arrow-shaped artifacts），即某个轴方向被极度拉伸以拟合个别视角的颜色，而非贴合真实表面几何；高斯基元可能不对齐环境表面，导致外推能力差并进一步退化相机跟踪；此外，建图过程中高斯基元数量的持续增长会不可避免地拖慢跟踪效率。CG-SLAM 的核心改进是构建一个一致性、几何稳定且不确定性感知的 3D 高斯场，具体包含三个层面的设计： （1）各向同性正则化。在标准 3DGS 中，高斯基元的协方差矩阵 $\Sigma = RSS^{}\top R^{}\top$ 中缩放矩阵 $S = \mathrm{diag}([s^x, s^y, s^z])$ 的三个分量可以独立优化，没有任何约束阻止某个方向极度拉伸。这种自由度在离线新视角合成中问题不大（多视角约束自然抑制极端形状），但在 SLAM 的在线增量建图场景中，早期帧数少、多视角约束不足，高斯极易退化为沿某一方向极度拉伸的"箭形"，既不贴合表面又产生误导梯度。为了抑制这种各向异性过拟合，CG-SLAM 引入软尺度正则化损失，鼓励高斯椭球趋向球形，在跟踪精度与渲染真实感之间取得平衡： 直观地说，这个损失惩罚尺度因子的最大/最小比值超过阈值 $\gamma$ 的高斯——当三个尺度因子接近相等时比值为 1，越拉伸比值越大。其中 $s_i^{x,y,z}$ 为第 $i$ 个高斯的尺度因子，$\gamma=1.0$ 控制允许的各向异性程度。
\textbf{方法信息：} 引用很多，Semantic Scholar 约 91 次。



% METHOD: 2408.12677
\paragraph{GSFusion（TSDF 几何与高斯外观的混合）（2024.08）}
GSFusion~\cite{2408_12677} 继承 SplaTAM 的 RGB-D 在线建图，但它认为纯高斯地图在几何一致性上仍不够稳：深度噪声会造成浮点体，局部高斯增密也可能在表面附近重复堆叠。它的改变是把 TSDF voxel geometry 和 3D Gaussian rendering 合成一个混合表示，让 TSDF 负责稳健几何融合，高斯负责照片级外观。

系统每来一帧先用深度更新 TSDF，再根据视觉和几何信息在 3D 空间分配新高斯。为了减少无效高斯，它还用 image quadtree downsampling 在图像上自适应采样，把更多资源放在纹理或几何变化明显区域。优化阶段同时使用颜色、深度和正则项约束。项目仓库材料与 GSFusion 项目一致。相对 SplaTAM，GSFusion 的核心不是重新设计 tracking，而是把“地图几何是否可信”交给 TSDF 累积，减少高斯数量和表面伪影。
\textbf{实现细节：} GSFusion（2024.08） 的核心转变在于地图表示层面。动机。SplaTAM 将 3DGS 作为唯一地图表示，但纯高斯地图难以提供精确的几何约束（如用于机器人导航的障碍物检测）；同时，SplaTAM 的逐像素致密化会产生大量冗余高斯，导致内存持续膨胀。SplaTAM 将 3DGS 作为唯一地图表示，GSFusion 则构建双地图系统——在 Supereight2 体积映射引擎上同时维护一张 TSDF 体积地图和一张 3DGS 地图。TSDF 地图提供准确的几何结构，3DGS 地图负责照片级渲染，两者互补而非替代。这一设计带来的核心收益是：可以利用 TSDF 体素的权重信息来去重——查询待初始化位置最近的体素，若权重为 1 表示首次观测、可安全添加高斯；若权重已超过 1 表示该位置已被访问过、无需重复初始化。每个叶节点中心仅生成一个高斯，初始旋转为单位矩阵，缩放由叶节点中心到其角的反投影长度确定。这一策略使高斯数量从逐像素采样降低一个数量级，且紧凑的初始分布直接消除了 SplaTAM 所需的后续致密化和剪枝过程。在线优化策略。SplaTAM 对所有帧统一执行固定次数的建图优化。GSFusion 则利用"新初始化高斯数量"作为信息增益指标来区分关键帧与非关键帧——新高斯数超过阈值时标记为关键帧。关键帧执行 $m$ 次优化迭代，非关键帧仅执行 $n$ 次（$n < m$），节省的计算量被随机分配给历史关键帧列表中的帧，以缓解在线优化中的遗忘问题（catastrophic forgetting）。
\textbf{设计动机与边界：} 相对原始 3DGS 以独立高斯和图像重建误差为主的设定，GSFusion 把问题落在以下可实现的改动上：GSFusion采用混合场景表示：基于八叉树的TSDF体素网格用于几何重建，3D Gaussian用于高质量渲染。核心创新在于增量建图流程：(1)先对输入深度帧执行TSDF融合更新体素；(2)利用视觉和几何信息（而非逐像素采样）在3D空间中分配新Gaussian，通过图像四叉树结构在保留结构区域像素的同时对均匀区域降采样，大幅减少初始化Gaussian数量；(3)对Gaussian参数进行在线优化，包含颜色损失、深度损失和正则化损失。该设计避免了纯Gaussian方法中因密集采样导致的存储开销和效率问题，同时TSDF提供的几何约束有助于减少浮动伪影。其设计目的不是单独增加一个模块，而是用上述表示、监督或渲染约束消除原流程在目标任务中的具体失配。提出GSFusion，首次将3D Gaussian嵌入到体素建图系统中，利用TSDF提供的几何信息指导Gaussian初始化，并使用图像四叉树结构大幅减少splat数量，实现同时生成紧凑的3D Gaussian地图和体素地图，显著提升计算效率且不牺牲渲染质量。适用时还应保留论文给出的边界：论文未明确提及具体局限，但从方法设计看，依赖TSDF质量可能影响Gaussian初始化效果；四叉树降采样可能在极精细纹理区域丢失细节。
\textbf{方法信息：} 代码：\url{https://github.com/goldoak/gsfusion}；GitHub 关注度较高，约 228 stars。



% METHOD: 2501.05242
\paragraph{SEGS-SLAM（结构点云与外观运动嵌入）（2025.01）}
SEGS-SLAM~\cite{2501_05242} 的局部锚点是 RGB-D 3DGS SLAM，但它指出 SplaTAM 一类方法如果只依赖输入深度和渲染残差，结构边界和重复纹理区域仍容易漂移。方法的改变是引入 structure-enhanced initialization：使用视觉 SLAM 或 ORB-SLAM3 产生的结构化点云初始化高斯，使高斯分布从一开始就贴近稳定几何，而不是完全依赖后续 densification。

在建图中，SEGS-SLAM 用 geometry-driven densification 控制高斯增密，并用 Appearance-from-Motion Embedding 处理随相机运动变化的外观。不同于每帧存一个 appearance embedding，AfME 根据相机位姿预测外观调制，从而减少历史帧绑定和存储。项目实现围绕结构增强和 appearance embedding 展开。相对 GSFusion 的 TSDF 混合，SEGS-SLAM 更强调从 SLAM 结构点和运动外观中给高斯一个更稳的初始化与颜色解释。
\textbf{实现细节：} SEGS-SLAM（2025.01） 沿用 Photo-SLAM 的跟踪--建图解耦架构（ORB-SLAM3 前端定位 + 3DGS 后端渲染），但识别出一个被先前工作忽视的关键问题：3DGS-SLAM 系统在优化过程中高斯的底层结构会逐渐退化——即使 Photo-SLAM 用 ORB-SLAM3 点云初始化高斯，优化过程中高斯仍会偏离场景的真实几何结构，导致边缘模糊和结构失真。动机。标准 3DGS 的自适应致密化以光度损失梯度为唯一驱动，没有显式的结构保持机制，导致高斯在优化中逐渐偏离 ORB 特征点捕获的锐利边缘和几何结构。结构增强锚点（SEPM）。SEGS-SLAM 的核心观察是 Scaffold-GS 中基于锚点的 3DGS 组织方式能有效保持场景结构：锚点在训练中位置固定，其下属高斯的分布始终围绕锚点展开。SEGS-SLAM 将这一思想引入 SLAM——对每个关键帧的 ORB-SLAM3 点云 $P_k$ 做增量体素化 $V_k = \lfloor P_k / l \rfloor$（$l$ 为体素尺寸），将体素中心作为锚点。锚点继承点云的强结构先验，且因位置固定而持续保持场景结构，有效防止优化过程中的结构退化。高斯参数（位置偏移、颜色、旋转、缩放）通过轻量 MLP 从锚点特征解码，颜色分支额外融入外观嵌入。Appearance-from-Motion 嵌入（AfME）。AfME 改用相机位姿 $(R, t)$ 作为外观编码器的输入，通过轻量 MLP $M_a$ 将位姿映射为嵌入向量 $\phi^{(a)}_{R,t} = M_a(R, t)$，再送入颜色解码器。位姿作为输入有三个优势：每帧唯一、天然编码空间信息、比离散索引更连续。训练后固定高斯属性仅变化位姿输入时，渲染图像的光照条件随位姿平滑变化，验证了 AfME 能有效将外观编码到连续位姿空间中。频率金字塔正则化（FPR）。
\textbf{方法信息：} 代码：\url{https://github.com/leaner-forever/segs-slam}；GitHub 关注度较高，约 208 stars。



% METHOD: 2503.01109
\paragraph{FGS-SLAM（频率感知的稀疏-稠密地图融合）（2025.03）}
FGS-SLAM~\cite{2503_01109} 把 Fourier 频率分析接入 Gaussian SLAM。它继承稀疏几何跟踪加稠密高斯渲染的双地图结构，认为不同区域不应使用同样密度的高斯：低频平滑区域用少量大高斯即可，高频边缘或纹理区域才需要更密的小高斯。方法的改变是 frequency-based adaptive densification。

系统维护 sparse map 与 dense 3DGS map。稀疏地图通过高斯和 GICP 支持 tracking，稠密地图负责渲染；二者共享部分属性并联合优化。频率分析决定在哪些区域增加高斯、调整尺度或保留稀疏表示，从而减少盲目 densification。项目实现面向 Fourier-based real-time SLAM。相对 GSFusion 借助 TSDF 稳定几何，FGS-SLAM 主要从图像和场景频率角度决定地图分辨率。
\textbf{实现细节：} 在频域分析方向，FGS-SLAM（2025.03） 针对 SplaTAM 框架中高斯初始化的不确定性问题——空间域信息无法判别场景各区域需要"多少、多密"的高斯——提出用傅里叶频域分析驱动自适应致密化。动机。标准 3DGS-SLAM 的致密化基于轮廓值和深度误差等空间域指标，但这些指标无法区分"纹理丰富的细节区域需要密集高斯"和"平坦表面只需少量高斯"，导致高斯分配不均匀。与 SplaTAM 的核心差异体现在两方面。跟踪架构：SplaTAM 依赖可微渲染 + 轮廓掩码在图像空间优化位姿，需要维护稠密地图才能跟踪；FGS-SLAM 则构建独立的稀疏高斯地图，利用 GICP（广义迭代最近点）直接在 3D 空间中估计变换矩阵 $T$，将源点云与稀疏地图对齐。这一设计借鉴了 GICP 与 3DGS 共享高斯分布假设的数学基础，使跟踪不再受制于稠密地图质量。频域致密化：SplaTAM 基于轮廓值和深度误差判断"哪里缺高斯"；FGS-SLAM 先对当前帧做 2D 傅里叶变换，用高斯高通滤波器 $H(u,v) = 1 - e^{-D^{}2(u,v)/2D_0^{}2}$ 分离高/低频区域，再以三角阈值法自适应确定采样间隔——高频区域（纹理丰富的细节）以间隔 $m$ 密集采样并赋予小半径 $\lambda_h$，低频区域（平坦表面）以间隔 $n$（$n > m$）稀疏采样并赋予大半径 $\lambda_l$。此外，系统引入高斯缺失区域检测掩码 $M_m = M_i \cup M_d \cup M_c$（联合不透明度掩码、深度掩码和颜色掩码），仅在缺失区域新增高斯，避免盲目添加导致的爆炸式增长。稀疏地图与稠密地图共享高斯属性，通过联合渲染保持一致性。建图与跟踪并行运行于独立线程。系统整体以 36 FPS 运行，比 SplaTAM（约 4 FPS）快近一个数量级。
\textbf{设计动机与边界：} 相对原始 3DGS 以独立高斯和图像重建误差为主的设定，FGS-SLAM 把问题落在以下可实现的改动上：FGS-SLAM包含以下核心设计：1）傅里叶频域自适应致密化：利用傅里叶变换将图像空间域转为频域，分析颜色变化梯度与频率的关系。低频区域（颜色变化小）使用稀疏大高斯表示，高频区域（颜色变化剧烈）使用密集小高斯表示。建立高斯先验加速参数收敛，减少冗余高斯点；2）稀疏-密集地图融合：稀疏地图由3D高斯组成，通过GICP（广义迭代最近点）实现高效相机跟踪；密集地图同样基于3D高斯，通过3DGS优化实现高保真视觉表示。两种地图共享高斯属性集和联合优化；3）系统集成：两种地图功能独立但通过共享高斯属性统一，联合渲染实现属性优化一致性。密集地图自适应初始化新高斯，与稀疏地图共享缺失区域掩码。为了消除原始 3DGS 在该目标任务中的具体失配，方法把上述表示、监督或渲染约束作为可学习和可执行的改动，而不是单独堆叠模块。提出首个利用频域分析实现高质量3DGS建图的实时SLAM系统。基于傅里叶频域分析的自适应致密化方法建立高斯先验加速收敛，构建独立的稀疏地图（GICP跟踪）和密集地图（3DGS高保真表示），在Replica和TUM数据集上平均36 FPS。适用时还应保留论文给出的边界：频域分析增加了每帧计算开销；对快速运动的鲁棒性需进一步验证；仅支持RGB-D输入。
\textbf{方法信息：} 代码：\url{https://github.com/3dv-coder/fgs-slam}。



% METHOD: 2504.16545
\paragraph{ToF-Splatting（稀疏 ToF 深度补足 RGB-D SLAM）（2025.04）}
ToF-Splatting~\cite{2504_16545} 面向移动设备或低成本传感器上只有稀疏 Time-of-Flight 深度的场景。它继承 RGB-D 3DGS SLAM 的思路，但输入不再是完整深度图，而可能只有几十个稀疏深度点。直接把这些点喂给 SplaTAM 式系统会造成大面积未约束区域；方法的改变是通过多帧几何整合把稀疏 ToF 深度扩展成可用的 dense depth prior。

流程中，系统结合单目 RGB、稀疏 ToF 深度和多帧相机运动，在时间窗口内融合几何约束生成更稠密深度图，再作为高斯初始化和 tracking 监督。这样高斯地图仍可利用深度尺度，而不需要完整 RGB-D 设备。相对纯单目方法，它多了真实尺度锚点；相对标准 RGB-D 方法，它的关键增量是把稀疏深度通过多帧一致性补全到可建图程度。
\textbf{实现细节：} ToF-Splatting 的输入是单目 RGB、极低分辨率 ToF 深度、多帧相机位姿或短窗口运动；输出是补全后的 dense depth prior、高斯初始化和用于 tracking/mapping 的几何监督。它的核心不是新增高斯参数，而是在进入 3DGS SLAM 前先用多帧积分把稀疏深度变成稳定尺度约束。与之相邻的 PointSLAM++ 更偏向用 Scaffold-GS 风格神经锚点降低地图内存，DINO-SLAM 则把 DINO 特征与几何信息蒸馏到语义增强地图；三者分别回答“深度从哪来”“锚点如何管大地图”和“语义特征如何并入 SLAM”。
\textbf{实现补充：} 实现上还要核对输入、输出、参数、特征、初始化、训练损失、正则、渲染或推理接口，以及优化中的剪枝、分裂、采样、编码或解码步骤；这些抓手决定该方法到底改了 3DGS 管线的哪一环。
\textbf{方法信息：} Sony 参与了这篇论文。
\textbf{设计动机与边界：} 相对原始 3DGS 以独立高斯和图像重建误差为主的设定，ToF-Splatting 把问题落在以下可实现的改动上：ToF-Splatting针对AR/VR设备中低功耗ToF传感器（如仅64个深度点）的场景设计。核心创新是多帧集成模块，将三种信息源融合：(1)极稀疏ToF深度测量作为可靠的几何锚点；(2)单目色彩信息提供纹理和边缘线索；(3)多视角几何提供时间一致性和结构约束。融合后生成密集深度图，作为3DGS-SLAM管道的输入。管道基于3DGS表示场景，支持高效的追踪和建图。特别解决了极稀疏深度直接用于GS-SLAM时建图质量严重退化的问题。其设计目的不是单独增加一个模块，而是用上述表示、监督或渲染约束消除原流程在目标任务中的具体失配。提出首个专为极稀疏ToF深度数据设计的3DGS-SLAM管道ToF-Splatting，引入多帧集成模块，融合极稀疏ToF深度、单目色彩和多视角几何线索生成密集深度图，在合成和真实稀疏ToF数据上达到SOTA追踪和建图性能。适用时还应保留论文给出的边界：多帧集成增加计算延迟；对高速运动的鲁棒性有待验证；极稀疏深度下追踪精度仍受限。
% METHOD: 2510.06644
\paragraph{RTGS（面向边缘设备的冗余削减）（2025.10）}
RTGS~\cite{2510_06644} 仍然继承 RGB-D 3DGS SLAM，但它的核心问题变成工程可部署性：在线高斯建图的 bottleneck 不只在算法，也在 GPU 内存访问、反向传播缓存和 tile 调度。它的改变是 multi-level redundancy reduction，把冗余从算法层和硬件层同时去掉。

算法上，RTGS 重用 backward gradients 做 adaptive Gaussian pruning，避免单独计算复杂重要性；又根据关键帧和 alpha 信息做 dynamic downsampling，使低贡献区域少参与渲染和反传。硬件上，它设计 sub-block streaming、pixel pair scheduling、渲染与反向传播 buffer，以及 gradient merge unit，让边缘 GPU 也能保持 30 FPS 以上的 SLAM 更新。代码已公开，系统实现可以直接核查。相对 SplaTAM 和 GSFusion，RTGS 的锚点价值在于把“高斯 SLAM 是否实时”具体落到冗余渲染、梯度合并和内存流水线。
\textbf{实现细节：} 动机。3DGS-SLAM 在跟踪阶段对全部高斯执行完整渲染和反向传播，但绝大多数高斯对位姿梯度的贡献可忽略不计；同时，非关键帧与关键帧使用相同分辨率处理，造成大量冗余计算，严重限制了系统在边缘 GPU 上的实时性能。RTGS 在不改变 SplaTAM 高斯表示和渲染管线的前提下，引入两层冗余消除。自适应高斯剪枝利用跟踪过程已有的位置梯度 $\nabla_{\boldsymbol{\mu}}L$ 与协方差梯度 $\nabla_{\Sigma}L$ 计算重要性分数 $s = \|\nabla_{\boldsymbol{\mu}}L\|_2 + \beta\|\nabla_{\Sigma}L\|_2$（复用跟踪阶段的梯度，不引入额外开销），采用 mask-then-prune 策略：先在 $K$ 次迭代中屏蔽低重要性高斯（排除渲染但仍保留以计算 tile-Gaussian 变化率），第 $K+1$ 次迭代时永久移除。剪枝上限设为 50\%，超过此值 ATE 急剧上升。动态下采样将关键帧保持原始分辨率 $R_0$，非关键帧则根据距最近关键帧的时序距离逐步降采样：紧跟关键帧的非关键帧降至 $R_0/16$，后续每帧以因子 $m=2$ 递增，最高不超过 $R_0/4$，保证分辨率变化的平滑过渡以避免轨迹漂移。
\textbf{实现补充：} 实现上还要核对输入、输出、参数、特征、初始化、训练损失、正则、渲染或推理接口，以及优化中的剪枝、分裂、采样、编码或解码步骤；这些抓手决定该方法到底改了 3DGS 管线的哪一环。
\textbf{方法信息：} Amazon 参与了这篇论文；代码：\url{https://github.com/umn-zhaolab/rtgs}。



% METHOD: 2603.06989
\paragraph{MipSLAM（抗锯齿高斯 SLAM）（2026.03）}
MipSLAM~\cite{2603_06989} 把 Mip-Splatting 的抗锯齿思想引入 SLAM。普通 3DGS SLAM 在相机尺度、分辨率或内参变化时，细小高斯会产生 aliasing，进而影响 tracking residual 和地图复用。MipSLAM 的改变是 elliptical adaptive anti-aliasing：在椭圆域内以自适应采样点和权重积分高斯贡献，让渲染图在不同尺度下更稳定。

系统还加入 spectral-aware pose graph optimization，把高频噪声和真实几何错位区分开来，减少由 aliasing 引发的漂移。它支持在不同相机内参、分辨率和尺度下复用同一高斯地图。项目实现包含 alias-free SLAM 组件。相对 VarSplat 的不确定性路线，MipSLAM 处理的是另一类 tracking 噪声：渲染采样本身带来的频率混叠。
\textbf{实现细节：} MipSLAM（2026.03） 修改了标准 3DGS 的像素级高斯贡献计算——用椭圆自适应数值积分替代像素中心点采样——$\alpha$-混合 渲染管线和深度渲染保持不变。动机。标准 3DGS 在像素中心做点采样来评估高斯贡献，当相机分辨率、焦距或缩放发生变化时，这种点采样违反 Nyquist-Shannon 采样定理，产生锯齿伪影；在 SLAM 的在线场景中，高斯基元持续投影且相机参数可能变化，锯齿会直接导致位姿估计漂移。该方法系统性地解决了 3DGS-SLAM 在相机参数（分辨率、焦距、缩放）变化时的锯齿与漂移问题，包含 EAA 抗锯齿和 SA-PGO 位姿图优化两个核心模块。（1）EAA：椭圆自适应抗锯齿。\textit{问题}。标准 3DGS 在计算像素 $u$ 处的高斯贡献时，直接在像素中心点处评估投影 2D 高斯 $G_i^{2D}(u)$，代入 $\alpha$-混合 公式得到颜色 $C(u) = \sum_{i=1}^{N} c_i \alpha_i G_i^{2D}(u) \prod_{j=1}^{i-1}(1 - \alpha_j G_j^{2D}(u))$ 和深度 $\hat{D}(u) = \sum_{i=1}^{N} z_i \alpha_i G_i^{2D}(u) \prod_{j=1}^{i-1}(1 - \alpha_j G_j^{2D}(u))$。
\textbf{实现补充：} 实现上还要核对输入、输出、参数、特征、初始化、训练损失、正则、渲染或推理接口，以及优化中的剪枝、分裂、采样、编码或解码步骤；这些抓手决定该方法到底改了 3DGS 管线的哪一环。
\textbf{方法信息：} 代码：\url{https://github.com/yzli1998/mipslam}。
\textbf{设计动机与边界：} 相对原始 3DGS 以独立高斯和图像重建误差为主的设定，MipSLAM 把问题落在以下可实现的改动上：MipSLAM解决3DGS-SLAM的核心问题——重建地图与固定相机参数绑定，换配置就产生严重伪影。(1)椭圆自适应抗混叠(EAA)：将Gaussian投影到椭圆域而非矩形采样，基于几何属性自适应选择采样点和权重，优先边界区域，在Riemann框架下高效近似积分，避免昂贵的解析计算。与MipSplatting的2D/3D滤波和Analytic-Splatting的解析积分不同。(2)频谱感知位姿图优化(SA-PGO)：将整个轨迹建模为时空信号，通过图拉普拉斯谱分析在频域中抑制高频噪声和漂移，而非独立处理每个位姿。(3)整体框架：联合估计相机位姿和Gaussian地图，支持不同内参、分辨率、缩放的相机配置间地图复用。其设计目的不是单独增加一个模块，而是用上述表示、监督或渲染约束消除原流程在目标任务中的具体失配。提出MipSLAM——首个频率感知3DGS SLAM框架，支持跨不同相机配置的高保真抗混叠NVS和鲁棒位姿估计。引入椭圆自适应抗混叠(EAA)算法通过几何感知数值积分近似Gaussian贡献，以及频谱感知位姿图优化(SA-PGO)在频域中抑制高频噪声和漂移。适用时还应保留论文给出的边界：EAA的采样精度与计算开销需要权衡。SA-PGO对长轨迹的频谱分析可能有较高内存需求。
% METHOD: 2603.09673
\paragraph{VarSplat（把不确定性渲染进每个像素）（2026.03）}
VarSplat~\cite{2603_09673} 继承 RGB-D 高斯 SLAM 的 frame-to-model tracking，但它指出 SLAM 中不是所有渲染误差都同样可信。深度噪声、遮挡边界和未充分优化的高斯都会产生不同不确定性；若仍用固定光度或深度损失，系统会把噪声当作真实几何优化。方法的改变是让每个 splat 学习 appearance variance，并通过总方差定律渲染 per-pixel uncertainty。

这些不确定性不只是可视化结果，而是进入 tracking、submap registration 和 loop closure。系统按子图维护高斯地图，优化相机、子图位姿、高斯属性和 variance；高不确定区域对位姿残差贡献较低，稳定区域则提供强约束。项目实现包含 uncertainty-aware mapping。相对 CG-SLAM 的深度不确定性，VarSplat 更进一步把方差作为高斯属性学习并可微渲染，使不确定性成为贯穿 SLAM 各阶段的权重。
\textbf{实现细节：} VarSplat（2026.03） 沿用 MipSLAM 所在的子地图式 3DGS-SLAM 管线（场景划分为多个子地图，每个子地图包含一组 3D 高斯基元），但切换了鲁棒性增强的切入角度。动机。在 SLAM 的在线增量建图中，高斯基元的参数误差不可避免——低纹理区域、反射表面和深度不连续处的高斯尤其不可靠，但这些区域在标准 3DGS 中与其他区域被同等对待，不可靠观测的梯度会误导位姿优化和地图更新。MipSLAM 在像素级贡献计算上做文章（用椭圆自适应积分替代点采样），而 VarSplat 则在标准 $\alpha$-混合 渲染之上为每个高斯基元增加一个外观方差参数 $\sigma_i^{}2$，通过全方差定律（Law of Total Variance）将其合成为逐像素不确定性图 $V$，从而让系统能"知道自己不知道什么"——低纹理区域、反射表面、深度不连续处自然产生高不确定性。该方法系统性地解决了 3DGS-SLAM 在不可靠观测下的位姿漂移与建图退化问题。（1）外观方差的直觉。
\textbf{实现补充：} 实现上还要核对输入、输出、参数、特征、初始化、训练损失、正则、渲染或推理接口，以及优化中的剪枝、分裂、采样、编码或解码步骤；这些抓手决定该方法到底改了 3DGS 管线的哪一环。
\textbf{技术实现补充：} VarSplat（2026.03） 沿用 MipSLAM 所在的子地图式 3DGS-SLAM 管线（场景划分为多个子地图，每个子地图包含一组 3D 高斯基元），但切换了鲁棒性增强的切入角度。动机。在 SLAM 的在线增量建图中，高斯基元的参数误差不可避免——低纹理区域、反射表面和深度不连续处的高斯尤其不可靠，但这些区域在标准 3DGS 中与其他区域被同等对待，不可靠观测的梯度会误导位姿优化和地图更新。MipSLAM 在像素级贡献计算上做文章（用椭圆自适应积分替代点采样），而 VarSplat 则在标准 $\alpha$-混合 渲染之上为每个高斯基元增加一个外观方差参数 $\sigma_i^2$，通过全方差定律（Law of Total Variance）将其合成为逐像素不确定性图 $V$，从而让系统能"知道自己不知道什么"——低纹理区域、反射表面、深度不连续处自然产生高不确定性。该方法系统性地解决了 3DGS-SLAM 在不可靠观测下的位姿漂移与建图退化问题。（1）外观方差的直觉。标准 3DGS 中每个高斯基元 $G_i$ 由位置 $\boldsymbol{\mu}_i$、不透明度 $\alpha_i$、尺度 $\mathbf{s}_i$、协方差矩阵 $\Sigma_i$ 和球谐系数 $c_i$（编码平均颜色）定义。VarSplat 在此基础上增加外观方差 $\sigma_i^2 \in \mathbb{R}^3$——注意它\textit{不是}定义几何范围的协方差 $\Sigma_i$，也\textit{不是}定义平均外观的球谐系数 $c_i$，而是建模围绕该平均颜色的观测离散程度。直觉上，即使球谐颜色正确建模了平均外观，不同视点下的实际颜色观测仍可能因遮挡边界变化、半透明表面混合等因素产生波动；$\sigma_i^2$ 正是量化这种波动的统计量。在深度不连续处和遮挡边界附近，微小的视角变化会改变重叠高斯的可见性和 $\alpha$ 权重，导致颜色观测不一致。因此这些区域自然学到更大的 $\sigma_i^2$。（2）方差渲染：全方差定律 $+$ $\alpha$ 合成。VarSplat 的核心创新在于：不依赖额外的预训练不确定性预测器，而是直接利用 3DGS 光栅化器已有的 $\alpha$-混合 管线，通过全方差定律将逐基元方差合成为逐像素方差。
\textbf{方法信息：} 代码：\url{https://github.com/anhthuan1999/varsplat}。
\textbf{设计动机与边界：} 相对原始 3DGS 以独立高斯和图像重建误差为主的设定，VarSplat 把问题落在以下可实现的改动上：VarSplat在标准3DGS参数（位置、不透明度、缩放、旋转、SH系数）基础上增加外观方差参数sigma\^{}2，通过全方差定律和alpha合成将其传播为per-pixel不确定性图V。系统采用子图策略，将场景分解为多个子图。不确定性在三个阶段发挥作用：(1)跟踪阶段：V图提供帧间姿态更新的可靠性度量；(2)子图配准阶段：per-pixel可靠性支持重叠子图间的中程对齐；(3)回环检测阶段：per-splat方差调节子图相似度以校正长程漂移。与传统预训练预测器不同，VarSplat从零开始学习方差参数，与外观和几何参数联合优化。系统端到端在线训练，随着地图增长同时优化姿态、高斯参数和sigma\^{}2。其设计目的不是单独增加一个模块，而是用上述表示、监督或渲染约束消除原流程在目标任务中的具体失配。提出首个学习per-splat外观方差的3DGS-SLAM系统，通过全方差定律渲染可微分per-pixel不确定性图，将不确定性应用于跟踪、子图配准和回环检测三个阶段，在Replica/TUM-RGBD/ScanNet/ScanNet++上实现鲁棒的跟踪、建图和新视角合成。适用时还应保留论文给出的边界：仅适用于RGB-D输入，不支持纯RGB场景；方差参数增加存储开销；未在超大规模户外场景中验证。

\section{多传感器 SLAM 与户外建图}

多模态方法的共同锚点是：高斯擅长表达可渲染外观，但位姿和几何初始化可以由 LiDAR、IMU、深度、雷达或事件相机提供。与 RGB-D 室内不同，户外和机器人场景经常有远距离、稀疏纹理、快速运动、光照变化和传感器失配。本节方法的差异主要体现在融合层级：有的把 LiDAR 或 IMU 作为前端位姿先验，有的把点云、SDF 或雷达回波直接变成高斯优化约束，有的则在农业、室内雷达或多季节环境中重写地图维护规则。
% METHOD: 2404.00923
\paragraph{MM3DGS SLAM（视觉、深度与 IMU 的多模态高斯 SLAM）（2024.04）}
MM3DGS SLAM~\cite{2404_00923} 继承 SplaTAM 式 track-and-map，但把输入扩展到 vision、depth 和 inertial measurements。它的关键问题是，当深度缺失或视觉快速运动时，单靠图像残差优化当前位姿不够稳定。方法的改变是让 IMU preintegration 提供初始位姿和尺度，再用高斯渲染的光度与深度损失做精化。

系统通过视觉 overlap 和 motion threshold 选择关键帧。tracking 阶段固定地图，优化 inverse camera transform，使输入图像和渲染结果一致；mapping 阶段联合优化高斯和最近关键帧位姿。高斯初始化可以来自深度，也可以来自运动估计推断出的结构。项目实现包含多模态输入与 UT-MM 数据集支持。相对 GS-LIVO 的 LiDAR 主导，MM3DGS 更接近室内或机器人 RGB-D-IMU 系统，重点是把惯性约束接到高斯 SLAM 的前端初始化。
\textbf{实现细节：} 最后，上述所有方法都假设传感器已预先标定。MM3DGS（2024.04） 挑战了这一前提，首次展示用未标定相机图像和惯性测量实现准确 SLAM——无需预先标定相机内参和相机-IMU 外参，系统在运行过程中在线估计所有标定参数，跟踪精度比当时最佳提升 3 倍。非传统传感器融合。本节讨论雷达、事件相机与热红外等非传统传感器与 3DGS SLAM 的融合方法。
\textbf{实现补充：} 实现上还要核对输入、输出、参数、特征、初始化、训练损失、正则、渲染或推理接口，以及优化中的剪枝、分裂、采样、编码或解码步骤；这些抓手决定该方法到底改了 3DGS 管线的哪一环。
\textbf{设计动机与边界：} 相对原始 3DGS 以独立高斯和图像重建误差为主的设定，MM3DGS SLAM 把问题落在以下可实现的改动上：框架包含四个主要阶段：位姿优化（跟踪）、关键帧选择、Gaussian初始化和建图。跟踪阶段固定3D Gaussian地图，通过反向相机变换优化位姿，利用光度损失和深度损失联合优化，使用不透明度阈值掩码处理地图未覆盖区域。惯性测量通过预积分提供初始位姿估计，加速收敛并提高跟踪精度。关键帧选择基于视觉重叠和运动阈值。Gaussian初始化使用深度信息或运动估计。建图阶段优化Gaussian属性和最近关键帧的位姿。框架支持单目相机+IMU和RGB-D+IMU两种配置。惯性预积分提供尺度感知，解决单目设置的尺度模糊问题。其设计目的不是单独增加一个模块，而是用上述表示、监督或渲染约束消除原流程在目标任务中的具体失配。提出首个使用3D Gaussians作为地图表示、融合无位姿相机图像和惯性测量的SLAM系统MM3DGS。通过融合相对位姿变换、深度估计和光度渲染质量的多模态损失函数实现关键帧建图和跟踪。发布多模态数据集UT-MM，相比SOTA 3DGS SLAM实现3倍跟踪精度提升和5%渲染质量提升。适用时还应保留论文给出的边界：IMU标定和时间同步增加系统复杂度；在高度动态场景中仍有限；需要相对准确的IMU内参标定；长时间运行需回环检测。
\textbf{方法信息：} 代码：\url{https://github.com/vita-group/mm3dgs-slam}；GitHub 关注度较高，约 210 stars。



% METHOD: 2404.06926
\paragraph{Gaussian-LIC（连续时间 LiDAR-inertial-camera 高斯 SLAM）（2024.04）}
Gaussian-LIC~\cite{2404_06926} 继承 LiDAR-inertial-camera fusion 的连续时间 odometry 锚点。它使用 Coco-LIC 类似的连续时间轨迹和 factor graph，每隔短时间更新位姿，再让 3DGS 建图线程生成照片级地图。相比离散关键帧系统，连续时间轨迹更适合处理 rolling motion、异步传感器和快速移动平台。

系统有两个互相解耦但数据相连的模块。前端是 LiDAR-inertial-camera odometry：LiDAR 点云、IMU 预积分和相机视觉约束进入连续时间因子图，得到稳定位姿，并兼容固态和机械旋转 LiDAR。后端是在线 3DGS 建图：地图初始化同时使用 LiDAR 点和在线视觉三角化 SfM 点，前者提供稀疏但可靠的几何骨架，后者补充 LiDAR 覆盖不到的视觉区域。为了解决户外天空和曝光变化，系统加入 background Gaussians 和 exposure-aware rendering；为保证实时性，代码以 C++/CUDA 管线实现增量高斯更新、关键帧选择和并行优化。

相对 GS-LIVO 的 hash octree 和 IESKF，Gaussian-LIC 的局部增量是连续时间 LIC 前端和高斯背景建模，使多传感器数据在时间和外观上更一致。相对室内 RGB-D 的 SplaTAM，它不依赖完整深度图，而是把 LiDAR、IMU 和相机组合成户外无界场景可用的位姿与初始化来源。公开仓库 README 明确对应 ICRA 2025 版本，并说明代码已发布。
\textbf{实现细节：} 本路线介绍 LiDAR 点云与视觉图像在 3DGS SLAM 中的紧耦合方法，以 Gaussian-LIC 系列为代表。LiDAR 提供精确的几何测量和远距离感知能力，与视觉的丰富纹理信息互补。Gaussian-LIC（2024.04） 是首个实时照片级 LiDAR-惯性-相机 SLAM 方法，旨在解决纯视觉 SLAM 在无界户外场景中的三大痛点：弱纹理导致几何退化、天空和变化曝光导致渲染伪影、以及大规模场景下实时性能的瓶颈。其核心思路是利用 LiDAR 的精确点云为高斯基元提供比视觉反投影更精确的初始化先验，从而加速收敛并提升重建质量。多模态里程计与数据准备。系统由两个并行模块组成：LiDAR-惯性-相机里程计（负责位姿估计和建图数据准备）以及 3DGS 在线建图（负责照片级重建）。里程计模块基于先前工作 Coco-LIC（一种连续时间因子图优化框架），每 0.1 秒增量优化一段轨迹，融合三种因子：点到地图的 LiDAR 因子、帧到地图的视觉因子和惯性因子，即使在剧烈运动和光照变化等挑战条件下也能输出鲁棒、实时的位姿。系统支持固态 LiDAR 和机械旋转 LiDAR 两种类型。对每个时间窗口内的 LiDAR 点，系统以每 $N_l$ 个点随机保留一个的比例降采样后，投影到最新图像上着色并剔除相机视场外的点，与位姿标注的图像一起作为一帧发送给建图线程。LiDAR-视觉联合高斯初始化。尽管 LiDAR 点提供了精确的几何先验，但其覆盖范围远小于相机——特别是机械旋转 LiDAR 的重复扫描模式。
\textbf{锚点细化：} Gaussian-LIC：针对纯视觉 SLAM 在无界户外弱纹理、天空、变化曝光和大场景实时性瓶颈的局限，用\textbf{LiDAR 点云联合视觉 SfM 点高斯初始化}替代纯视觉反投影，核心思想是\textbf{LiDAR 精确几何先验加速收敛并提升重建质量}——连续时间因子图（factor graph，一种将各类传感器约束建模为图节点间边的统一优化框架，位姿和传感器观测作为节点、约束作为边，最小化所有边的残差即得最优估计）融合 LiDAR-视觉-惯性三种因子输出鲁棒位姿，视觉滑窗三角化 SfM 点补偿 LiDAR 盲区（近距离物体 LiDAR 测不到），上半球采样天空高斯、处理无界场景，三项 CUDA 加速（逐高斯反传、瓦片剔除、稀疏 Adam）实现实时。
\textbf{方法信息：} Baidu 参与了这篇论文；代码：\url{https://github.com/APRIL-ZJU/Gaussian-LIC}；GitHub 关注度很高，约 563 stars。



% METHOD: 2410.17084
\paragraph{GS-LIVM（大规模 LiDAR-inertial-visual 照片级建图）（2024.10）}
GS-LIVM~\cite{2410_17084} 与 GS-LIVO 相邻，但更强调 large-scale outdoor mapping。户外 LiDAR 点云常稀疏且非均匀，直接用点初始化高斯会造成表面空洞；GS-LIVM 的改变是引入 Voxel-GPR，用 Gaussian Process Regression 在体素内补全 LiDAR 点云，得到更均匀的 completion points。

高斯初始化时，它用 covariance-centered Gaussian initialization 让高斯尺度和局部点云统计匹配；渲染优化中加入 similarity regularization，减少相机方向偏置导致的外观不一致。项目实现围绕 LiDAR-inertial-visual mapping。相对 GS-LIVO 的实时 odometry 与 hash map，GS-LIVM 的增量更像建图质量修正：先把稀疏 LiDAR 几何补成更适合高斯覆盖的表面，再做照片级渲染优化。
\textbf{设计动机与边界：} 相对原始 3DGS 以独立高斯和图像重建误差为主的设定，GS-LIVM 把问题落在以下可实现的改动上：方法基于LiDAR-惯性-视觉融合SLAM框架进行鲁棒状态估计和点云坐标变换，核心创新：(1)体素级高斯过程回归(Voxel-GPR)：对每帧扫描的体素使用GPR将不均匀的点云变换为均匀分布的补全点云，通过训练子集的PCA特征向量构建核函数，基于GP后验生成均匀采样点。所有操作通过自定义CUDA内核并行加速；(2)协方差中心化设计：估计的协方差用于初始化3D高斯体的缩放和旋转参数加速收敛，同时更新GPR参数；(3)相似性正则化：缓解户外单向运动导致的高斯体向相机方向优化偏差问题；(4)整体框架支持实时密集建图，体素化3DGS表示便于大规模场景管理。其设计目的不是单独增加一个模块，而是用上述表示、监督或渲染约束消除原流程在目标任务中的具体失配。提出首个面向大规模室外场景的实时LiDAR-惯性-视觉 3DGS建图框架，通过Voxel-GPR处理稀疏点云，协方差中心化设计优化高斯体初始化适用时还应保留论文给出的边界：依赖LiDAR传感器增加硬件成本；GPR处理增加计算开销；对极端天气条件的鲁棒性未验证
\textbf{方法信息：} 代码：\url{https://github.com/xieyuser/gs-livm}；GitHub 关注度较高，约 339 stars。



% METHOD: 2411.02703
\paragraph{LVI-GS（紧耦合 LVI 与高斯超基元）（2024.11）}
LVI-GS~\cite{2411_02703} 把 LiDAR、visual 和 inertial SLAM 与 3DGS 通过 hyper-primitive module 连接起来。它不是简单把 LVI 轨迹导出给高斯训练，而是在 odometry thread 和 3DGS optimization thread 之间共享点云、体素和 Gaussians。这样做的目标是让位姿估计和照片级建图都能使用同一套局部几何状态。

高斯初始化来自 colored LiDAR，训练阶段采用 pyramid training、LiDAR depth loss 和 map expansion，CUDA 优化支持在线运行。相对 Gaussian-LIC，它更强调共享地图数据结构；相对 PINGS，它没有引入 SDF 解码，而是以 LiDAR 几何和高斯外观为核心。它适合作为 LVI 系统中“如何把稀疏几何变成可渲染地图”的代表方法。
\textbf{方法信息：} ByteDance 参与了这篇论文。
\textbf{设计动机与边界：} 相对原始 3DGS 以独立高斯和图像重建误差为主的设定，LVI-GS 把问题落在以下可实现的改动上：框架通过两个并行线程实现：(1) 里程计线程——处理LiDAR-Visual-Inertial里程计；(2) 3DGS优化线程——实时优化3D高斯。两者共同维护超原语模块（包含3D点云、体素和3D高斯）。关键技术：(1) 超原语管理——将地图点组织到固定大小体素中，通过体素激活判断最近活动，控制点密度；(2) 金字塔训练——多尺度特征学习提高渲染质量；(3) 深度损失——利用LiDAR测量值提供深度监督信号增强几何感知；(4) 自定义CUDA加速确保实时性能。3D高斯从着色LiDAR点云初始化，通过可微渲染优化。其设计目的不是单独增加一个模块，而是用上述表示、监督或渲染约束消除原流程在目标任务中的具体失配。提出LVI-GS——首个紧耦合LiDAR-Visual-Inertial SLAM框架与3DGS结合。核心创新包括：(1) 从着色LiDAR点云初始化3D高斯；(2) 金字塔训练方法学习多级特征；(3) LiDAR深度损失增强几何特征感知；(4) 高斯地图扩展、关键帧选择、线程管理和CUDA加速的完整策略，实现实时照片级建图。适用时还应保留论文给出的边界：依赖LiDAR硬件；多线程管理和数据同步的复杂性；实时性能可能受场景规模影响。
% METHOD: 2501.08672
\paragraph{GS-LIVO（LiDAR-inertial-visual 高斯建图锚点）（2025.01）}
GS-LIVO~\cite{2501_08672} 是 LiDAR、IMU 和视觉融合路线中的代表性锚点。它继承 LIO 或 LIVO 的紧耦合位姿估计，再把 Gaussian map 作为可渲染稠密地图。单纯 RGB 高斯 SLAM 在户外远距离和低纹理场景中容易失去尺度和几何；GS-LIVO 的改变是用 LiDAR 与 IMU 提供稳定运动和稀疏几何，用视觉光度优化补足外观。

实现上，它维护 hash-indexed octree global Gaussian map，利用 LiDAR 和视觉快速初始化高斯。odometry 侧采用 IESKF 紧耦合融合惯性、LiDAR 和视觉观测；mapping 侧只优化当前可见高斯，并把哈希存储转换成 contiguous GPU memory 以提高渲染和反传效率。项目实现面向实时 LIVO Gaussian mapping。相对室内 RGB-D SplaTAM，GS-LIVO 的关键变化是传感器融合和地图索引：高斯不再只由深度图逐帧生成，而是在全局 LiDAR-inertial 坐标系中维护。
\textbf{实现细节：} GS-LIVO（2025.01） 在上述 Gaussian-LIC 的 LiDAR-视觉融合路线上做了根本性的架构升级——将里程计框架从连续时间因子图替换为迭代误差状态卡尔曼滤波器（IESKF），并用全局八叉树 + GPU 滑动窗口的双层架构取代了 Gaussian-LIC 的单层全量高斯地图，从而首次将实时高斯 SLAM 部署到资源受限的嵌入式平台。其动机直指 Gaussian-LIC 等现有方法的两项瓶颈：一是地图构建线程远慢于位姿估计线程，机器人无法快速感知环境变化；二是高斯基元数量随场景增长导致的 GPU 显存膨胀，使得现有方法在嵌入式平台完全无法运行。全局八叉树地图与双层存储架构。与 Gaussian-LIC 将所有高斯直接存储在 GPU 显存中不同，GS-LIVO 引入了空间哈希索引的递归八叉树作为全局地图的 CPU 端存储结构。根体素通过空间哈希键 $k = \lfloor \mathbf{p} / v_s \rfloor$ 索引，其中 $v_s$ 为根体素边长。八叉树可根据环境复杂度递归细分，天然支持多级细节（LoD）自适应：室内场景采用 0.03\,m 细粒度根体素、户外大规模场景采用 1.0\,m 粗粒度根体素，均支持 2 级细分。这一设计的核心考量在于：CPU 内存（RAM）容量远大于 GPU 显存（VRAM），且可通过交换空间轻松扩展，因此将全局八叉树驻留在 CPU 端，仅在 GPU 端维护当前视场（FoV）内的高斯子集，从根本上将显存占用与场景总规模解耦。LiDAR-相机联合初始化：从深度自适应到表面贴合。Gaussian-LIC 对每个 3D 点创建各向同性高斯（旋转设为单位矩阵），尺度按深度自适应分配。
\textbf{设计动机与边界：} 相对原始 3DGS 以独立高斯和图像重建误差为主的设定，GS-LIVO 把问题落在以下可实现的改动上：GS-LIVO系统架构包含四个核心模块：(1)全局高斯地图：使用哈希索引的八叉树结构，通过空间哈希键高效覆盖稀疏空间体积，支持不同尺度和细节层级的自适应细分。根体素的索引通过空间哈希键计算。(2)高斯快速初始化与在线优化：基于LiDAR和视觉信息进行高斯的快速初始化，使用光度梯度进行在线优化。多传感器融合确保初始化质量。(3)滑动窗口高斯维护：增量维护滑动窗口内的局部高斯，仅优化视场内的高斯，大幅减少GPU计算和显存消耗。将非连续的哈希存储转换为连续GPU内存布局以实现并行优化。(4)IESKF紧耦合里程计：使用迭代误差状态卡尔曼滤波器进行多传感器融合位姿优化，利用高斯地图的实时更新和渲染提供精确的视觉约束。系统在NVIDIA Jetson Orin NX平台上验证了实时性能。其设计目的不是单独增加一个模块，而是用上述表示、监督或渲染约束消除原流程在目标任务中的具体失配。提出首个可在资源受限嵌入式系统（NVIDIA Jetson Orin NX）上实时运行的GS-SLAM框架。设计哈希索引八叉树全局高斯地图和滑动窗口高斯优化，结合IESKF紧耦合多传感器融合里程计，实现实时高保真建图。适用时还应保留论文给出的边界：依赖LiDAR传感器，成本较高；滑动窗口策略可能导致历史区域的细节退化；嵌入式平台上的渲染分辨率受限。
\textbf{方法信息：} 代码：\url{https://github.com/hkust-aerial-robotics/gs-livo}；GitHub 关注度较高，约 452 stars。



% METHOD: 2502.05752
\paragraph{PINGS（SDF 与高斯辐射的点式隐式地图）（2025.02）}
PINGS~\cite{2502_05752} 不把 3DGS 单独作为地图，而是把 point-based implicit neural map、SDF 和 Gaussian radiance 统一起来。它的锚点是神经点地图：每个 neural point 保存位置、方向、几何特征和外观特征；几何特征解码 SDF，外观特征和高斯参数用于照片级渲染。这样做是为了同时得到可优化几何表面和高质量视图合成。

方法的改变在于互相监督。LiDAR 或深度观测约束 SDF 几何，视觉观测通过 Gaussian splatting 约束外观；SDF 提供的表面信息又可以让高斯更贴近几何，渲染误差则帮助补充纯几何地图缺失的纹理。项目实现包含 neural point、SDF 和 Gaussian rendering 模块。相对 GS-LIVO 的显式 octree Gaussian map，PINGS 更像把高斯嵌入到点式隐式地图中，用几何隐式场和显式 splat 各管一部分问题。
\textbf{实现细节：} 实验验证。在轨迹精度上，LIVE-GS 在 Botanic Garden 三个序列上均取得最优 ATE（0.528/0.529/0.498\,m），优于 Gaussian-LIC 和 GS-LIVM，证实 surfel bundle adjustment 对抑制长程漂移的有效性。在鲁棒性上，GS-ICP-SLAM 和 GS-SDF 在 airport03 大规模户外场景中完全失败，而 LIVE-GS 正常运行。PINGS（2025.02） 在统一的弹性点基隐式神经地图内融合有符号距离场（SDF）和高斯泼溅辐射场，通过几何一致性约束实现两种模态的相互增强——SDF 提供精确的表面几何，辐射场提供照片级外观，两者的联合表示支持同时满足导航和渲染的下游需求。
\textbf{实现补充：} 实现上还要核对输入、输出、参数、特征、初始化、训练损失、正则、渲染或推理接口，以及优化中的剪枝、分裂、采样、编码或解码步骤；这些抓手决定该方法到底改了 3DGS 管线的哪一环。
\textbf{方法信息：} 代码：\url{https://github.com/prbonn/pings}；GitHub 关注度较高，约 233 stars。



% METHOD: 2510.26358
\paragraph{AgriGS-SLAM（跨季节果园高斯地图）（2025.10）}
AgriGS-SLAM~\cite{2510_26358} 把多传感器高斯 SLAM 放到农业果园。它的锚点不是通用室内 RGB-D，而是多季节、多相机、强植被变化下的 orchard mapping。方法用 Direct LiDAR Odometry 前端和 second-order motion prior 提供运动估计，用 iSAM2 后端和 Scan Context 回环维护全局一致性。

高斯地图侧，多相机观测共享同一组 Gaussians，并通过 probabilistic LiDAR depth consistency 把 LiDAR 深度反投影到相机视图参与高斯优化。这样，树干、枝条和果园行列的几何由 LiDAR 稳定，颜色和细节由多视角图像补充。项目系统面向 dormant、flower 和 harvest 多季节数据。相对 GS-LIVM 的通用户外建图，AgriGS-SLAM 的增量是把高斯地图维护规则适配农业场景的季节变化和多相机采集。
\textbf{实现细节：} 最后，上述所有方法都假设传感器已预先标定。MM3DGS（2024.04） 挑战了这一前提，首次展示用未标定相机图像和惯性测量实现准确 SLAM——无需预先标定相机内参和相机-IMU 外参，系统在运行过程中在线估计所有标定参数，跟踪精度比当时最佳提升 3 倍。非传统传感器融合。本节讨论雷达、事件相机与热红外等非传统传感器与 3DGS SLAM 的融合方法。
\textbf{实现补充：} 进一步拆开来看，输入侧对应图像、点云、位姿、语义提示或已训练高斯资产；输出侧对应新的高斯参数、渲染图像、编辑结果、地图状态或评测指标。表示层需要说明位置、尺度、旋转、不透明度、特征、材质或 MLP 解码器怎样改；训练层需要说明损失、正则、采样、剪枝、分裂或压缩何时发生；推理层则看渲染、查询、规划或系统接口如何接到原始 3DGS / Scaffold-GS / 4DGS 管线上。
\textbf{设计动机与边界：} 相对原始 3DGS 以独立高斯和图像重建误差为主的设定，AgriGS-SLAM 把问题落在以下可实现的改动上：AgriGS-SLAM由三个核心模块组成：(1) SLAM前端使用Direct LiDAR Odometry (DLO)进行扫描到子图的里程估计，引入二阶运动先验（含加速度）替代恒速假设，增强拖拉机在粗糙地形上的收敛稳定性；(2) SLAM后端维护基于iSAM2的因子图，支持关键帧间里程因子和Scan Context描述子的回环检测，保证全局轨迹一致性；(3) 多视角3DGS模块将多相机视角联合优化，所有相机共享同一组高斯基元，通过批处理光栅化恢复遮挡下的果园结构。位姿优化通过概率LiDAR深度一致性项反向传播到相机投影，紧耦合几何-外观。统一的梯度驱动地图生命周期在关键帧间执行，保留精细细节同时限制内存增长。为了消除原始 3DGS 在该目标任务中的具体失配，方法把上述表示、监督或渲染约束作为可学习和可执行的改动，而不是单独堆叠模块。提出首个面向果园环境的3DGS-SLAM系统，融合LiDAR里程计和多相机3DGS渲染，通过概率LiDAR深度一致性项引导位姿优化，并在休眠期、花期和收获期三个季节部署验证，实现实时高保真重建。适用时还应保留论文给出的边界：依赖LiDAR硬件增加成本；多相机标定需求较高；目前仅在果园环境验证，泛化到其他户外场景需进一步验证。
\textbf{方法信息：} 代码：\url{https://github.com/airlab-polimi/agri-gs-slam}。
% METHOD: 2512.02293
\paragraph{VIGS-SLAM（视觉惯性高斯 SLAM）（2025.12）}
VIGS-SLAM~\cite{2512_02293} 聚焦 visual-inertial Gaussian Splatting SLAM。它继承 DROID-SLAM 式 dense correspondence 和 GRU 更新，但把 IMU preintegration 残差、速度和 bias 一起放入状态优化。单纯 RGB 高斯 SLAM 在快速运动和纹理不足时容易跟踪不稳；VIGS-SLAM 的改变是让惯性约束贯穿初始化、局部优化和回环后的地图更新。

具体流程包括鲁棒 IMU initialization、time-varying bias 建模、基于 GRU 的 correspondence confidence，以及对 poses、disparities、velocity 和 bias 的联合优化。loop closure 触发后，系统同步更新高斯地图而不只是修正轨迹。相对 MM3DGS 的 RGB-D-IMU 路线，VIGS-SLAM 更强调纯视觉惯性输入下的尺度和运动稳定性。
\textbf{实现细节：} 问题背景与设计动机。尽管近年来基于 3DGS 的 SLAM 方法已实现稠密、照片级的地图构建，但纯视觉设计在运动模糊、弱纹理和曝光变化等挑战条件下显著退化。VIGS-SLAM（2025.12） 的核心思路是将 IMU 的惯性约束紧密融入视觉跟踪和高斯建图的统一优化框架，使系统在上述挑战场景中保持鲁棒。具体而言，系统以 RGB 帧序列 $\{I_i\}_{i=1}^N$ 和原始 IMU 测量 $\{\mathbf{a}_k, \boldsymbol{\omega}_k\}_{k=1}^M$（加速度和角速度）为输入，通过两个并行模块协同工作：前端跟踪模块负责相机位姿估计与关键帧管理，高斯泼溅建图模块负责增量构建照片级场景地图。回环闭合检测在后台并行执行，以保证全局一致性。前端跟踪：视觉-惯性联合优化。跟踪模块在线处理图像流，当光流幅度超过阈值时创建新关键帧；为限制 IMU 预积分在长时间间隔上的漂移，系统同时强制每 $t_{kf}$ 秒至少创建一个关键帧。视觉残差项沿用 DROID-SLAM 的学习型 GRU 更新算子，迭代预测关键帧对 $(i, j)$ 之间的稠密对应 $\mathbf{p}^{*}_{ij}$ 及其置信度 $\mathbf{w}_{ij}$，通过最小化重投影误差联合优化位姿和深度： 其中 $E$ 为局部帧图的共视边集合，$\pi$ 和 $\pi^{-1}$ 分别为投影和反投影函数，$\mathbf{u}_i$ 为帧 $i$ 中的 2D 像素坐标，$\| \cdot \|_{\mathbf{W}}$ 为按置信度加权的 Mahalanobis 距离——置信度越高（即对应越可靠）的像素对优化贡献越大。
\textbf{实现补充：} 实现上还要核对输入、输出、参数、特征、初始化、训练损失、正则、渲染或推理接口，以及优化中的剪枝、分裂、采样、编码或解码步骤；这些抓手决定该方法到底改了 3DGS 管线的哪一环。
\textbf{设计动机与边界：} 相对原始 3DGS 以独立高斯和图像重建误差为主的设定，VIGS-SLAM 把问题落在以下可实现的改动上：系统接受RGB帧和IMU测量(角速度、线性加速度)作为输入，同时执行跟踪和高斯映射。(1)IMU初始化：分阶段恢复度量尺度、重力方向和IMU参数。(2)视觉残差：基于DROID-SLAM的GRU更新算子迭代优化帧间对应关系和置信度图，联合优化相机位姿和视差。(3)惯性残差：连续关键帧间IMU预积分提供相对旋转、位置和速度约束，配合Jacobian和协方差传播。时变偏差建模通过连续参数化避免固定偏差假设的局限。(4)联合优化：最小化视觉和惯性残差的加权和，同时优化位姿、视差、速度和IMU偏差。(5)回环检测：检测并执行回环闭合，配合对应高斯更新确保全局一致性。前端关键帧选择基于光流幅度和时间阈值。为了消除原始 3DGS 在该目标任务中的具体失配，方法把上述表示、监督或渲染约束作为可学习和可执行的改动，而不是单独堆叠模块。提出首个紧耦合视觉惯性3DGS SLAM系统VIGS-SLAM，在统一优化框架中联合优化相机位姿、深度和IMU状态。具备鲁棒IMU初始化、时变偏差建模和回环检测(含一致高斯更新)功能。在五个挑战性数据集上优于SOTA方法，在运动模糊、弱纹理和曝光变化条件下保持鲁棒。适用时还应保留论文给出的边界：需要IMU硬件支持；IMU初始化需要一定运动激励；回环检测的计算开销较大；长序列的内存管理需优化。
\textbf{方法信息：} Microsoft 参与了这篇论文。



% METHOD: 2604.13492
\paragraph{RadarSplat-RIO（雷达图像的高斯束调整）（2026.04）}
RadarSplat-RIO~\cite{2604_13492} 把高斯 splatting 推到 radar-inertial odometry。它的锚点不是 RGB 渲染，而是 Doppler-aware radar rendering：给定雷达位姿和高斯地图，系统渲染雷达图像并与真实雷达观测比较。这样 3DGS 的可微渲染思想被迁移到雷达 bundle adjustment。

流程中，MRIO 前端提供初始位姿，后端用 RadarSplat++ 渲染雷达回波并优化 poses 和 map；IMU preintegration 与 Doppler velocity 残差共同约束运动。相对视觉或 LiDAR 高斯 SLAM，RadarSplat-RIO 的关键差异是观测域变了：颜色不再是监督，雷达强度、速度和成像几何成为高斯优化目标。它说明 3DGS 的局部核表示可以进入非视觉 SLAM，但正文中应把它视为传感器迁移案例，而不是 RGB 渲染锚点。
\textbf{实现细节：} RadarSplat-RIO（2026.04） 首次将高斯泼溅的光束调整范式引入雷达 SLAM，填补了这一空白。动机与问题定义：雷达里程计（如 MRIO）的典型管线是：从雷达原始数据经 FFT 生成方位-距离（RA）图像，通过 CFAR 检测提取特征点，利用 Doppler 速度估计雷达自运动速度（将每个特征点的 Doppler 速度建模为自运动速度在视线方向上的投影，通过最小二乘求解），再结合 IMU 预积分进行位姿图优化。这一管线的根本问题在于帧间里程计缺乏全局一致性校正——长序列运行中漂移不可避免，而传统回环闭合需要 revisit 同一地点，限制了纠正时机。核心方法由前端和后端两部分构成。前端沿用 MRIO 的雷达-惯性里程计获取初始位姿，后端执行位姿精化、建图和光束调整三个过程，三者均基于 RadarSplat++ 场景表示的可微渲染。RadarSplat++ 是对 RadarSplat 的关键扩展。每个高斯基元存储 3D 位置 $\mu$、旋转四元数 $q$、缩放向量 $S$ 和雷达散射截面（RCS）属性 $\sigma$（替代传统 3DGS 的 RGB 颜色）。基于雷达方程的渲染过程将 3D 高斯投影到极坐标空间，沿仰角轴进行正交加权平均，根据 RCS 值 $\sigma$ 渲染出 RA 图像。随后，通过天线增益加权和 Doppler 速度的高斯核软分箱，将 RA 空间的功率和 Doppler 图像转换为 RD 图像——这一转换对 RA 图像和 Doppler 速度完全可微，使梯度能够回传至高斯参数和位姿。RadarSplat++ 因此成为首个支持雷达全维度（距离-方位-Doppler）测量的可微渲染管线。后端三个过程依次运行。位姿精化仅优化最新关键帧的位姿 $T_i$，损失函数联合 RA 和 RD 渲染误差：$\mathcal{L}_{\mathrm{pose}} = \sum_{i \in W} \left[ L\!\left(\hat{I}_{RA}, I_{RA}\right) + L\!\left(\hat{I}_{RD}, I_{RD}\right) \right]$，其中自运动速度 $v_i$ 由前端初始化并在优化过程中通过相邻位姿差递归更新。
\textbf{方法信息：} Meta 参与了这篇论文。
\section{回环、退化输入与系统优化}

SLAM 系统最终会遇到比单帧建图更硬的问题：长序列必须回环，模糊或事件输入没有清晰 RGB 帧，单目深度或稠密化会让地图膨胀，跟踪损失在无重叠时可能没有梯度。本节方法都在修补高斯 SLAM 的系统边界：LoopSplat 给全局一致性，Deblur-SLAM 与 MBA-SLAM 处理曝光时间内的运动，IncEventGS 处理事件流，SpectralSplats 改优化目标，RGS-SLAM 与 MVS-GS 改初始化和深度来源。


% METHOD: 2408.10154
\paragraph{LoopSplat（基于高斯子图配准的回环锚点）（2024.08）}
LoopSplat~\cite{2408_10154} 是 3DGS SLAM 回环闭合路线的代表性锚点。许多早期高斯 SLAM 系统只做局部 tracking 和 mapping，长序列中全局漂移会直接固化到高斯地图里。LoopSplat 的改变是把场景拆成 3DGS submaps，并用高斯子图之间的直接配准来产生回环约束，而不是把高斯退化成普通点云后再配准。

系统在局部阶段做 frame-to-model tracking 和 dense mapping；当相机运动超过阈值时，当前子图固定并开启新子图。后端持续检测是否回到已访问区域，一旦发现回环，就选择相关 3DGS 子图，通过子图间直接 registration 得到相对位姿约束。随后鲁棒 pose graph optimization 计算全局一致轨迹，并把各个子图刚性对齐。代码 README 说明其前端、后端和 PGO 都围绕 Gaussian splat submaps 组织。相对 SplaTAM 的在线局部地图，LoopSplat 的关键贡献是把高斯地图纳入 pose graph 全局一致性，使高保真渲染不再只依赖前端位姿不漂移。
\textbf{锚点细化：} LoopSplat：针对缺乏回环检测的 SLAM 在长序列必然累积漂移、此前方法回环约束依赖传统点云配准无法利用场景表示本身信息的局限，用\textbf{直接在 3DGS 子地图间配准提取回环约束}替代 FPFH + RANSAC + ICP 点云配准。核心思想是\textbf{将 3DGS 作为跟踪、建图和全局一致性维护的统一场景表示}——所谓\textbf{子地图}（submap）是将大场景按空间或时间划分为若干局部地图单元（每个子地图独立优化、子地图间通过回环约束对齐），既降低单次优化规模又便于分布式处理；LoopSplat 是首个将 3DGS 同时用于跟踪、建图和回环闭合三个环节的耦合式 RGB-D SLAM 系统，回环约束不再借用点云配准，而是直接对齐两幅 3DGS 子地图的渲染图像和深度图。
\textbf{实现细节：} LoopSplat（2024.08） 是首个将 3DGS 表示同时用于跟踪、建图和回环闭合的耦合式 RGB-D SLAM 系统。其出发点是一个根本性问题：能否直接利用地图表示（即 3DGS）本身完成回环闭合，而无需回退到传统点云配准？此前的方法各有局限——GO-SLAM 需要代价高昂的哈希网格特征重训练来变形地图，Photo-SLAM 需要额外的 3DGS 参数优化来消化 ORB-SLAM 追踪器的位姿更新，而最相近的 Loopy-SLAM 虽采用子地图避免存储所有帧，但其回环约束仍依赖传统的 FPFH + RANSAC + ICP 全局到局部点云配准流水线，既缓慢又无法利用场景表示本身的丰富信息。LoopSplat 的回答是肯定的——它提出直接在 3DGS 子地图间进行配准以提取回环约束，从而将 3DGS 作为跟踪、建图和全局一致性维护的统一场景表示。在 LoopSplat 中，场景以子地图（submap）为单位组织。每个子地图以 3DGS 点云 $P^s = \{(\mu_i, \Sigma_i, o_i, C_i)\}$ 建模若干关键帧所观测的局部区域。子地图的划分采用动态触发策略：当当前帧相对子地图首帧的位移或旋转超过阈值 $d_{\max}$ 或 $\theta_{\max}$ 时初始化新子地图，避免了固定关键帧数策略可能导致的子地图过大或过小。前端采用帧到模型跟踪——先按匀速运动假设初始化位姿，再最小化渲染颜色/深度与输入的光度-深度联合损失优化位姿，并用 alpha 掩码和内点掩码剔除重建不完整区域的粗差。LoopSplat 的核心创新是3DGS 子地图间的直接配准，将回环约束的计算分解为三个阶段。重叠估计。两个子地图 $P$ 和 $Q$ 之间的近似重叠是鲁棒配准的前提。作者发现直接在高斯均值点云上匹配局部特征效果很差——高斯中心分布不均匀且缺乏判别性描述子。
\textbf{锚点细化：} LoopSplat：针对缺乏回环检测的 SLAM 在长序列必然累积漂移、此前方法回环约束依赖传统点云配准无法利用场景表示本身信息的局限，用\textbf{直接在 3DGS 子地图间配准提取回环约束}替代 FPFH + RANSAC + ICP 点云配准，核心思想是\textbf{将 3DGS 作为跟踪、建图和全局一致性维护的统一场景表示}——所谓\textbf{子地图}（submap）是将大场景按空间或时间划分为若干局部地图单元（每个子地图独立优化、子地图间通过回环约束对齐），既降低单次优化规模又便于分布式处理；LoopSplat 是首个将 3DGS 同时用于跟踪、建图和回环闭合三个环节的耦合式 RGB-D SLAM 系统，回环约束不再借用点云配准，而是直接对齐两幅 3DGS 子地图的渲染图像和深度图。
\textbf{方法信息：} 代码：\url{https://github.com/gradientspaces/loopsplat}；GitHub 关注度较高，约 391 stars；引用很多，Semantic Scholar 约 84 次。



% METHOD: 2410.08107
\paragraph{IncEventGS（无位姿事件相机增量重建）（2024.10）}
IncEventGS~\cite{2410_08107} 处理的是单个事件相机、无姿态先验的 3DGS 重建。事件相机没有常规 RGB 帧，观测是异步亮度变化；原始 3DGS 的图像重建损失不能直接使用。IncEventGS 的改变是把事件流分成时间窗口，并为每个窗口建立连续时间 $SE(3)$ 轨迹参数化，再用高斯渲染的亮度变化和事件测量对齐。

tracking 阶段，系统从事件块中采样两个邻近时间戳，固定当前地图，只优化该事件块的相机运动。mapping 阶段，在滑动窗口内联合优化高斯参数和相机轨迹，并根据新探索区域添加高斯、修剪透明高斯。项目实现围绕 event-based Gaussian splatting 展开。相对 RGB SLAM，IncEventGS 的局部锚点是连续时间事件建模：它把“图像重建”改成“亮度变化重建”，从而让高斯可微渲染服务于事件流 SLAM。
\textbf{实现细节：} IncEventGS（2024.10） 在 EGS-SLAM 的基础上进一步去掉了 RGB-D 帧和深度传感器，仅凭单个事件相机完成增量式 3DGS 重建与位姿估计。去掉 RGB-D 后，IncEventGS 面临双重难题：既无密集图像帧可提供绝对亮度监督，也无深度观测可约束几何——跟踪和建图必须完全依赖事件流的亮度变化信号驱动。事件 chunk 与对数亮度变化监督。与 EGS-SLAM 利用模糊图像建模（曝光区间内 $K$ 个采样点渲染积分）不同，IncEventGS 将事件流按 50\,ms 固定窗口切分为 chunk，每个 chunk 仅用两个端点 SE(3) 位姿 $\mathbf{T}_\mathrm{start}$、$\mathbf{T}_\mathrm{end}$ 通过 $\mathfrak{se}(3)$ 空间线性插值参数化连续轨迹。在每个 chunk 内随机采样两个邻近时刻 $t_k$ 和 $t_k{+}\Delta t$，从 3DGS 渲染对应亮度图像 $\hat{I}(t_k)$、$\hat{I}(t_k{+}\Delta t)$，合成对数亮度变化 $\hat{\mathbf{E}}(\mathbf{x}) = \log \hat{I}(t_k{+}\Delta t, \mathbf{x}) - \log \hat{I}(t_k, \mathbf{x})$，与事件流累积测量的亮度变化 $\mathbf{E}(\mathbf{x}) = C \cdot \sum_{t_k < t_i < t_k{+}\Delta t} p_i$（$C$ 为对比度阈值）求 $\ell_2$ 距离作为监督。这里的关键差异是：EGS-SLAM 的亮度变化来自灰度图像的 BT.601 变换后对数差分（事件图与模糊帧图的混合监督），而 IncEventGS 的亮度变化完全从 3DGS 渲染的灰度图直接取对数差分，无需任何帧式输入。增量跟踪-建图与滑动窗口 BA。IncEventGS 采用了经典 SLAM 的 tracking-mapping 交替范式，而非 EGS-SLAM 所使用的跟踪（模糊帧匹配）+ 建图（联合优化高斯与关键帧位姿）的分离流程。
\textbf{设计动机与边界：} 相对原始 3DGS 以独立高斯和图像重建误差为主的设定，IncEventGS 把问题落在以下可实现的改动上：IncEventGS采用SLAM式的tracking-mapping范式：(1)事件处理：将异步事件流按固定时间窗口分块，为每个块关联SE(3)空间中的连续时间轨迹参数化；(2)Tracking阶段：从事件块中采样两个邻近时间戳，利用已有3DGS渲染亮度变化并与事件测量对比，仅优化当前事件块对应的相机运动轨迹；(3)Mapping阶段：利用tracking阶段提供的初始轨迹，在滑动窗口内联合优化3DGS参数（高斯位置、协方差、不透明度、颜色）和相机轨迹。通过连续时间轨迹参数化可从事件流中任意采样时间点计算亮度变化。密度控制策略自适应增加新探索区域的高斯并修剪透明高斯。其设计目的不是单独增加一个模块，而是用上述表示、监督或渲染约束消除原流程在目标任务中的具体失配。提出首个基于单个事件相机的无姿态先验3DGS增量重建算法IncEventGS，通过tracking-mapping范式联合恢复相机轨迹和密集3D场景表示，性能优于需要GT姿态的NeRF-based方法适用时还应保留论文给出的边界：事件数据的稀疏性限制了纹理贫乏区域的重建质量；滑动窗口BA可能丢失长期约束；对快速相机运动的鲁棒性有待提升
\textbf{方法信息：} 代码：\url{https://github.com/wu-cvgl/inceventgs}；GitHub 关注度较高，约 117 stars。



% METHOD: 2411.08279
\paragraph{MBA-SLAM（运动模糊感知 bundle adjustment）（2024.11）}
MBA-SLAM~\cite{2411_08279} 同样处理运动模糊，但它把方法抽象成 motion blur aware tracking 和 bundle adjustment，可以服务 NeRF 或 3DGS 地图。它的锚点是物理运动模糊成像：每帧不再只有一个位姿，而由曝光开始和结束的连续 $SE(3)$ 轨迹表示。

追踪器估计每个模糊帧曝光期间的相机运动；建图器通过渲染多个虚拟清晰图像并积分成合成模糊图像，与真实模糊图像做光度一致性，联合优化场景表示和相机轨迹。项目实现包含 motion blur aware Gaussian SLAM。相对 Deblur-SLAM，MBA-SLAM 更像通用 BA 框架：核心不是先去模糊再跟踪，而是在 tracking 和 mapping 目标函数里直接承认曝光时间内的运动。
\textbf{实现细节：} 手持设备或快速运动场景中的运动模糊是 3DGS-SLAM 的主要退化因素，其影响通过两个机制恶性循环：运动模糊破坏了密集 SLAM 所依赖的光度一致性假设，导致位姿估计不准；不准确的位姿又引入多视角几何不一致，进一步恶化地图重建质量。MBA-SLAM（2024.11） 是首个系统性地将运动模糊的物理成像过程嵌入完整 RGB-D SLAM 管线的方法，其核心思路可用一句话概括：与其从模糊图像中去除模糊，不如从清晰渲染中合成模糊。物理模型与问题定义。运动模糊图像 $I(\mathbf{u})$ 本质上是曝光时间 $\tau$ 内所有虚拟清晰图像 $\hat{I}_t(\mathbf{u})$ 的积分： 实践中用离散近似：对 $n$ 个均匀采样的虚拟位姿分别渲染清晰图像 $\hat{I}_i(\mathbf{u})$，取平均得到合成模糊图像 $I(\mathbf{u}) \approx \frac{1}{n}\sum_{i=0}^{n-1} \hat{I}_i(\mathbf{u})$。前端跟踪：运动模糊感知的直接对齐。传统直接法跟踪假设参考帧与当前帧满足光度一致性，但运动模糊使得对应点外观差异巨大，这一假设完全失效。MBA-SLAM 选择重模糊（re-blurring）而非去模糊作为应对策略——重模糊只需沿轨迹渲染并取平均，而去模糊需要从严重退化图像中恢复细节，前者在数学上更良态且对强模糊更具鲁棒性。
\textbf{实现补充：} 实现上还要核对输入、输出、参数、特征、初始化、训练损失、正则、渲染或推理接口，以及优化中的剪枝、分裂、采样、编码或解码步骤；这些抓手决定该方法到底改了 3DGS 管线的哪一环。
\textbf{方法信息：} 代码：\url{https://github.com/wu-cvgl/mba-slam}。
% METHOD: 2412.19130
\paragraph{MVS-GS（在线 MVS 深度驱动的高质量建图）（2024.12）}
MVS-GS~\cite{2412_19130} 的问题是：RGB 图像流在线建图时，常规 SLAM 深度粗糙，直接初始化 3DGS 会限制最终渲染质量。它的改变是把在线 MVS 深度估计接入高斯建图。前端使用 DROID-SLAM 跟踪位姿，并在局部时间窗口内用 MVSFormer 估计关键帧高分辨率深度。

这些深度不会直接使用，而是经过深度一致性检查、尺度对齐和离群值过滤。后端把精化深度转换成点云初始化高斯，并持续优化高斯参数。前端和后端并行运行，保证新关键帧到来时地图能及时更新。项目实现围绕 online MVS 和 Gaussian mapping。相对 RGS-SLAM 用特征匹配一次性初始化，MVS-GS 借助 learned MVS 产生更 dense 的深度图，重点在提高 RGB-only 在线建图的几何质量。
\textbf{实现细节：} 与 DenseSplat 和 LoopSplat 采用运动立体或单目深度估计不同，MVS-GS（2024.12） 将预训练的多视角立体网络 MVSFormer 引入在线 SLAM 管线，从局部时间窗口内的相邻关键帧直接估计高分辨率深度图和置信图。然而，在线 MVS 因源视角有限而产生大量离群点（尤其出现在背景和边界区域），为此 MVS-GS 设计了两级深度精化管线：先用 V-Fuse 深度融合网络将参考帧与邻域帧的深度-置信图融合为时序一致的精化深度图，再通过多视角几何一致性检查——将参考帧深度投影到邻域帧并计算相对深度差异，剔除未通过至少三视图一致性检验的像素——进一步滤除离群点。精化后的稠密深度图直接用于高斯点初始化，避免了 Photo-SLAM 等方法从稀疏特征扩展高斯的低质量问题。系统架构上，前端使用 DROID-SLAM 的稠密光束法平差进行跟踪，后端以独立线程并行运行，采用广义指数 splatting（GES）替代标准 3DGS 表示以提升大场景下的效率，且仅对深度图中未探索区域生成新高斯，避免重复初始化。
\textbf{方法信息：} 代码：\url{https://github.com/lbg030/mvs-gs}。
\textbf{设计动机与边界：} 相对原始 3DGS 以独立高斯和图像重建误差为主的设定，MVS-GS 把问题落在以下可实现的改动上：系统分为前端和后端两个并行线程。前端使用DROID-SLAM跟踪相机位姿，利用MVSFormer在局部时间窗口内估计关键帧的高分辨率深度图，并通过深度一致性检查、深度对齐和离群值过滤等综合精化技术提升深度质量。后端将精化后的深度图转换为点云初始化3DGS高斯基元，并持续优化高斯参数。前端和后端并行运行，确保每来一个新关键帧时3DGS模型及时更新。其设计目的不是单独增加一个模块，而是用上述表示、监督或渲染约束消除原流程在目标任务中的具体失配。首次将在线MVS深度估计与3DGS建模深度结合，通过MVSFormer生成高分辨率深度图并综合精化，实现高质量3DGS在线建模，在挑战性户外环境中显著超越现有密集SLAM方法。适用时还应保留论文给出的边界：MVS深度估计在纹理缺失区域可能失效；系统需要GPU支撑MVSFormer的实时推理；快速运动场景下深度估计质量可能下降。
% METHOD: 2503.12572
\paragraph{Deblur Gaussian Splatting SLAM（运动模糊的子帧建模）（2025.03）}
Deblur Gaussian Splatting SLAM~\cite{2503_12572} 面向运动模糊输入。普通 Gaussian SLAM 假设每帧对应一个瞬时相机位姿，但真实模糊图像是曝光时间内多个位姿渲染结果的积分；若直接跟踪模糊帧，位姿和地图都会被拉偏。方法的改变是桥接 frame-to-frame 与 frame-to-model：前者提供初始相机位姿和高斯初始化，后者学习曝光期间的子帧轨迹。

具体做法是把模糊图像近似成多个虚拟清晰帧的平均，通过 LERP 和 SLERP 插值得到曝光时间内的虚拟相机，再最小化观测模糊图像与渲染虚拟帧平均值之间的误差。系统还引入 OmniDepth 之类的单目深度先验，并使用可变形 3DGS 保持全局一致性。相对 MBA-SLAM，Deblur-SLAM 更明确地把 Splat-SLAM 式 frame-to-frame 轨迹、frame-to-model 子帧优化和模糊分解三者拼接成 RGB 高斯 SLAM 管线。
\textbf{方法信息：} Microsoft 参与了这篇论文。
\textbf{设计动机与边界：} 相对原始 3DGS 以独立高斯和图像重建误差为主的设定，Deblur Gaussian Splatting SLAM 把问题落在以下可实现的改动上：Deblur-SLAM桥接帧到帧和帧到模型方法的优势：(1) 利用SplatSLAM的帧到帧追踪提供精确的相机位姿和高斯初始化；(2) 叠加帧到模型追踪器学习子帧轨迹，在子帧精度渲染清晰图像（虚拟图像）；(3) 运动模糊物理建模：将运动模糊图像建模为曝光时间内清晰图像的积分，通过M个虚拟相机线性插值（LERP/SLERP）近似离散化，最小化观测模糊图像与渲染虚拟帧平均值的误差；(4) 采用可变形3DGS表示实现全局一致性，支持在线变形优化；(5) 集成单目深度估计器(OmniDepth)提供深度先验。三个核心组件（帧到帧追踪器、帧到模型追踪器、模糊分解）协同工作。其设计目的不是单独增加一个模块，而是用上述表示、监督或渲染约束消除原流程在目标任务中的具体失配。提出Deblur-SLAM，首个鲁棒的RGB SLAM管线，从运动模糊输入恢复清晰重建。创新点：(1) 结合帧到帧和帧到模型方法建模子帧相机轨迹实现高保真重建；(2) 集成在线回环闭合和全局BA实现密集精确的全局轨迹；(3) 物理运动模糊图像形成建模，最小化观测模糊图像与渲染模糊图像的误差；(4) 利用单目深度估计器和在线高斯变形确保精确建图。适用时还应保留论文给出的边界：子帧虚拟视图数量M影响去模糊质量和计算量；依赖单目深度估计器的精度；回环闭合检测在高度动态场景中可能失效。
% METHOD: 2601.00705
\paragraph{RGS-SLAM（一次性密集初始化替代残差增密）（2026.01）}
RGS-SLAM~\cite{2601_00705} 重新审视 3DGS SLAM 的密度控制。原始 3DGS 和许多 SLAM 版本依赖残差驱动 densification，在线场景中这会带来慢收敛、拓扑抖动和冗余高斯。RGS-SLAM 的改变是用 training-free correspondence-to-Gaussian initialization 做一次性密集三角化，先生成空间均衡、结构感知的高斯集合，再让后续优化只精炼属性。

具体流程在关键帧窗口内使用 DINOv3 描述子做密集特征匹配，通过 confidence-aware inlier classifier 过滤错误对应，然后一次性三角化得到高斯种子。之后优化均值、协方差、不透明度和颜色，但保持拓扑相对固定。项目实现可替换 MonoGS 的 densification 模块。相对 MotionGS 的压缩剪枝，RGS-SLAM 更强调初始化阶段：先给地图一个可靠高斯拓扑，再减少在线残差增密带来的不稳定。
\textbf{实现细节：} 从"绕过可微渲染"回到可微渲染管线内部，RGS-SLAM（2026.01） 保留 SplaTAM 的可微渲染管线与关键帧管理框架，但将核心致密化机制从"逐帧残差驱动增密"替换为"DINOv3 密集特征匹配 $+$ 多视角三角化的一次性初始化"。动机。SplaTAM 的残差驱动致密化在稀疏视角下收敛缓慢且容易产生空间不一致——每帧仅根据当前渲染残差局部添加高斯，缺乏全局视野；而基于语义特征的三角化能为后续光度优化提供更稳定、覆盖更均匀的初始高斯分布。具体而言，每个关键帧触发时，用 DINOv3 描述子建立密集多视角对应，经置信度过滤后批量三角化出约 1000 个各向异性椭球高斯种子（而非 SplaTAM 的各向同性球状高斯），一次插入地图，后续仅做轻量合并与剪枝，不再迭代增密。这一设计的动机是：残差驱动增密在稀疏视角下收敛缓慢且容易产生空间不一致，而基于语义特征的三角化能为后续光度优化提供更稳定、覆盖更均匀的初始高斯分布。在跟踪端，RGS-SLAM 放弃了 SplaTAM 的轮廓掩码和深度损失，仅用光度损失（加仿射亮度模型补偿曝光变化），并采用解析 $SE(3)$ Jacobian 替代通用自动微分加速位姿更新。
\textbf{设计动机与边界：} 相对原始 3DGS 以独立高斯和图像重建误差为主的设定，RGS-SLAM 把问题落在以下可实现的改动上：方法核心是单步密集初始化替代渐进密集化：(1) 在关键帧窗口内使用密集特征匹配（DFM），通过DINOv3描述子获取密集多视角对应；(2) 置信度感知内点分类器精炼对应质量；(3) 一次性三角化生成结构感知的高斯种子集合。后续优化仅精炼均值、协方差、不透明度和颜色，保持拓扑固定，实现稳定平稳的优化。直接替换MonoGS的密集化模块，完全兼容现有GS-SLAM流水线。为了消除原始 3DGS 在该目标任务中的具体失配，方法把上述表示、监督或渲染约束作为可学习和可执行的改动，而不是单独堆叠模块。用训练无关的对应-到-高斯初始化替代GS-SLAM的残差驱动密集化。通过DINOv3特征匹配+置信度感知内点分类实现一次性密集三角化，生成空间均衡、结构感知的高斯集合。收敛加速约20%，位姿漂移减少30%以上，渲染吞吐量提升20%，达到925 FPS实时建图。适用时还应保留论文给出的边界：依赖DINOv3特征质量；特征匹配在弱纹理区域可能失败；对快速运动的鲁棒性待验证。
\textbf{方法信息：} 代码：\url{https://github.com/breeze1124/rgs-slam}。



% METHOD: 2603.24036
\paragraph{SpectralSplats（频谱矩监督的可微跟踪）（2026.03）}
SpectralSplats~\cite{2603_24036} 关注的是 3DGS tracking 的梯度消失。当渲染高斯和目标结构在空间上完全不重叠时，普通像素域 L2 损失没有有效方向，优化器不知道该把高斯往哪里推。它的改变是把空间域监督换成 frequency-domain spectral moments，用全局复正弦特征给渲染图提供方向性梯度。

方法推导了空间光度误差中 self-term 与 target supervision 的分解，说明无重叠时 target supervision 项可能为零。频谱矩通过低频到高频的全局分量扩大吸引域，frequency annealing 先用低频形成平滑目标，再逐步引入高频实现精确对齐。相对 MipSLAM 的抗混叠，SpectralSplats 不只是减少采样噪声，而是直接改 tracking loss 的形状，解决优化早期或大位姿误差下没有梯度的问题。
\textbf{方法信息：} NVIDIA 参与了这篇论文。
\textbf{设计动机与边界：} 相对原始 3DGS 以独立高斯和图像重建误差为主的设定，SpectralSplats 把问题落在以下可实现的改动上：方法核心是将标准空间L2损失替换为频谱目标。当高斯原语的渲染足迹与目标结构无空间重叠时，标准空间损失的梯度为零。SpectralSplats使用一组全局复正弦特征（频谱矩）监督渲染图像，构建全局吸引域。数学推导：将光度误差分解为Self-Term和Target Supervision，证明当渲染高斯与目标无重叠时Target Supervision项为零。频谱矩通过全局频率分量提供方向性梯度。频率退火策略从低频（全局凸性）逐渐过渡到高频（精确空间对齐），避免高频引入的周期性局部最小值。其设计目的不是单独增加一个模块，而是用上述表示、监督或渲染约束消除原流程在目标任务中的具体失配。提出了SpectralSplats框架，通过将优化目标从空间域转移到频域，解决了3DGS跟踪中的梯度消失问题。引入频谱矩构建全局吸引域，确保即使像素重叠完全不存在时也存在有效梯度。推导了频率退火策略从全局凸性平滑过渡到精确空间对齐。适用时还应保留论文给出的边界：频谱变换增加计算开销；频率退火调度需要调参。
% METHOD: 2603.26810
\paragraph{Unblur-SLAM（运动与散焦模糊的双阶段处理）（2026.03）}
Unblur-SLAM~\cite{2603_26810} 处理比 Deblur-SLAM 更宽的退化输入：运动模糊和散焦模糊。它的改变是双阶段策略。第一阶段用前馈去模糊网络把大多数模糊帧变成可供 Droid-SLAM 和单目深度网络使用的图像；第二阶段对去模糊失败的帧，不强行交给前端，而是通过全局 3DGS 表示和 blur MLP 在三维空间中模拟模糊形成。

系统先构建 blur detection benchmark，选择可靠指标把帧分成清晰、成功去模糊和失败去模糊三类。成功帧进入常规 tracking 和高斯形变优化；失败帧则跳过易出错的跟踪器，通过 MLP 生成多个子帧并平均渲染，计算与真实模糊帧的光度损失。相对 Deblur-SLAM 主要建模运动模糊曝光积分，Unblur-SLAM 增加了模糊识别、物理约束去模糊训练和自适应计算量调节，更像面向复杂模糊输入的系统级补丁。
\textbf{实现细节：} Unblur-SLAM（2026.03） 从三个维度对 MBA-SLAM 的运动模糊范式进行了根本性拓展：将模糊类型从单一运动模糊扩展到运动模糊 + 散焦模糊，将固定计算量改为自适应于模糊程度，将纯物理轨迹建模改为前馈去模糊网络 + 逐像素模糊核的混合管线。系统基于 SplatSLAM 的纯 RGB 框架构建，含回环闭合和全局光束法平差。
\textbf{技术实现补充：} Unblur-SLAM（2026.03） 从三个维度对 MBA-SLAM 的运动模糊范式进行了根本性拓展：将模糊类型从单一运动模糊扩展到运动模糊 + 散焦模糊，将固定计算量改为自适应于模糊程度，将纯物理轨迹建模改为前馈去模糊网络 + 逐像素模糊核的混合管线。系统基于 SplatSLAM 的纯 RGB 框架构建，含回环闭合和全局光束法平差。
\textbf{方法信息：} Meta 参与了这篇论文；代码仓库：\url{https://github.com/slammate/unblur-slam}，当前未确认完整实现。
\textbf{设计动机与边界：} 相对原始 3DGS 以独立高斯和图像重建误差为主的设定，Unblur-SLAM 把问题落在以下可实现的改动上：Unblur-SLAM包含以下核心模块：(1)模糊量化与识别：构建综合模糊检测基准（含真实和半合成数据集），评估39种图像质量指标，选择最优模糊检测器将帧分为三类：清晰帧、成功去模糊帧、失败去模糊帧。(2)前馈去模糊：提出物理约束的训练方案训练3D一致的单图去模糊模型，增强Droid-SLAM跟踪器和单目深度估计网络的模糊鲁棒性。(3)全局3DGS表示：结合Splat-SLAM的在线形变和渲染优化。对去模糊成功的帧通过Droid-SLAM处理后用高斯形变和光度/深度误差优化；对去模糊失败的帧跳过跟踪器，通过MLP建模多个子帧的平均来计算与运动模糊帧的光度损失。(4)局部-全局多视图优化和回环检测。其设计目的不是单独增加一个模块，而是用上述表示、监督或渲染约束消除原流程在目标任务中的具体失配。提出首个能同时处理运动模糊和散焦模糊的RGB SLAM系统。设计双阶段策略：第一阶段用前馈去模糊网络处理模糊帧；第二阶段对去模糊失败的帧通过3DGS全局表示和模糊MLP建模多个模糊子帧来模拟3D空间中的模糊形成过程。还提出基于模糊量的自适应计算量调节机制。适用时还应保留论文给出的边界：依赖去模糊基础模型的质量，对极端模糊可能仍然失败，模糊MLP增加了优化复杂度，回环检测可能受模糊影响。

\section{视觉定位与重定位}

定位方法把高斯地图的用途从“生成新视角”转为“帮助求 6-DoF 位姿”。这里有两条路线：一条直接优化渲染深度或颜色残差，另一条把描述符、地标概率或 2D-3D 特征匹配嵌入高斯。后者往往比纯光度定位更能处理较大的初始误差，因为它先建立可验证对应，再做 PnP 或局部精化。
% METHOD: 2409.14067
\paragraph{SplatLoc（地标选择与 3D 描述符定位锚点）（2024.09）}
SplatLoc~\cite{2409_14067} 是 3DGS 视觉定位中的代表性锚点。它继承原始 3DGS 的可渲染场景表示，但给每个高斯额外关联 3D descriptor 和 landmark probability。普通 3DGS 能渲染图像，却不直接告诉定位系统哪些 splats 是稳定地标；SplatLoc 的改变是把高斯从外观原语升级成可匹配、可筛选的 3D 地标。

方法先从特征体蒸馏 scene-specific descriptor 到高斯，使查询图像中的 2D 特征能和 3D 高斯建立对应。随后通过显著 3D 地标选择保留定位贡献大的高斯子集，同时起到模型压缩作用。各向异性正则约束关键高斯不要过度拉长，避免地标形状影响定位精度。公开代码和复核材料说明其可验证性较强。相对 GSplatLoc 的深度残差优化，SplatLoc 更偏 feature-based localization：先找到可靠高斯地标，再求相机位姿。
\textbf{锚点细化：} SplatLoc：面向 AR 应用场景，在 STDLoc 的 Feature Gaussian 框架基础上重新审视了特征学习和地标选取两个核心环节。STDLoc 通过 $\alpha$-混合光栅化将高斯特征渲染为 2D 特征图并与真值对齐，这一机制虽然有效，却引入了描述符偏差——$\alpha$-混合将单个高斯的特征与邻近高斯纠缠，优化后的特征并不严格对应该高斯本身的视觉身份。SplatLoc 的第一项改进是学习更接近单个高斯身份的无偏 3D descriptor field；第二项改进是为每个高斯学习地标概率，用显著 3D landmarks 支撑 2D--3D matching 和 PnP；第三项改进是对关键高斯施加各向异性正则，避免被选作地标的高斯拉成长条后带来定位几何误差。推理时，查询图像先提取 2D 特征，与高斯描述符建立对应，再用地标子集估计 6-DoF 位姿，因此它不是单纯渲染残差定位，而是“高斯地图 + 可匹配地标”的定位路线。
\textbf{实现细节：} SplatLoc（2024.09） 面向 AR 应用场景，在 STDLoc 的 Feature Gaussian 框架基础上重新审视了特征学习和地标选取两个核心环节。STDLoc 通过 $\alpha$-混合光栅化将高斯特征渲染为 2D 特征图并与真值对齐，这一机制虽然有效，却引入了描述符偏差——$\alpha$-混合将单个高斯的特征与邻近高斯纠缠，优化后的特征并不严格对应该高斯本身的视觉身份。SplatLoc 从这一观察出发，提出三项针对 STDLoc 的增量改进。无偏描述符学习——绕过 $\alpha$-混合的特征体蒸馏。回顾 STDLoc 为每个高斯直接存储特征属性 $f_i$ 并通过 $\alpha$-混合渲染优化。SplatLoc 完全放弃这一路线，转而将多视角 2D 特征图投影融合到 3D 特征体 $\mathcal{V}$ 中，再用场景专用的 3D 解码器（多分辨率哈希编码 + 浅层 MLP）从 3D 位置直接回归高维描述符——不经过任何 $\alpha$-混合，因此是无偏的。蒸馏仅在场景表面流形上执行（通过 Marching Cubes 提取表面后采样），避免在空洞区域学习无效描述符。这一设计使描述符学习与高斯基元参数解耦，模型存储仅需 112MB，而同样基于点式特征存储的 PNeRFLoc 需要 788MB——减少约 7 倍。显著性地标选择——三维评分替代匹配评分。STDLoc 用多视角匹配评分 $s(g_i)$ 评价高斯质量，再通过最远点采样保证空间覆盖。SplatLoc 重新定义了更全面的三维显著性评分 $\mathcal{H}(p)$，融合三项指标：显著性（可学习的地标概率分数）、泛化性（训练过程中被观察到的最大视角张角）和几何一致性（高斯中心到多视角表面点的距离误差）。
\textbf{锚点细化：} SplatLoc：面向 AR 应用场景，在 STDLoc 的 Feature Gaussian 框架基础上重新审视了特征学习和地标选取两个核心环节。STDLoc 通过 $\alpha$-混合光栅化将高斯特征渲染为 2D 特征图并与真值对齐，这一机制虽然有效，却引入了描述符偏差——$\alpha$-混合将单个高斯的特征与邻近高斯纠缠，优化后的特征并不严格对应该高斯本身的视觉身份。SplatLoc 从这一观察出发，提出三项针对 STDLoc 的增量改进。
\textbf{设计动机与边界：} 相对原始 3DGS 以独立高斯和图像重建误差为主的设定，SplatLoc 把问题落在以下可实现的改动上：核心组件：(1) 3D高斯场景表示：标准3DGS表示场景，每个高斯额外关联3D描述符和地标概率分数；(2) 无偏3D描述符特征场学习：从构建的特征体中蒸馏场景特定描述符到高斯基元，确保精确的2D-3D对应用于位姿估计；(3) 显著3D地标选择：基于显著性分数选择合适的高斯子集用于定位，同时实现模型压缩；(4) 各向异性正则化：约束关键高斯基元的形状，防止过度拉长影响定位性能，同时不影响渲染质量；(5) 增量重建流程。其设计目的不是单独增加一个模块，而是用上述表示、监督或渲染约束消除原流程在目标任务中的具体失配。提出基于3DGS的视觉定位方法SplatLoc，实现高质量渲染和精确6-DoF定位的双重能力，通过显著地标选择和各向异性正则化提升定位精度和地图紧凑性适用时还应保留论文给出的边界：大规模场景的可扩展性；动态场景支持有限；对初始SfM质量有依赖
\textbf{方法信息：} Microsoft 参与了这篇论文；代码：\url{https://github.com/zhaihongjia/splatloc}；GitHub 关注度较高，约 145 stars；引用很多，Semantic Scholar 约 60 次。



% METHOD: 2409.16502
\paragraph{GSplatLoc（XFeat 描述符蒸馏定位）（2024.09）}
另一篇 GSplatLoc~\cite{2409_16502} 走的是描述符嵌入路线。它把 XFeat dense keypoint descriptors 蒸馏到 3DGS 中，使每个高斯除颜色外还携带局部特征。普通光度定位容易被曝光、纹理重复或初始位姿影响；该方法的改变是先通过 2D-3D cosine similarity 匹配得到粗位姿，再用渲染特征和光度变形损失做精化。

训练时，系统用颜色损失、深度损失和特征蒸馏损失共同优化高斯。推理时，查询图像的 XFeat 特征与高斯特征匹配，经 PnP 和 RANSAC 得到初始位姿；随后渲染特征图做关键点匹配精化，再用光度 deformation loss 微调。其定位流程可以复核。相对 SplatLoc，它的差异在于具体采用 XFeat 和两阶段特征加光度精化，而不是以显著地标选择为核心。
\textbf{设计动机与边界：} 相对原始 3DGS 以独立高斯和图像重建误差为主的设定，GSplatLoc 把问题落在以下可实现的改动上：两阶段流程：(1)场景建模：使用基于特征的3DGS方法，结合XFeat深度特征提取器训练3DGS模型。XFeat生成(W/8)*(H/8)*64特征图，双线性上采样到原始分辨率。通过最小化包含颜色损失、深度损失和特征蒸馏损失的联合损失函数训练。3DGS点云中每个点存储64维XFeat特征。(2)位姿估计：粗位姿通过2D-3D余弦相似度匹配+PnP+RANSAC获得；精化分两步：基于特征的位姿精化（渲染特征图进行关键点匹配）和基于变形的位姿精化（最小化查询图像与渲染图像间的光度变形损失）。其设计目的不是单独增加一个模块，而是用上述表示、监督或渲染约束消除原流程在目标任务中的具体失配。提出GSplatLoc，首次将XFeat密集关键点描述符蒸馏到3DGS中实现视觉定位。两阶段方法：粗位姿通过2D-3D余弦相似度匹配获得，精化阶段通过渲染光度变形损失优化。在室内外数据集上超越NeRFMatch和PNeRFLoc等最新方法。适用时还应保留论文给出的边界：特征蒸馏增加3DGS存储开销；粗位姿估计依赖充分的关键点匹配；精化阶段的渲染损失在纹理缺失区域可能不稳定。
\textbf{方法信息：} 代码：\url{https://github.com/be2rlab/gsplatloc}；GitHub 关注度较高，约 145 stars。



% METHOD: 2412.20056
\paragraph{GSplatLoc（可微深度渲染的超精确定位）（2024.12）}
GSplatLoc~\cite{2412_20056} 的局部锚点是可微渲染定位。给定预构建 3DGS 场景和查询深度图，它直接最小化查询深度与高斯渲染深度之间的差异，用渲染器反传梯度优化相机 $R,t$。相比需要训练场景坐标回归网络的方法，它不需要为每个新场景重新学习定位模型。

系统核心是 GPU 加速深度渲染、鲁棒深度损失和 6-DoF 位姿梯度下降。它专注定位而非完整 SLAM，因此可作为已有高斯地图上的 relocalization 模块。项目实现与 ultra-precise camera localization 对应。相对 SplatLoc 的地标描述符路线，GSplatLoc 依赖初始位姿附近的可微几何残差，适合高精度微调而非大范围检索。
\textbf{实现细节：} GSplatLoc（2024.12） 沿用 GSLoc 的梯度优化范式，但做了三项关键改变：将优化信号从 RGB 光度差异替换为深度图差异，将位姿参数化从 $\mathfrak{se}(3)$ 李代数改为四元数-平移表示，并借助 gsplat 库的全可微渲染管线直接反向传播梯度，无需像 GSLoc 那样手动推导解析梯度。目标函数设计体现了深度优化的特点：除了标准的深度 $L_1$ 损失 $\mathcal{L}_d$ 外，额外引入等高线损失 $\mathcal{L}_c$——用 Sobel 算子提取渲染深度图与观测深度图的梯度（边缘），约束二者的结构对齐——两项以 $\lambda_1\!=\!0.8$、$\lambda_2\!=\!0.2$ 加权组合。等高线损失专注于深度不连续处的对齐，对绝对深度尺度的偏差不敏感，在纹理缺失但几何结构清晰的区域提供关键梯度信号。场景初始化阶段，GSplatLoc 对参考深度图生成的点云施加 PCA 对齐，将点云主轴与坐标轴对齐以统一优化起点，显著提高了不同数据集间优化损失的收敛稳定性。求解采用 Adam 优化器，旋转学习率 $5\!\times\!10^{-4}$、平移学习率 $10^{-3}$，通常约 100 次迭代即收敛——远快于 GSLoc 的数百至数千步。
\textbf{设计动机与边界：} 相对原始 3DGS 以独立高斯和图像重建误差为主的设定，GSplatLoc 把问题落在以下可实现的改动上：GSplatLoc将相机定位问题转化为可微优化：(1) 给定预构建的3DGS场景表示和一组参考深度图像，将查询深度图像与3DGS渲染深度图之间的差异作为优化目标；(2) 利用3DGS的可微渲染管线计算损失相对于相机位姿(R,t)的梯度，通过梯度下降直接优化6-DoF位姿；(3) GPU加速的深度渲染管线支持实时位姿估计；(4) 设计了鲁棒的损失函数和优化策略以处理复杂相机运动。该方法专注于定位而非完整SLAM，可集成到现有3DGS SLAM框架中。其设计目的不是单独增加一个模块，而是用上述表示、监督或渲染约束消除原流程在目标任务中的具体失配。提出GSplatLoc，首个充分利用3DGS可微渲染进行超精确相机定位的方法，在Replica数据集上达到0.01cm平移误差和接近0的旋转误差，建立新的定位精度基准。适用时还应保留论文给出的边界：依赖预构建的高质量3DGS场景表示；对初始位姿估计有一定敏感性；实时性能受3DGS模型规模影响。
\textbf{方法信息：} 代码：\url{https://github.com/be2rlab/gsplatloc}；GitHub 关注度较高，约 145 stars。



% METHOD: 2504.20379
\paragraph{GSFeatLoc（渲染 RGB-D 后做特征对应）（2025.04）}
GSFeatLoc~\cite{2504_20379} 提出一个更轻量的替代思路：不一定要在高斯内部学习描述符，也不一定要直接最小化光度残差，而是从初始位姿附近用 3DGS 渲染合成 RGB-D，再在查询图像和合成图像之间做普通 2D-2D 特征匹配。深度图把合成图像特征提升到 3D，最终用 PnP 求解查询相机。

这一路线相对 iComMa 式光度优化的优势是容忍更大的初始位姿误差，推理也更快；代价是仍需要一个可用初值和稳定图像特征。项目实现围绕 feature correspondence on 3DGS。相对 SplatLoc，GSFeatLoc 不把地标永久写入高斯属性，而是把高斯当成可快速合成参考视图的地图。
\textbf{实现细节：} GSFeatLoc（2025.04） 与上述方法形成鲜明对比：LSGS-Loc、Hi$^{}2$-GSLoc 等 3DGS 定位方法都基于光度优化范式——迭代渲染图像、计算光度损失、反向传播更新位姿——通常需要数十至上百次渲染迭代，单帧推理超过 10 秒。GSFeatLoc 将这一范式替换为特征对应范式：仅从初始位姿渲染一次 RGBD 图像，用 SuperPoint 提取查询图像与渲染图像的特征点、SuperGlue 建立 2D-2D 匹配，再利用渲染深度图将渲染图像侧的 2D 匹配点反投影为 3D 坐标，将 2D-2D 对应提升为 2D-3D 对应，最终通过 PnP + RANSAC 一次性求解位姿。整个过程无需迭代渲染、无需梯度优化、无需训练场景特定的特征提取器——SuperPoint/SuperGlue 作为预训练模型直接操作。
\textbf{方法信息：} 代码：\url{https://github.com/jongwonjlee/gsfeatloc}。
\textbf{设计动机与边界：} 相对原始 3DGS 以独立高斯和图像重建误差为主的设定，GSFeatLoc 把问题落在以下可实现的改动上：GSFeatLoc采用三步管线：(1) 渲染合成图像：使用3DGS在初始位姿估计处渲染合成RGBD图像和深度图。(2) 建立2D-2D特征对应：在查询图像和合成图像之间提取并匹配特征点。(3) 建立2D-3D对应并求解PnP：利用深度图将2D特征点提升到3D空间，建立2D-3D对应关系，通过PnP问题求解最终位姿估计。方法概念上是iComMa（基于光度损失的3DGS定位）的特征对应替代方案。方法虽假设有初始位姿估计，但容忍初始误差高达55度旋转和1.1单位平移。其设计目的不是单独增加一个模块，而是用上述表示、监督或渲染约束消除原流程在目标任务中的具体失配。提出GSFeatLoc，利用3DGS渲染合成RGBD图像并通过特征对应（而非光度损失最小化）进行视觉定位，推理时间从超过10秒缩短至0.1秒（加速超100倍），同时容忍高达55度的初始位姿误差。适用时还应保留论文给出的边界：需要初始位姿估计（虽容忍较大误差）；对严重遮挡或纹理缺失区域效果下降；在Tanks and Temples等复杂场景上成功率有限。
% METHOD: 2508.18242
\paragraph{GSVisLoc（可泛化的 2D-3D 跨模态定位）（2025.08）}
GSVisLoc~\cite{2508_18242} 关注可泛化视觉定位：给定一个普通 3DGS 场景，不修改、不重训，也不额外生成参考图像，直接建立查询图像与高斯集合的匹配。它的改变是用 KPConv-based encoder 对 3D Gaussians 编码和下采样，再用图像编码器提取粗细两级 2D 特征，通过交错 self-attention 和 cross-attention 建立跨模态对应。

粗匹配得到初始 2D-3D 对应，精匹配细化局部位置，最后用 PnP 求位姿并通过 ICP-like 步骤精化。它相对 SplatLoc 和 GSplatLoc 的区别在于，不依赖场景专属特征蒸馏，而是希望模型能泛化到新高斯地图。正文中应把它看作“高斯地图作为 3D token 集合”的定位路线，而不是传统渲染残差定位。
\textbf{实现细节：} GSFeatLoc（2025.04） 与上述方法形成鲜明对比：LSGS-Loc、Hi$^{}2$-GSLoc 等 3DGS 定位方法都基于光度优化范式——迭代渲染图像、计算光度损失、反向传播更新位姿——通常需要数十至上百次渲染迭代，单帧推理超过 10 秒。GSFeatLoc 将这一范式替换为特征对应范式：仅从初始位姿渲染一次 RGBD 图像，用 SuperPoint 提取查询图像与渲染图像的特征点、SuperGlue 建立 2D-2D 匹配，再利用渲染深度图将渲染图像侧的 2D 匹配点反投影为 3D 坐标，将 2D-2D 对应提升为 2D-3D 对应，最终通过 PnP + RANSAC 一次性求解位姿。整个过程无需迭代渲染、无需梯度优化、无需训练场景特定的特征提取器——SuperPoint/SuperGlue 作为预训练模型直接操作。
\textbf{方法信息：} NVIDIA 参与了这篇论文。
\textbf{设计动机与边界：} 相对原始 3DGS 以独立高斯和图像重建误差为主的设定，GSVisLoc 把问题落在以下可实现的改动上：方法包含三个步骤：(1)3DGS编码：使用KPConv-based编码器对3D高斯进行编码和下采样（过滤低不透明度高斯后随机采样子集），输入高斯参数（不透明度、球谐系数、方向、缩放），输出代表3D空间区域的特征；(2)2D图像编码：使用图像编码器提取粗（8x8）和细（2x2）两级特征；(3)粗-精匹配：通过交错自注意力和交叉注意力层对齐3D和2D特征，粗匹配建立初始对应关系，精匹配细化对应关系；(4)位姿精化：使用PnP求解器估计初始位姿，再通过ICP-like精化提升精度。其设计目的不是单独增加一个模块，而是用上述表示、监督或渲染约束消除原流程在目标任务中的具体失配。提出GSVisLoc方法，直接利用3DGS场景表示进行视觉定位，无需修改、重训练或额外参考图像。设计三步流程（粗匹配-精匹配-位姿精化），使用KPConv编码器将3D高斯下采样编码为3D特征，图像编码器提取2D特征，通过交错自注意力和交叉注意力层建立跨模态匹配。可泛化到新场景无需额外训练。适用时还应保留论文给出的边界：对高斯数量的限制可能影响大场景精度；粗匹配阶段可能遗漏正确对应；对极端视角变化的鲁棒性有待验证。
% METHOD: 2605.04730
\paragraph{ULF-Loc（消除 alpha blending 特征偏差）（2026.05）}
ULF-Loc~\cite{2605_04730} 延续高斯特征场定位，但指出一个更细的偏差来源：3DGS 的 alpha blending 会让每个高斯特征为了补偿邻域贡献而偏离真实地标特征，最终影响 2D-3D 匹配。它的改变不是再换一个描述符，而是从理论上分析 alpha-blending feature optimization 的偏差，并用几何加权特征融合替代偏差特征优化。

具体模块包括关键点一致性地标采样、几何加权特征融合和局部几何一致性验证。前者选择多视角中稳定匹配的高斯，后者用几何关系而非 alpha 权重聚合特征，并拒绝由渲染伪影造成的误匹配。项目实现面向 unbiased landmark feature。相对 XFeat 版 GSplatLoc，ULF-Loc 的增量在于解释并修正高斯特征场自身的混合偏差。
\textbf{设计动机与边界：} 相对原始 3DGS 以独立高斯和图像重建误差为主的设定，ULF-Loc 把问题落在以下可实现的改动上：核心创新在于理论分析和偏差消除。首先数学推导证明alpha-blending使得渲染特征偏离真值(因为每个Gaussian的特征需补偿邻域贡献)。然后提出三个关键组件：(1)关键点一致性地标采样，选择在多视角中稳定匹配的可靠Gaussian作为地标；(2)几何加权特征融合，用几何信息(而非alpha权重)进行特征聚合，避免偏差；(3)局部几何一致性验证，拒绝由渲染伪影导致的误匹配。训练时间仅需STDLoc的1/10，GPU内存为1/6。其设计目的不是单独增加一个模块，而是用上述表示、监督或渲染约束消除原流程在目标任务中的具体失配。首次从理论上揭示了3DGS中alpha-blending特征优化引入偏差的机制。提出几何加权特征融合替代偏差特征优化，结合关键点一致性采样和几何一致性验证。在Cambridge Landmarks上减少17%平移误差，训练时间仅STDLoc的1/10。适用时还应保留论文给出的边界：依赖初始SfM点云质量；对退化场景(如弱纹理)的鲁棒性有限。
\textbf{方法信息：} 代码：\url{https://github.com/cyril-gyd/ulf-loc}；GitHub 关注度较高，约 48 stars。



\section{语义、对象与动态 SLAM}

语义和动态场景把 Gaussian SLAM 的地图语义从“颜色加几何”扩展到“类别、对象、可查询实例和运动状态”。这类方法的锚点不再只是 SplaTAM，而是语义 SLAM、开放词汇建图、对象级 SLAM 和 dynamic SLAM。它们共同改变的是高斯属性和优化权重：高斯不仅有颜色、尺度和不透明度，还可能携带语义特征、层次标签、物体位姿、动态概率或长期 motion memory。


% METHOD: 2409.12518
\paragraph{Hier-SLAM（层次类别树压缩大规模语义）（2024.09）}
Hier-SLAM~\cite{2409_12518} 关注语义类别数量扩展。若每个高斯用平面 softmax 表示数百个类别，存储、训练和渲染都会膨胀；Hier-SLAM 的改变是把类别组织成层次树。它用 LLM 生成并验证语义层次，例如从根节点到叶节点形成 Background、Structure、Plane、Wall 这样的路径。

训练中，层内交叉熵约束同层类别，跨层一致性约束父子节点，使高斯语义预测既有粗类别也有细类别。标准平面 softmax 可看作单层树的特例。项目实现面向 large-scale semantic SLAM。相对 GS3LAM 的多模态语义追踪，Hier-SLAM 主要解决语义表示压缩和可扩展性：当类别超过数百个时，语义高斯需要结构化标签空间。
\textbf{实现细节：} 从显式标签到层次化编码：语义 SLAM 的可扩展性瓶颈。OpenGS-SLAM 用一维离散标签替代高维语义特征，在物体级交互和存储效率上迈出关键一步。然而，这一路线面临一个新瓶颈：当场景语义类别数从数十增长到数百甚至上千时，每个高斯的语义嵌入维度必须随类别数线性增长（扁平 Softmax 编码），导致存储和计算成本急剧膨胀。以 ScanNet 的 550 个语义类别为例，扁平编码所需的参数量在现有硬件上已无法运行。Hier-SLAM（2024.09） 的核心洞察在于：语义信息天然具有层次结构——"椅子"是"家具"的子类，"家具"又是"物体"的子类。这一从"扁平"到"层次"的范式转换，使语义 SLAM 首次具备了处理 500+ 类别复杂场景的可扩展能力。
\textbf{实现补充：} 实现上还要核对输入、输出、参数、特征、初始化、训练损失、正则、渲染或推理接口，以及优化中的剪枝、分裂、采样、编码或解码步骤；这些抓手决定该方法到底改了 3DGS 管线的哪一环。
\textbf{设计动机与边界：} 相对原始 3DGS 以独立高斯和图像重建误差为主的设定，Hier-SLAM 把问题落在以下可实现的改动上：核心创新是层次化树表示G=(V,E)：(1) LLM层次树生成：使用GPT-4o-mini将语义类别迭代聚类为层次树，引入循环验证器确保聚类正确性；(2) 树参数化：将语义标签编码为从根到叶的路径，如Wall=Background>Structure>Plane>Wall；(3) 层次语义损失：包含层内优化（同层节点间交叉熵）和跨层优化（父子节点一致性约束）；(4) 增强的SLAM系统：改进跟踪和建图性能。标准平面Softmax表示可视为单层树编码的特例。其设计目的不是单独增加一个模块，而是用上述表示、监督或渲染约束消除原流程在目标任务中的具体失配。提出层次化类别高斯泼溅表示，利用LLM自动构建语义层次树，将语义信息紧凑编码到3DGS中，实现500+语义类别的大规模场景理解，渲染FPS达2000（含语义），同时存储和训练时间显著减少适用时还应保留论文给出的边界：LLM生成的层次树可能不完全符合特定任务需求；依赖LLM的质量；超大规模场景的可扩展性有待进一步验证
\textbf{方法信息：} 代码：\url{https://github.com/leeby68/hier-slam}；GitHub 关注度较高，约 169 stars。



% METHOD: 2411.08373
\paragraph{DG-SLAM（运动掩码与自适应高斯管理）（2024.11）}
DG-SLAM~\cite{2411_08373} 直接针对动态环境中高斯地图污染问题。它的改变是三部分：运动 mask generation、自适应高斯点管理和 hybrid camera tracking。系统用深度 warping 与在线实例视频分割结合，利用深度残差和分割结果生成动态掩码，过滤运动物体区域。

建图时，运动掩码决定哪些高斯可以插入、保留或删除；跟踪时，系统结合光度一致性和几何约束，并用累积不透明度判断渲染可靠性。关键帧插入还根据光流距离调整，在运动剧烈区域增加观测。相对 WildGS-SLAM 的不确定性预测，DG-SLAM 更显式地把动态区域识别为 mask，并将 mask 接入高斯生命周期管理。
\textbf{实现细节：} DG-SLAM（2024.11） 是首个基于 3DGS 的鲁棒动态视觉 SLAM，针对已有动态 SLAM 一类仅能处理小幅运动、另一类依赖语义分割存在域间隙的局限，用几何一致性 + 语义先验互补融合 + 粗-精两级跟踪替代单一先验，核心思想是深度变形一致性捕获语义无法识别的非常规移动物体——运动掩码取滑动窗口内多帧变形掩码交集缓解单次噪声，粗跟踪用 DROID-VO 在掩码过滤后位姿优化中置零动态像素置信度，精跟踪引入累积不透明度作为渲染可靠性指标双重过滤。
\textbf{实现补充：} 为避免只把它读成名称或数据点，还应把实现接口明确拆开：输入是什么，输出是高斯参数、图像、轨迹、标注还是评测指标，哪些参数、网络模块、训练损失或渲染接口真正接到原始 3DGS / Scaffold-GS / 4D-GS 管线上，以及剪枝、分裂、采样、解码、优化或评测协议在哪一步发生。若该条目是survey、benchmark 或 dataset，这里的“实现”对应论文集合、任务划分、相机/轨迹/标注组织、评价指标和基线配置，而不是虚构新的 rasterizer；若是算法条目，则对应高斯位置、尺度、旋转、不透明度、特征或 MLP 解码如何进入训练和渲染流程。
\textbf{方法信息：} 代码：\url{https://github.com/fudan-zvg/dg-slam}；引用很多，Semantic Scholar 约 52 次。



% METHOD: 2412.01217
\paragraph{RGBDS-SLAM（三维多级金字塔语义高斯）（2024.12）}
RGBDS-SLAM~\cite{2412_01217} 继承 RGB-D semantic dense SLAM，重点是细节恢复和多模态一致性。普通语义高斯 SLAM 若只在单一分辨率优化，细小结构和语义边界容易被平滑；RGBDS-SLAM 的改变是引入 3D multi-level pyramid Gaussian Splatting，从不同分辨率级别提取并优化细节。

系统在颜色、深度和语义三种输出上做紧耦合多特征重建优化，使 RGB、depth 和 semantic map 在渲染过程中互相增强。多级金字塔先保证全局一致，再逐层补充局部细节。项目实现围绕 RGB-D semantic dense SLAM。相对 GS3LAM，它的重点不是解析雅可比或遗忘控制，而是多尺度表示如何改善语义和几何边界。
\textbf{方法信息：} 代码：\url{https://github.com/zhenzhongcao/rgbds-slam}。
\textbf{设计动机与边界：} 相对原始 3DGS 以独立高斯和图像重建误差为主的设定，RGBDS-SLAM 把问题落在以下可实现的改动上：RGBDS-SLAM针对现有3DGS-based SLAM在细节恢复和重建一致性上的不足进行改进。核心是三维多级金字塔高斯Splatting方法，通过构建多级图像金字塔从不同分辨率级别提取丰富细节信息，逐级优化保证全局一致性。设计紧耦合多特征重建优化机制，使RGB、深度和语义图的重建精度在渲染优化过程中相互增强。通过合理的约束实现多特征信息的有效融合。系统支持RGB、深度和语义三模态的同时高质量重建。其设计目的不是单独增加一个模块，而是用上述表示、监督或渲染约束消除原流程在目标任务中的具体失配。提出三维多级金字塔高斯Splatting方法，通过多级图像金字塔提取细节信息，设计紧耦合多特征重建优化机制使RGB、深度和语义重建相互增强，在Replica和ScanNet上超越SOTA。适用时还应保留论文给出的边界：多级金字塔增加了计算和内存开销；实时性能可能受多模态联合优化影响；对语义标签质量有一定依赖。
% METHOD: 2503.02223
\paragraph{DQO-MAP（双二次曲面与对象级高斯）（2025.03）}
DQO-MAP~\cite{2503_02223} 从多物体 SLAM 切入。它不是给所有高斯预测语义类别，而是要估计每个对象的 6-DoF 位姿和高保真外观。方法把 dual quadrics 与 3DGS 结合：双二次曲面负责粗对象几何、位姿和尺度，高斯负责对象表面细节和渲染。

系统基于 ORB-SLAM2，使用 YOLOv10 生成 2D bounding boxes。对象关联先用 IoU 做粗匹配，再用 quadric distance 和遮挡面积重叠做精筛；高斯优化只在对象框区域增密和更新，CPU 管理对象初始化与关联，GPU 负责高斯优化。项目实现围绕 dual quadrics multi-object mapping。相对 OpenGS-Fusion 的开放词汇查询，DQO-MAP 的局部锚点是对象位姿和实例级地图管理。
\textbf{实现细节：} 其他相关工作。层次化语义与物体级交互方向的工作围绕"如何让 3DGS 表示支持结构化推理和交互"展开。沿层次化编码路线，Hier-SLAM 用层次树编码语义类别（深度 10 的二叉树仅需 20 维覆盖 1024 类）大幅压缩存储，Hier-SLAM++ 扩展为支持单目输入并用 DUSt3R 提供几何先验、MeshGPT 提取几何层次。沿物体级表示路线，DQO-MAP 利用 3DGS 高保真物体重建和二次曲面精确物体位姿估计实现快速物体提取，Go-SLAM 为每个高斯分配唯一物体标识符支持开放词汇查询和最优路径生成。沿概率语义建图路线，Bayesian Fields 提出任务驱动的概率语义建图方法用贝叶斯更新聚合多视角观测。三条路线分别从压缩、交互和不确定性三个角度补足了语义 SLAM 的物体级能力。
\textbf{方法信息：} 代码：\url{https://github.com/lihaoy-ux/dqo-map}。
\textbf{设计动机与边界：} 相对原始 3DGS 以独立高斯和图像重建误差为主的设定，DQO-MAP 把问题落在以下可实现的改动上：系统基于ORB-SLAM2构建，包含两大核心模块：(1) 物体位姿与尺度估计：使用YOLOv10生成2D边界框，采用双二次曲面作为几何基元估计物体6-DOF位姿。双二次曲面的形状与3D高斯相同，可自然分解为旋转、平移和尺度参数，实现位姿估计。(2) 物体级关联：设计从粗到精的关联策略，先用IoU过滤粗匹配结果，再用二次曲面标准距离（QD）精确筛选。针对遮挡情况，通过面积重叠比判定是否为同一物体的不同部分。(3) 高斯增量更新：仅在物体边界框区域进行高斯增密与优化，大幅提升效率。CPU负责物体管理（初始化、关联），GPU负责高斯优化，两者并行执行。其设计目的不是单独增加一个模块，而是用上述表示、监督或渲染约束消除原流程在目标任务中的具体失配。提出首个无需先验位姿的多物体3D映射系统，将双二次曲面（dual quadrics）与3DGS结合，实现物体6-DOF位姿估计与高保真重建的紧耦合。通过CPU-GPU并行架构（CPU管理物体、GPU优化高斯），实现从场景中快速提取独立物体。设计了从粗到精的3D-2D物体关联策略。适用时还应保留论文给出的边界：依赖RGB-D输入；物体关联在极端遮挡或物体高度相似时可能失败；未在大规模室外场景中验证。
% METHOD: 2504.03886
\paragraph{WildGS-SLAM（不确定性引导的动态单目 SLAM）（2025.04）}
WildGS-SLAM~\cite{2504_03886} 是动态场景单目高斯 SLAM 的重点方法。它不依赖显式语义标签或深度输入，而是用不确定性感知几何映射处理动态干扰。普通静态 3DGS SLAM 会把行人、车辆等动态区域写入地图，导致 tracking 和渲染都受污染；WildGS-SLAM 的改变是用 DINOv2 特征和浅层 MLP 在线预测逐像素 uncertainty map。

tracking 中，不确定性指导 dense BA 优先使用可靠静态区域；mapping 中，不确定性降低动态物体对渲染损失的影响，并帮助移除动态伪影。系统还加入深度不确定性损失增强预测。相对 DG-SLAM 的运动 mask 和实例视频分割，WildGS-SLAM 更偏纯几何与特征不确定性，不要求预定义动态类别。它适合作为“动态区域先降低权重再建静态高斯地图”的锚点。
\textbf{实现细节：} WildGS-SLAM（2025.04） 沿 DG-SLAM 的"检测动态-排除-重建静态"路线，但将动态检测从二值掩码替换为逐像素连续不确定性，用 DINOv2 特征驱动的浅层 MLP 在线预测不确定性图，无需语义分割标签。这一设计的动机在于：DG-SLAM 的语义分割仅能识别预定义类别，深度变形掩码在物体边缘易误判，两者取交集仍会漏检或误检。该不确定性图在两个环节同时发挥作用：在跟踪中，基于 DROID-SLAM 的光流跟踪框架引入不确定性加权的密集光束法平差（用马氏距离替代标准 $L_2$ 范数，使高不确定性像素对位姿优化的影响最小化），并额外加入 Metric3D 单目深度正则化以稳定在线训练早期不确定性尚未收敛时的位姿估计；在建图中，不确定性图作为渲染光度损失和深度损失的像素级权重，自动降低动态区域的优化贡献。值得注意的实现细节是梯度流分离——3DGS 地图优化与不确定性 MLP 训练使用各自独立的损失函数，互不回传梯度，避免两者互相干扰导致训练不稳定。系统仅需单目 RGB 输入，在多个动态环境数据集上实现无伪影视角合成。
\textbf{锚点细化：} WildGS-SLAM：沿 DG-SLAM 的"检测动态-排除-重建静态"路线，但将动态检测从\emph{二值掩码}替换为\emph{逐像素连续不确定性}，用 DINOv2 特征驱动的浅层 MLP 在线预测不确定性图，无需语义分割标签。这一设计的动机在于：DG-SLAM 的语义分割仅能识别预定义类别，深度变形掩码在物体边缘易误判，两者取交集仍会漏检或误检。WildGS-SLAM 的不确定性 MLP 随流式帧在线训练——对每帧提取 DINOv2 特征（采用注入了 3D 感知的微调版本），经浅层 MLP 解码为逐像素不确定性图，以修改版 SSIM 损失、方差正则项和深度一致性损失联合训练。该不确定性图在两个环节同时发挥作用：在跟踪中，基于 DROID-SLAM 的光流跟踪框架引入不确定性加权的密集光束法平差（用马氏距离替代标准 $L_2$ 范数，使高不确定性像素对位姿优化的影响最小化），并额外加入 Metric3D 单目深度正则化以稳定在线训练早期不确定性尚未收敛时的位姿估计；在建图中，不确定性图作为渲染光度损失和深度损失的像素级权重，自动降低动态区域的优化贡献。值得注意的实现细节是梯度流分离——3DGS 地图优化与不确定性 MLP 训练使用各自独立的损失函数，互不回传梯度，避免两者互相干扰导致训练不稳定。系统仅需单目 RGB 输入，在多个动态环境数据集上实现无伪影视角合成。
\textbf{设计动机与边界：} 相对原始 3DGS 以独立高斯和图像重建误差为主的设定，WildGS-SLAM 把问题落在以下可实现的改动上：WildGS-SLAM的核心是不确定性预测和集成机制：(1)不确定性预测：使用预训练DINOv2特征提取器获取图像特征，通过浅层MLP在线训练预测逐像素不确定性图，增量更新以适应流式输入帧。引入自定义深度不确定性损失增强训练。(2)跟踪模块：利用不确定性信息增强光流跟踪中的密集BA，优先处理可靠区域。(3)建图模块：不确定性预测指导渲染损失进行高斯地图优化，移除动态干扰物，精炼重建质量。整个系统仅使用单目RGB输入，无需显式深度或语义标签，采用纯几何方法处理动态内容。为了消除原始 3DGS 在该目标任务中的具体失配，方法把上述表示、监督或渲染约束作为可学习和可执行的改动，而不是单独堆叠模块。提出WildGS-SLAM，首个利用不确定性感知几何映射处理动态环境的单目RGB高斯SLAM系统。通过DINOv2特征和浅层MLP预测逐像素不确定性图，引导动态物体移除，增强密集BA和高斯地图优化。在多个数据集上实现无伪影视角合成，在动态环境中超越SOTA方法。适用时还应保留论文给出的边界：依赖DINOv2特征质量；对极快运动物体的处理可能不足；不确定性MLP需要在线训练时间。
\textbf{方法信息：} Microsoft 参与了这篇论文；引用很多，Semantic Scholar 约 69 次。



% METHOD: 2508.01150
\paragraph{OpenGS-Fusion（开放词汇对象级融合）（2025.08）}
OpenGS-Fusion~\cite{2508_01150} 把 3DGS 与 TSDF 结合，用于开放词汇 dense mapping 和对象级理解。普通语义 SLAM 类别集合固定，难以响应任意文本查询；OpenGS-Fusion 的改变是把在线语义特征无损融合进混合 3DGS-TSDF 表示，并用 MLLM-assisted adaptive threshold 根据查询语义和场景上下文调整相似度阈值。

框架中，TSDF 提供几何稳定性，高斯提供可渲染外观和语义特征承载。具体流程是在线输入 RGB-D 或多视图观测，先更新 TSDF 几何，再把开放词汇语义特征融合到 3DGS/TSDF 混合地图中；查询阶段，MAAT 模块用多模态大模型根据文本语义和场景上下文调整特征相似度阈值，替代固定阈值，改善物体边界分割和文本查询响应。相对 Hier-SLAM 的层次固定类别，OpenGS-Fusion 面向开放词汇交互：读者应把它理解为机器人在线感知和语言查询场景中的高斯语义融合路线。
\textbf{方法信息：} 代码：\url{https://github.com/YOUNG-bit/OpenGS-Fusion}。
\textbf{设计动机与边界：} 相对原始 3DGS 以独立高斯和图像重建误差为主的设定，OpenGS-Fusion 把问题落在以下可实现的改动上：OpenGS-Fusion框架包含三个核心设计：(1)3DGS+TSDF混合表示：将3D高斯表征与截断符号距离场结合，实现在线语义特征的无损融合，兼顾渲染质量和几何精度；(2)MLLM辅助自适应阈值（MAAT）：利用多模态大语言模型（如GPT-4V）根据查询语义和场景上下文自适应调整特征相似度阈值，替代固定阈值策略，显著提升物体边界分割精度；(3)在线管线：支持实时感知，无需离线预处理或预训练，适合机器人探索和交互场景。框架能响应文本查询提供精确的3D物体模型。其设计目的不是单独增加一个模块，而是用上述表示、监督或渲染约束消除原流程在目标任务中的具体失配。提出OpenGS-Fusion框架，将3DGS与TSDF结合实现在线语义特征无损融合。引入MLLM辅助自适应阈值（MAAT）方法，通过多模态大语言模型自适应调整相似度阈值，3D mIoU提升17%。支持在线感知和语言引导场景交互。适用时还应保留论文给出的边界：依赖MLLM进行阈值调整，引入额外延迟；TSDF分辨率限制了细节表达；大规模场景的内存开销可能较大。
% METHOD: 2603.27781
\paragraph{GS3LAM（语义高斯场的完整 SLAM 锚点）（2026.03）}
GS3LAM~\cite{2603_27781} 是语义 Gaussian SLAM 的代表性锚点。它把场景建模为 Semantic Gaussian Field，其中每个高斯包含位置、旋转、尺度、不透明度、RGB 颜色和语义特征。普通 RGB-D 高斯 SLAM 只能渲染颜色和深度，不能直接输出一致语义地图；GS3LAM 的改变是把颜色、深度和语义一起放入 splatting 渲染和位姿优化。

tracking 阶段，方法推导 SG-Field 关于相机位姿的解析雅可比，用 RGB、depth 和 semantic losses 联合优化当前位姿；语义输出通过 splatting 渲染后再解码，使同一个高斯集合同时服务外观、几何和类别对齐。mapping 阶段，深度自适应尺度正则解决高斯与几何表面对齐问题，随机采样关键帧映射缓解模型 SLAM 的灾难性遗忘，自适应高斯扩展逐步补全语义地图。论文还专门讨论 model-based SLAM 的优化偏置：如果只反复优化局部共视关键帧，语义特征会加剧遗忘。相对 RGBDS-SLAM 的多级金字塔，GS3LAM 的锚点意义在于完整定义了语义高斯场如何参与 tracking、mapping 和遗忘控制。
\textbf{锚点细化：} GS3LAM：针对 PanoSLAM 和 GSFF-SLAM 在增量建图场景下高斯尺度失配和灾难性遗忘的局限，用\textbf{深度自适应尺度正则化 + 随机采样关键帧建图}替代 LCKM 局部共视性策略。核心思想是\textbf{DSR 用 $2\sigma$ 规则约束全局尺度避免边缘模糊，RSKM 均匀采样全部关键帧消除优化偏置}——DSR 使跟踪误差降低 16
\textbf{实现细节：}
\textbf{实现补充：} 实现上还要核对输入、输出、参数、特征、初始化、训练损失、正则、渲染或推理接口，以及优化中的剪枝、分裂、采样、编码或解码步骤；这些抓手决定该方法到底改了 3DGS 管线的哪一环。
\textbf{锚点细化：} GS3LAM：针对 PanoSLAM 和 GSFF-SLAM 在增量建图场景下高斯尺度失配和灾难性遗忘的局限，用\textbf{深度自适应尺度正则化 + 随机采样关键帧建图}替代 LCKM 局部共视性策略，核心思想是\textbf{DSR 用 $2\sigma$ 规则约束全局尺度避免边缘模糊，RSKM 均匀采样全部关键帧消除优化偏置}——DSR 使跟踪误差降低 16%、PSNR 提升 1.17 dB、语义 mIoU 提升 3.36\%，RSKM 使 PSNR 提升 5.49 dB、方差降低 24.42 倍。
\textbf{方法信息：} 代码：\url{https://github.com/lif314/gs3lam}；GitHub 关注度较高，约 98 stars；引用很多，Semantic Scholar 约 38 次。



% METHOD: 2604.12837
\paragraph{GGD-SLAM（可泛化运动模型与动态特征增强）（2026.04）}
GGD-SLAM~\cite{2604_12837} 也处理动态单目高斯 SLAM，但它强调动态性需要时间上下文。单帧图像很难判断一个人是静止背景还是暂时停下的动态物体；GGD-SLAM 的改变是用 FIFO queue 和 sequence attention 从输入帧序列中提取动态语义特征，再通过 dynamic feature enhancer 分离静态和动态组件。

对被动态物体遮挡的区域，系统用 KD-tree 构建静态高斯点云并进行静态信息采样填充，减少动态干扰写入地图。训练损失中加入 interferer-adaptive SSIM，使系统对动态遮挡更稳。相对 DG-SLAM 的显式 mask，GGD-SLAM 更依赖可泛化运动模型和序列注意力；相对 WildGS-SLAM，它把不确定性扩展为时序动态特征。
\textbf{实现细节：} DyPho-SLAM（2025.09） 集成先验图像信息生成精化掩码，自适应特征提取策略在移除动态障碍后增强优化约束。PG-SLAM（2024.11） 同时建图动态前景（包括非刚体人体和刚体物品）和静态背景，利用静态背景和动态前景增加观测进行噪声补偿。LVD-GS（2025.10） 融合开放世界分割和隐式残差约束生成细粒度动态掩码，由 DINO-Depth 特征的不确定性估计引导。DynaGSLAM（2025.03） 是首个在动态场景中同时实现在线 GS 渲染、跟踪和运动预测的实时系统——而非简单地检测和移除动态区域。
\textbf{实现补充：} 实现上还要核对输入、输出、参数、特征、初始化、训练损失、正则、渲染或推理接口，以及优化中的剪枝、分裂、采样、编码或解码步骤；这些抓手决定该方法到底改了 3DGS 管线的哪一环。
\textbf{方法信息：} Huawei 参与了这篇论文。
\textbf{设计动机与边界：} 相对原始 3DGS 以独立高斯和图像重建误差为主的设定，GGD-SLAM 把问题落在以下可实现的改动上：GGD-SLAM包含以下核心设计：（1）FIFO队列结构：管理渐进式输入帧，通过序列注意力机制从时间序列中提取动态语义特征（单帧缺乏运动信息，需要多帧时序上下文）；（2）动态特征增强器：结合序列注意力提取的动态语义特征，分离静态和动态组件；（3）静态信息采样填充遮挡区域：通过KD树构建静态高斯点云，最小化动态干扰物对静态组件的影响；（4）干扰物自适应SSIM损失：针对动态环境设计的损失函数，增强系统对动态干扰的韧性。系统无需深度输入或预定义语义标签。其设计目的不是单独增加一个模块，而是用上述表示、监督或渲染约束消除原流程在目标任务中的具体失配。提出GGD-SLAM，利用可泛化运动模型解决动态环境下的单目3DGS SLAM问题，无需预定义语义标注或深度输入。通过FIFO队列和序列注意力机制提取动态语义特征，结合动态特征增强器分离动静组件，设计干扰物自适应SSIM损失，在真实动态数据集上实现SOTA位姿估计和密集重建。适用时还应保留论文给出的边界：依赖足够的帧间运动信息进行动态识别；FIFO队列大小影响性能和计算开销；极端快速运动可能导致动态识别失败。
% METHOD: 2606.29237
\paragraph{MoPe（运动持久性的地图级记忆）（2026.06）}
MoPe~\cite{2606_29237} 的增量很具体：它认为动态性不是瞬时视觉属性，而应作为地图级 memory 维护。普通 dynamic SLAM 如果只看当前帧运动，行人暂时停下、遮挡后重现或慢速移动时很容易被吸收到静态高斯地图。MoPe 的改变是在 monocular Gaussian Splatting SLAM 中记录并传播区域或高斯的 motion permanence。

当某个高斯或区域多次表现出动态行为，系统会把这种信息写入地图状态，并在后续帧继续传播；即使行人短暂停下、被遮挡后重现，或慢速物体一度看似静止，地图更新模块也不会立刻把它当作可靠静态结构融合。输入仍是单目 SLAM 帧序列，训练和建图时多维护一个 motion permanence 变量，输出是更少动态污染的静态高斯地图。相对 DG-SLAM 和 GGD-SLAM，MoPe 的篇幅可以较短，但它提供了一个有用概念：动态判断应从逐帧 mask 上升为地图状态变量。
\textbf{实现细节：} MoPe（2026.06） 延续 3DGS-SLAM 在"Dynamic-scene robustness in Gaussian SLAM"上的基本设定，主要改变是：Motion Permanence：将动态性建模为跨时间的持久属性而非单帧观测。3DGS 的显式高斯表示和可微渲染能力保持为基础。其设计动机是：指出动态性不是瞬时视觉属性，而应作为地图级记忆变量维护，从表示层避免动态内容被吸收到静态地图。具体而言，在 monocular Gaussian Splatting SLAM 中记录和传播区域或高斯的 motion permanence，当行人暂停、遮挡后重现等情况发生时仍保持动态判断，减少错误建图。适用边界是：依赖动态性估计和时间关联质量，对长期遮挡、缓慢移动或语义不确定对象仍可能误判。
\textbf{实现补充：} 实现上还要核对输入、输出、参数、特征、初始化、训练损失、正则、渲染或推理接口，以及优化中的剪枝、分裂、采样、编码或解码步骤；这些抓手决定该方法到底改了 3DGS 管线的哪一环。
\textbf{技术实现补充：} MoPe。MoPe（2026.06） 延续 3DGS-SLAM 在"Dynamic-scene robustness in Gaussian SLAM"上的基本设定，主要改变是：Motion Permanence：将动态性建模为跨时间的持久属性而非单帧观测。3DGS 的显式高斯表示和可微渲染能力保持为基础。其设计动机是：指出动态性不是瞬时视觉属性，而应作为地图级记忆变量维护，从表示层避免动态内容被吸收到静态地图。具体而言，在 monocular Gaussian Splatting SLAM 中记录和传播区域或高斯的 motion permanence，当行人暂停、遮挡后重现等情况发生时仍保持动态判断，减少错误建图。适用边界是：依赖动态性估计和时间关联质量，对长期遮挡、缓慢移动或语义不确定对象仍可能误判。
\textbf{方法信息：} 代码：\url{https://github.com/chloeqxq/mope}。
\textbf{设计动机与边界：} 相对原始 3DGS 以独立高斯和图像重建误差为主的设定，MoPe 把问题落在以下可实现的改动上：在 monocular Gaussian Splatting SLAM 中记录和传播区域或高斯的 motion permanence，当行人暂停、遮挡后重现等情况发生时仍保持动态判断，减少错误建图。其设计目的不是单独增加一个模块，而是用上述表示、监督或渲染约束消除原流程在目标任务中的具体失配。指出动态性不是瞬时视觉属性，而应作为地图级记忆变量维护，从表示层避免动态内容被吸收到静态地图。适用时还应保留论文给出的边界：依赖动态性估计和时间关联质量，对长期遮挡、缓慢移动或语义不确定对象仍可能误判。

\section{大规模、协作与相邻任务}

这一节收纳不完全落在前面路线中的高斯地图系统。它们共同关心地图如何在大规模、资源受限、协作或非典型任务中被维护：RTG-SLAM 和 Pocket-SLAM 控制地图冗余，DiskChunGS 把大场景放到磁盘分块，MAC-Ego3D 处理多智能体共识，4D Gaussian Splatting SLAM 把动态物体显式建成四维高斯，Mesh-Learner 与 GS-STVSR 则说明某些 SLAM 相邻任务会借用高斯思想但目标并非标准定位建图。
% METHOD: 2311.11700
\paragraph{GS-SLAM（SLAM 与在线建图）（2023.11）}
GS-SLAM~\cite{2311_11700} 关注的问题是：把 3DGS 用于 SLAM、定位或稠密建图 GS-SLAM 是宽召回补全中发现的论文，主要根据标题、摘要和 arXiv 类别判定其与 3DGS 的关系。
相对原始 3DGS，GS-SLAM 的变化不是只替换损失函数，而是把表示、初始化、约束或网络预测路径重新组织到任务流程中：该工作被召回是因为论文文本包含 Gaussian/3D Gaussian 相关表述；当前分类先确定其任务域和与 3DGS 的相关性，后续证据流水线再补代码、企业和引用信息。
\textbf{实现细节：} 具体实现可以按输入、核心模块和渲染/优化路径来读。输入端首先取得任务所需的图像、相机、文本、时间序列、语义或几何线索，并把这些线索转换成可以约束高斯中心、尺度、旋转、不透明度、颜色或附加特征的条件。核心结构中，方法使用的主要机制是：从摘要层面看，GS-SLAM 围绕“把 3DGS 用于 SLAM、定位或稠密建图”展开。若其正文或代码后续可读，应进一步核对它具体改动的是表示、初始化、渲染、优化、损失、动态建模、压缩或应用接口中的哪一环。 训练或推理时，这些模块通常先给出初始化或中间特征，再通过可微 Gaussian rasterization、重建损失、感知/语义约束、几何正则或任务损失回传到高斯参数；如果方法包含前馈网络，则网络输出的是高斯属性、变形场、压缩码、语义特征或控制变量，而不是只输出最终 2D 图像。渲染路径仍以显式高斯投影、深度排序和 alpha 合成为基础，只是在采样、属性预测、局部编辑、动态变形、压缩编码或下游任务接口处加入额外约束。这样写入本章时，它与前后方法的关系是：相对原始 3DGS 和本章已有锚点方法，它可以理解为 3DGS extends\_or\_uses 宽召回分类阶段的关系判断，后续需要按论文方法部分复核。
\textbf{证据与边界：} 本轮召回分类尚未逐项核对实验表格；后续证据补全阶段需要读取论文源码/PDF或项目页确认数据集、指标和对比方法。 该条来自宽召回批次；源码下载受 arXiv e-print 网络限制，部分论文暂时只能依据标题、摘要和类别分类，需要在入可信书前做人工/LLM精读复核。 本书只采用本地 paper pack、数据库分类、源码/引用登记和论文摘要中可核验的信息；没有本地证据支撑的速度、排行榜或开源完整性不在正文中强化。\textbf{方法信息：} 代码：\url{https://github.com/zerokhong1/ros2-gs-slam-demo}；引用信号：485 citations (15.3/month, domain-relative cutoff 3.4/month, semantic\_scholar)。
\textbf{原始方法补充：} GS-SLAM 将原始 3DGS 从离线 SfM 场景改成 RGB-D 在线系统：输入是带已知内参的连续 RGB-D 流，输出是每帧相机位姿和可持续更新的稠密高斯地图。映射端在关键帧位姿确定后以可微 splatting 渲染 RGB 与深度，并采用自适应扩张策略加入解释新几何的高斯、删除噪声高斯；跟踪端固定地图，依据 RGB-D 渲染损失的解析反向优化位姿，并以粗到细调度扩大收敛域。为了让新 RGB-D 观测在不重训整段序列的条件下同时校正位姿和地图，GS-SLAM 把“先定位再建图”的 SLAM 约束直接写进高斯的训练和渲染循环；静态 3DGS 的全局联合训练并不具备这一在线接口。
% METHOD: 2312.10070
\paragraph{Gaussian-SLAM（SLAM 与在线建图）（2023.12）}
Gaussian-SLAM~\cite{2312_10070} 关注的问题是：把 3DGS 用于 SLAM、定位或稠密建图 Gaussian-SLAM 是宽召回补全中发现的论文，主要根据标题、摘要和 arXiv 类别判定其与 3DGS 的关系。
相对原始 3DGS，Gaussian-SLAM 的变化不是只替换损失函数，而是把表示、初始化、约束或网络预测路径重新组织到任务流程中：该工作被召回是因为论文文本包含 Gaussian/3D Gaussian 相关表述；当前分类先确定其任务域和与 3DGS 的相关性，后续证据流水线再补代码、企业和引用信息。
\textbf{实现细节：} 具体实现可以按输入、核心模块和渲染/优化路径来读。输入端首先取得任务所需的图像、相机、文本、时间序列、语义或几何线索，并把这些线索转换成可以约束高斯中心、尺度、旋转、不透明度、颜色或附加特征的条件。核心结构中，方法使用的主要机制是：从摘要层面看，Gaussian-SLAM 围绕“把 3DGS 用于 SLAM、定位或稠密建图”展开。若其正文或代码后续可读，应进一步核对它具体改动的是表示、初始化、渲染、优化、损失、动态建模、压缩或应用接口中的哪一环。 训练或推理时，这些模块通常先给出初始化或中间特征，再通过可微 Gaussian rasterization、重建损失、感知/语义约束、几何正则或任务损失回传到高斯参数；如果方法包含前馈网络，则网络输出的是高斯属性、变形场、压缩码、语义特征或控制变量，而不是只输出最终 2D 图像。渲染路径仍以显式高斯投影、深度排序和 alpha 合成为基础，只是在采样、属性预测、局部编辑、动态变形、压缩编码或下游任务接口处加入额外约束。这样写入本章时，它与前后方法的关系是：相对原始 3DGS 和本章已有锚点方法，它可以理解为 3DGS extends\_or\_uses 宽召回分类阶段的关系判断，后续需要按论文方法部分复核。
\textbf{证据与边界：} 本轮召回分类尚未逐项核对实验表格；后续证据补全阶段需要读取论文源码/PDF或项目页确认数据集、指标和对比方法。 该条来自宽召回批次；源码下载受 arXiv e-print 网络限制，部分论文暂时只能依据标题、摘要和类别分类，需要在入可信书前做人工/LLM精读复核。 本书只采用本地 paper pack、数据库分类、源码/引用登记和论文摘要中可核验的信息；没有本地证据支撑的速度、排行榜或开源完整性不在正文中强化。\textbf{方法信息：} 代码：\url{https://github.com/VladimirYugay/Gaussian-SLAM}；引用信号：243 citations (7.8/month, domain-relative cutoff 3.4/month, semantic\_scholar)；代码关注：vladimiryugay/gaussian-slam has 1170 stars (37.6/month, domain-relative cutoff 20.9/month)。
\textbf{设计动机与边界：} 相对原始 3DGS 以独立高斯和图像重建误差为主的设定，Gaussian-SLAM 把问题落在以下可实现的改动上：从摘要层面看，Gaussian-SLAM 围绕“把 3DGS 用于 SLAM、定位或稠密建图”展开。若其正文或代码后续可读，应进一步核对它具体改动的是表示、初始化、渲染、优化、损失、动态建模、压缩或应用接口中的哪一环。其设计目的不是单独增加一个模块，而是用上述表示、监督或渲染约束消除原流程在目标任务中的具体失配。该工作被召回是因为论文文本包含 Gaussian/3D Gaussian 相关表述；当前分类先确定其任务域和与 3DGS 的相关性，后续证据流水线再补代码、企业和引用信息。适用时还应保留论文给出的边界：该条来自宽召回批次；源码下载受 arXiv e-print 网络限制，部分论文暂时只能依据标题、摘要和类别分类，需要在入可信书前做人工/LLM精读复核。
% METHOD: 2402.13255
\paragraph{How NeRFs and 3D Gaussian Splatting are Reshaping SLAM（SLAM 与在线建图）（2024.02）}
How NeRFs and 3D Gaussian Splatting are Reshaping SLAM~\cite{2402_13255} 关注的问题是：把 3DGS 用于 SLAM、定位或稠密建图 How NeRFs and 3D Gaussian Splatting are Reshaping SLAM 是宽召回补全中发现的论文，主要根据标题、摘要和 arXiv 类别判定其与 3DGS 的关系。
相对原始 3DGS，How NeRFs and 3D Gaussian Splatting are Reshaping SLAM 的变化不是只替换损失函数，而是把表示、初始化、约束或网络预测路径重新组织到任务流程中：该工作被召回是因为论文文本包含 Gaussian/3D Gaussian 相关表述；当前分类先确定其任务域和与 3DGS 的相关性，后续证据流水线再补代码、企业和引用信息。
\textbf{实现细节：} 具体实现可以按输入、核心模块和渲染/优化路径来读。输入端首先取得任务所需的图像、相机、文本、时间序列、语义或几何线索，并把这些线索转换成可以约束高斯中心、尺度、旋转、不透明度、颜色或附加特征的条件。核心结构中，方法使用的主要机制是：从摘要层面看，How NeRFs and 3D Gaussian Splatting are Reshaping SLAM 围绕“把 3DGS 用于 SLAM、定位或稠密建图”展开。若其正文或代码后续可读，应进一步核对它具体改动的是表示、初始化、渲染、优化、损失、动态建模、压缩或应用接口中的哪一环。 训练或推理时，这些模块通常先给出初始化或中间特征，再通过可微 Gaussian rasterization、重建损失、感知/语义约束、几何正则或任务损失回传到高斯参数；如果方法包含前馈网络，则网络输出的是高斯属性、变形场、压缩码、语义特征或控制变量，而不是只输出最终 2D 图像。渲染路径仍以显式高斯投影、深度排序和 alpha 合成为基础，只是在采样、属性预测、局部编辑、动态变形、压缩编码或下游任务接口处加入额外约束。这样写入本章时，它与前后方法的关系是：相对原始 3DGS 和本章已有锚点方法，它可以理解为 3DGS extends\_or\_uses 宽召回分类阶段的关系判断，后续需要按论文方法部分复核。
\textbf{证据与边界：} 本轮召回分类尚未逐项核对实验表格；后续证据补全阶段需要读取论文源码/PDF或项目页确认数据集、指标和对比方法。 该条来自宽召回批次；源码下载受 arXiv e-print 网络限制，部分论文暂时只能依据标题、摘要和类别分类，需要在入可信书前做人工/LLM精读复核。 本书只采用本地 paper pack、数据库分类、源码/引用登记和论文摘要中可核验的信息；没有本地证据支撑的速度、排行榜或开源完整性不在正文中强化。\textbf{方法信息：} 代码：\url{https://github.com/spla-tam/splatam}；引用信号：127 citations (4.4/month, domain-relative cutoff 3.4/month, semantic\_scholar)；代码关注：doongli/awesome-implicit-nerf-slam has 2064 stars (72.0/month, domain-relative cutoff 20.9/month)。
\textbf{设计动机与边界：} 相对原始 3DGS 以独立高斯和图像重建误差为主的设定，How NeRFs and 3D Gaussian Splatting are Reshaping SLAM 把问题落在以下可实现的改动上：从摘要层面看，How NeRFs and 3D Gaussian Splatting are Reshaping SLAM 围绕“把 3DGS 用于 SLAM、定位或稠密建图”展开。若其正文或代码后续可读，应进一步核对它具体改动的是表示、初始化、渲染、优化、损失、动态建模、压缩或应用接口中的哪一环。其设计目的不是单独增加一个模块，而是用上述表示、监督或渲染约束消除原流程在目标任务中的具体失配。该工作被召回是因为论文文本包含 Gaussian/3D Gaussian 相关表述；当前分类先确定其任务域和与 3DGS 的相关性，后续证据流水线再补代码、企业和引用信息。适用时还应保留论文给出的边界：该条来自宽召回批次；源码下载受 arXiv e-print 网络限制，部分论文暂时只能依据标题、摘要和类别分类，需要在入可信书前做人工/LLM精读复核。
% METHOD: 2403.12550
\paragraph{RGBD GS-ICP SLAM（SLAM 与在线建图）（2024.03）}
RGBD GS-ICP SLAM~\cite{2403_12550} 关注的问题是：把 3DGS 用于 SLAM、定位或稠密建图 RGBD GS-ICP SLAM 是宽召回补全中发现的论文，主要根据标题、摘要和 arXiv 类别判定其与 3DGS 的关系。
相对原始 3DGS，RGBD GS-ICP SLAM 的变化不是只替换损失函数，而是把表示、初始化、约束或网络预测路径重新组织到任务流程中：该工作被召回是因为论文文本包含 Gaussian/3D Gaussian 相关表述；当前分类先确定其任务域和与 3DGS 的相关性，后续证据流水线再补代码、企业和引用信息。
\textbf{实现细节：} 具体实现可以按输入、核心模块和渲染/优化路径来读。输入端首先取得任务所需的图像、相机、文本、时间序列、语义或几何线索，并把这些线索转换成可以约束高斯中心、尺度、旋转、不透明度、颜色或附加特征的条件。核心结构中，方法使用的主要机制是：从摘要层面看，RGBD GS-ICP SLAM 围绕“把 3DGS 用于 SLAM、定位或稠密建图”展开。若其正文或代码后续可读，应进一步核对它具体改动的是表示、初始化、渲染、优化、损失、动态建模、压缩或应用接口中的哪一环。 训练或推理时，这些模块通常先给出初始化或中间特征，再通过可微 Gaussian rasterization、重建损失、感知/语义约束、几何正则或任务损失回传到高斯参数；如果方法包含前馈网络，则网络输出的是高斯属性、变形场、压缩码、语义特征或控制变量，而不是只输出最终 2D 图像。渲染路径仍以显式高斯投影、深度排序和 alpha 合成为基础，只是在采样、属性预测、局部编辑、动态变形、压缩编码或下游任务接口处加入额外约束。这样写入本章时，它与前后方法的关系是：相对原始 3DGS 和本章已有锚点方法，它可以理解为 3DGS extends\_or\_uses 宽召回分类阶段的关系判断，后续需要按论文方法部分复核。
\textbf{证据与边界：} 本轮召回分类尚未逐项核对实验表格；后续证据补全阶段需要读取论文源码/PDF或项目页确认数据集、指标和对比方法。 该条来自宽召回批次；源码下载受 arXiv e-print 网络限制，部分论文暂时只能依据标题、摘要和类别分类，需要在入可信书前做人工/LLM精读复核。 本书只采用本地 paper pack、数据库分类、源码/引用登记和论文摘要中可核验的信息；没有本地证据支撑的速度、排行榜或开源完整性不在正文中强化。\textbf{方法信息：} 代码：\url{https://github.com/Lab-of-AI-and-Robotics/GS\_ICP\_SLAM}；引用信号：110 citations (4.0/month, domain-relative cutoff 3.4/month, semantic\_scholar)；代码关注：lab-of-ai-and-robotics/gs\_icp\_slam has 526 stars (19.0/month, domain-relative cutoff 20.9/month)。
\textbf{设计动机与边界：} 相对原始 3DGS 以独立高斯和图像重建误差为主的设定，RGBD GS-ICP SLAM 把问题落在以下可实现的改动上：从摘要层面看，RGBD GS-ICP SLAM 围绕“把 3DGS 用于 SLAM、定位或稠密建图”展开。若其正文或代码后续可读，应进一步核对它具体改动的是表示、初始化、渲染、优化、损失、动态建模、压缩或应用接口中的哪一环。其设计目的不是单独增加一个模块，而是用上述表示、监督或渲染约束消除原流程在目标任务中的具体失配。该工作被召回是因为论文文本包含 Gaussian/3D Gaussian 相关表述；当前分类先确定其任务域和与 3DGS 的相关性，后续证据流水线再补代码、企业和引用信息。适用时还应保留论文给出的边界：该条来自宽召回批次；源码下载受 arXiv e-print 网络限制，部分论文暂时只能依据标题、摘要和类别分类，需要在入可信书前做人工/LLM精读复核。
% METHOD: 2404.19706
\paragraph{RTG-SLAM（不透明与透明高斯的紧凑地图）（2024.04）}
RTG-SLAM~\cite{2404_19706} 面向大规模环境实时重建。它继承 RGB-D 或 dense SLAM 的在线建图需求，但认为普通 3DGS 在实时系统中高斯数量和优化范围过大。方法的改变是紧凑高斯表示：添加高斯时强制其接近不透明或近透明，不透明高斯拟合表面和主颜色，透明高斯补充残余颜色。

在线优化中，系统把像素分成新观测、大颜色误差和大深度误差三类，只在这些区域显式添加高斯；又把高斯分为 stable 和 unstable，只优化不稳定高斯，并只渲染相关像素。相对 RTGS 的硬件级冗余削减，RTG-SLAM 的局部锚点是地图表示和在线优化策略：通过高斯角色分化和稳定性分组减少内存与计算。
\textbf{实现细节：} RTG-SLAM（2024.04） 强制每个高斯要么不透明要么近乎透明——不透明高斯拟合表面和主色，透明高斯拟合残差颜色。这一二值化策略的核心优势在于：通过不同的深度渲染方式使单个不透明高斯即可良好拟合局部表面区域，大幅减少所需高斯数量（避免了多个半透明高斯堆叠拟合一个表面的冗余），配合高效在线优化方案实现大规模实时重建。
\textbf{实现补充：} 实现上还要核对输入、输出、参数、特征、初始化、训练损失、正则、渲染或推理接口，以及优化中的剪枝、分裂、采样、编码或解码步骤；这些抓手决定该方法到底改了 3DGS 管线的哪一环。
\textbf{方法信息：} Baidu 参与了这篇论文；代码仓库：\url{https://github.com/misety/rtg-slam}，当前未确认完整实现。
\textbf{设计动机与边界：} 相对原始 3DGS 以独立高斯和图像重建误差为主的设定，RTG-SLAM 把问题落在以下可实现的改动上：RTG-SLAM包含两个核心设计。紧凑高斯表示：强制每个高斯在添加时确定为不透明（0.99）或近乎透明（0.1），不透明高斯拟合3D表面和主导颜色，透明高斯拟合残余颜色。深度渲染采用与颜色渲染不同的方式，使单个不透明高斯能良好拟合局部表面区域而无需多个重叠高斯，大幅降低内存和计算成本。高效在线优化：对每帧的三类像素（新观测、大颜色误差、大深度误差）显式添加高斯；将高斯分为稳定（已充分优化）和不稳定（仍在优化），仅优化不稳定高斯且仅渲染不稳定高斯占据的像素，大幅减少优化和渲染的计算量。其设计目的不是单独增加一个模块，而是用上述表示、监督或渲染约束消除原流程在目标任务中的具体失配。提出紧凑高斯表示，强制每个高斯为不透明或近乎透明，不透明高斯拟合表面，透明高斯拟合残余颜色；设计高效在线优化方案，将高斯分为稳定/不稳定，仅优化不稳定高斯；以约2倍速度和一半内存实现与SOTA NeRF SLAM相当的高质量重建。适用时还应保留论文给出的边界：二值不透明度假设可能限制半透明物体建模；稳定/不稳定分类阈值需手动设定；大场景的长期累积误差。
% METHOD: 2412.09723
\paragraph{MAC-Ego3D（多智能体高斯共识）（2024.12）}
MAC-Ego3D~\cite{2412_09723} 把高斯地图用于多智能体协作重建。单机器人 3DGS SLAM 只需维护一条轨迹和一个地图，多智能体系统则要对齐多个局部高斯场。相对单主体高斯 SLAM 基线，方法的改变是 Gaussian Consensus：每个智能体先通过 Intra-Agent Gaussian Consensus 构建局部地图，再通过 Inter-Agent Gaussian Consensus 进行全局对齐。

跨智能体阶段包括 selection、association、alignment 和 optimization。系统输入各智能体的 RGB-D 观测和局部高斯地图，先识别重叠区域相关高斯，再根据空间和外观特征建立对应，随后最小化重叠高斯的空间差异并联合优化位姿和高斯属性。项目实现使用异步多进程。相对 LoopSplat 的单机器人回环，MAC-Ego3D 把回环思想扩展为多主体之间的高斯共识。
\textbf{实现细节：} MAC-Ego3D（2024.12） 提出多智能体高斯共识（Multi-Agent Gaussian Consensus）机制，将协作建图问题建模为联合优化：各智能体最小化局部目标 $f^{a_i}(T_{1:t}^{a_i}, M^{a_i})$（自身建图与位姿估计误差），同时通过全局目标 $g^{a_i,a_j}$ 对齐重叠区域的局部地图以形成一致的全局地图。与 DiskChunGS 将全局地图管理集中在单实体不同，MAC-Ego3D 将共识分解为智能体内（Intra-Agent）和智能体间（Inter-Agent）两个层次，统一遵循四步循环：选择（Selection）$\to$ 关联（Association）$\to$ 对齐（Alignment）$\to$ 优化（Optimization），两层通过异步多进程并行执行。智能体内高斯共识实现单智能体的增量建图与跟踪。\textit{选择}步从局部地图 $M^{a_i}$ 中提取当前视野内且位于估计位姿 $T_k^{a_i}$ 邻近范围内的候选高斯，利用 SLAM 的时序连续性将搜索空间聚焦于最可能匹配当前观测的地图区域。
\textbf{实现补充：} 实现上还要核对输入、输出、参数、特征、初始化、训练损失、正则、渲染或推理接口，以及优化中的剪枝、分裂、采样、编码或解码步骤；这些抓手决定该方法到底改了 3DGS 管线的哪一环。
\textbf{方法信息：} 代码：\url{https://github.com/xiaohao-xu/mac-ego3d}。
\textbf{设计动机与边界：} 相对原始 3DGS 以独立高斯和图像重建误差为主的设定，MAC-Ego3D 把问题落在以下可实现的改动上：MAC-Ego3D使用3D高斯基元表示环境。每个智能体独立捕获RGB-D观测，通过Intra-Agent Gaussian Consensus构建局部高斯地图。全局对齐通过Inter-Agent Gaussian Consensus实现，包含四个步骤：(1) Selection：识别重叠区域的相关高斯基元；(2) Association：基于空间和外观相似性建立对应关系；(3) Alignment：最小化重叠高斯的空间差异；(4) Optimization：微调对齐的高斯和位姿。通过异步多进程实现高效并行化。统一的高斯splat表示支持高效RGB-D渲染，加速智能体间高斯关联和对齐。将局部精度与全局一致性桥接。其设计目的不是单独增加一个模块，而是用上述表示、监督或渲染约束消除原流程在目标任务中的具体失配。提出MAC-Ego3D，首个基于高斯共识的多智能体实时照片级3D重建框架。包含Intra-Agent和Inter-Agent Gaussian Consensus两阶段，实现15倍推理加速，部分场景自我运动估计误差降低一个数量级，RGB PSNR提升4-10 dB。建立新SOTA。适用时还应保留论文给出的边界：依赖RGB-D输入；通信开销随智能体数量增加；对动态场景处理能力有限。
% METHOD: 2503.16710
\paragraph{4D Gaussian Splatting SLAM（动态对象的四维高斯辐射场）（2025.03）}
4D Gaussian Splatting SLAM~\cite{2503_16710} 不把动态物体简单排除，而是显式建模其时序变化。它继承 dynamic 3DGS 和 SLAM 的双重锚点：相机仍需定位，地图中还要区分静态高斯和动态高斯。方法先用 YOLOv9 与光流生成 motion masking，将高斯分为 static 和 dynamic 集合。

动态集合通过稀疏控制点和 MLP 学习时变 6DoF 变换，再用 KNN 和 linear blend skinning 插值得到密集变形场。系统还设计 2D optical flow reconstruction，将动态物体光流渲染出来监督运动学习；在线阶段则继续估计相机位姿并更新静态、动态两类高斯。相对 DG-SLAM 的“动态作为干扰物”，这篇工作的改变是把动态对象作为 4D 辐射场的一部分保留下来。
\textbf{实现细节：} 系统包含四个核心组件：(1) 运动 masking：使用YOLOv9和光流结合的方法检测动态区域，将高斯分类为静态和动态集合。(2) 变形场建模：使用沿运动区域初始化的稀疏控制点，通过MLP学习每个控制点的时变6DoF变换，再通过KNN和LBS插值获得动态高斯的密集变形场。(3) 光流监督：设计2D光流图重建算法，渲染动态物体的光流用于监督4D高斯辐射场的运动学习。(4) 增量式SLAM：在线增量估计相机姿态并更新高斯表示，同时利用光度、几何和光流约束进行优化。不将动态物体视为噪声，而是显式建模其时序变化。 提出首个在动态场景中同时进行相机定位和4D高斯辐射场重建的SLAM系统，通过分类静态/动态高斯、稀疏控制点建模变形场和2D光流重建算法，实现鲁棒跟踪和高质量渲染。 依赖YOLOv9的检测质量；光流计算增加了计算开销；对快速运动和严重遮挡的处理有待提升。
\textbf{实现补充：} 实现上还要核对输入、输出、参数、特征、初始化、训练损失、正则、渲染或推理接口，以及优化中的剪枝、分裂、采样、编码或解码步骤；这些抓手决定该方法到底改了 3DGS 管线的哪一环。
\textbf{设计动机与边界：} 相对原始 3DGS 以独立高斯和图像重建误差为主的设定，4D Gaussian Splatting SLAM 把问题落在以下可实现的改动上：系统包含四个核心组件：(1) 运动 masking：使用YOLOv9和光流结合的方法检测动态区域，将高斯分类为静态和动态集合。(2) 变形场建模：使用沿运动区域初始化的稀疏控制点，通过MLP学习每个控制点的时变6DoF变换，再通过KNN和LBS插值获得动态高斯的密集变形场。(3) 光流监督：设计2D光流图重建算法，渲染动态物体的光流用于监督4D高斯辐射场的运动学习。(4) 增量式SLAM：在线增量估计相机姿态并更新高斯表示，同时利用光度、几何和光流约束进行优化。不将动态物体视为噪声，而是显式建模其时序变化。其设计目的不是单独增加一个模块，而是用上述表示、监督或渲染约束消除原流程在目标任务中的具体失配。提出首个在动态场景中同时进行相机定位和4D高斯辐射场重建的SLAM系统，通过分类静态/动态高斯、稀疏控制点建模变形场和2D光流重建算法，实现鲁棒跟踪和高质量渲染。适用时还应保留论文给出的边界：依赖YOLOv9的检测质量；光流计算增加了计算开销；对快速运动和严重遮挡的处理有待提升。
\textbf{方法信息：} 代码：\url{https://github.com/yanyan-li/4dgs-slam}；GitHub 关注度较高，约 191 stars。



% METHOD: 2504.19938
\paragraph{Mesh-Learner（网格光栅化与球谐纹理的替代表达）（2025.04）}
Mesh-Learner~\cite{2504_19938} 严格说不是标准 3DGS SLAM，而是面向传统 rasterization pipeline 的可学习重建表示。它的局部锚点是把视角相关辐射从高斯转移到 mesh 上的 spherical harmonic textures。这样做的动机是，某些机器人或大场景系统已有网格和光栅化管线，直接使用高斯 renderer 不一定方便。

方法为每个网格分配 SH texture，通过 rasterization 找到当前帧涉及的 texels，再由 CUDA renderer 根据视角方向生成颜色。world-space EWA filtering 缓解远处混叠，自适应 SH density 根据网格大小和形状调整纹理分辨率，CPU 存储全局纹理、GPU 只加载视锥内纹理。项目实现对应 mesh-based reconstruction。相对 3DGS SLAM，它的重要性在于提供了一个相邻路线：保持可渲染学习外观，但底层几何从高斯换成网格。
\textbf{实现细节：} Mesh-Learner为每个网格分配球谐纹理(SH Texture)，利用光栅化管线确定每个帧涉及的SH纹理并进行采样。核心组件包括：(1) 混合插值方法：在SH Texel之间进行混合插值获取视角相关颜色。(2) 世界空间EWA滤波：将椭圆加权平均滤波从局部纹理空间扩展到世界空间的SH纹理上，缓解远处场景的混叠伪影。(3) 自适应SH密度策略：根据每个网格的形状和大小确定最优SH纹理分辨率。(4) 延迟渲染+CUDA渲染器：光栅化网格到屏幕空间，保存采样纹理坐标和视角方向等信息，CUDA渲染器读取这些信息进行最终渲染。(5) CPU/GPU分离存储：SH纹理存储在CPU RAM，仅传输视锥内纹理到GPU，实现超大场景训练。 提出Mesh-Learner，首个原生兼容传统光栅化管线的3D重建框架，通过在网格上学习球谐纹理实现视角相关辐射的端到端学习，支持无限大场景训练（GPU仅加载视锥内纹理），在Replica和FAST-LIVO2数据集上达到SOTA。 需要预先生成高质量网格（来自LiDAR方法）；球谐函数对极高频率视角变化的表达能力有限；训练时间可能较长。
\textbf{实现补充：} 实现上还要核对输入、输出、参数、特征、初始化、训练损失、正则、渲染或推理接口，以及优化中的剪枝、分裂、采样、编码或解码步骤；这些抓手决定该方法到底改了 3DGS 管线的哪一环。
\textbf{方法信息：} 代码：\url{https://github.com/hku-mars/mesh-learner}。
\textbf{设计动机与边界：} 相对原始 3DGS 以独立高斯和图像重建误差为主的设定，Mesh-Learner 把问题落在以下可实现的改动上：Mesh-Learner为每个网格分配球谐纹理(SH Texture)，利用光栅化管线确定每个帧涉及的SH纹理并进行采样。核心组件包括：(1) 混合插值方法：在SH Texel之间进行混合插值获取视角相关颜色。(2) 世界空间EWA滤波：将椭圆加权平均滤波从局部纹理空间扩展到世界空间的SH纹理上，缓解远处场景的混叠伪影。(3) 自适应SH密度策略：根据每个网格的形状和大小确定最优SH纹理分辨率。(4) 延迟渲染+CUDA渲染器：光栅化网格到屏幕空间，保存采样纹理坐标和视角方向等信息，CUDA渲染器读取这些信息进行最终渲染。(5) CPU/GPU分离存储：SH纹理存储在CPU RAM，仅传输视锥内纹理到GPU，实现超大场景训练。其设计目的不是单独增加一个模块，而是用上述表示、监督或渲染约束消除原流程在目标任务中的具体失配。提出Mesh-Learner，首个原生兼容传统光栅化管线的3D重建框架，通过在网格上学习球谐纹理实现视角相关辐射的端到端学习，支持无限大场景训练（GPU仅加载视锥内纹理），在Replica和FAST-LIVO2数据集上达到SOTA。适用时还应保留论文给出的边界：需要预先生成高质量网格（来自LiDAR方法）；球谐函数对极高频率视角变化的表达能力有限；训练时间可能较长。
% METHOD: 2511.23030
\paragraph{DiskChunGS（核外分块的大规模高斯 SLAM）（2025.11）}
DiskChunGS~\cite{2511_23030} 直接解决 3DGS-SLAM 的 GPU 内存瓶颈。长序列或大规模场景中，把所有高斯常驻显存会失败；DiskChunGS 的改变是 out-of-core chunk architecture。场景被划分为空间 chunks，当前相机附近或可见的活跃块加载到 GPU，不活跃块存储在磁盘上。

跟踪侧集成 ORB-SLAM3，支持 monocular、stereo 和 RGB-D，也能通过 ROS wrapper 接收外部位姿。采样侧利用 ORB-SLAM3 稀疏点云和直接基元采样，并可用 Depth Anything 2 或 Fast-ACVNet 提供深度监督。系统还设计 chunk-aware loop closure，保证大规模环境中的全局一致性。项目实现面向核外高斯地图。相对 RTG-SLAM，DiskChunGS 的关键不是单个高斯如何更紧凑，而是地图如何在磁盘和显存之间流动。
\textbf{实现细节：} 上述 Gaussian-LIC 系列和 GS-LIVO 虽然通过 CPU-GPU 双层存储缓解了显存压力，但全局地图仍须完整驻留在系统 RAM 中——当场景从单房间扩展到公里级城市街道时，高斯数量可达数百万乃至数千万，系统 RAM 容量同样成为硬上限。DiskChunGS（2025.11） 沿大规模场景建图方向做出根本性改进：借鉴操作系统虚拟内存的思想，将场景按空间划分为若干"块"（chunk），仅将当前相机视野涉及的活跃区域保留在 GPU 中，其余区域以序列化形式存储在磁盘上按需加载——与 GS-LIVO 将冷高斯留在 CPU RAM 不同，DiskChunGS 将冷数据下沉到磁盘，利用 NVMe SSD 的大容量和可扩展性彻底将可重建场景规模与任何内存上限解耦。同时，DiskChunGS 保留完整的三阶球谐各向异性高斯表示（不做简化），且前端采用 ORB-SLAM3 提供位姿估计而非 LiDAR-惯性融合，使其兼容纯视觉输入（单目、立体和 RGB-D）。技术架构围绕空间分块与 GPU-磁盘三级调度展开，包含以下核心组件： (1) 空间分块：将 3D 世界划分为大小为 $s$（实验中取 10m）的立方体块。每个高斯根据其位置 $\mathbf{p}$ 被分配到对应块： \[ \phi(\mathbf{p}) = \left\lfloor \frac{p\_i + s/2}{s} \right\rfloor, \quad i \in \{x,y,z\} \] 块坐标编码为 64 位整数——每轴分配 21 位，并加入 $2^{20}$ 的偏移量处理负坐标——从而将三维块定位转化为单次哈希查询，实现 $O(1)$ 的块查找。(2) 高斯采样策略：与 Gaussian-LIC 系列依赖 LiDAR 点云直接初始化高斯不同，DiskChunGS 面向纯视觉输入，采用直接基元采样（direct primitive sampling）替代 3DGS 原始的迭代致密化。
\textbf{实现补充：} 实现上还要核对输入、输出、参数、特征、初始化、训练损失、正则、渲染或推理接口，以及优化中的剪枝、分裂、采样、编码或解码步骤；这些抓手决定该方法到底改了 3DGS 管线的哪一环。
\textbf{方法信息：} 代码：\url{https://github.com/rootlesswater/DiskChunGS-RealSense}。
\textbf{设计动机与边界：} 相对原始 3DGS 以独立高斯和图像重建误差为主的设定，DiskChunGS 把问题落在以下可实现的改动上：DiskChunGS的核心是核外分块架构：(1)空间分块：将场景划分为空间区域(chunks)，根据当前相机位置动态加载可见块到GPU内存，不活跃块存储在磁盘上。类似虚拟内存系统的概念。(2)跟踪：集成ORB-SLAM3进行特征化视觉里程计，支持单目、立体和RGB-D输入。也可通过ROS wrapper接受外部位姿。(3)采样：使用ORB-SLAM3稀疏点云和直接基元采样方法（基于Laplacian范数估计高斯放置概率），通过Depth-Anything-2或Fast-ACVNet获取深度监督。(4)分块感知回环检测：确保大规模环境中的全局一致性。(5)多组件损失函数：图像损失(L1+SSIM)和深度损失联合优化。支持增量地图更新和跨硬件平台的序列化。其设计目的不是单独增加一个模块，而是用上述表示、监督或渲染约束消除原流程在目标任务中的具体失配。提出DiskChunGS，通过核外(out-of-core)分块架构克服3DGS-SLAM的内存瓶颈。将场景划分为空间块(chunks)，仅将活跃区域保持在GPU内存中，不活跃区域存储在磁盘上。与ORB-SLAM3集成实现位姿估计和回环检测，在KITTI全部11个序列上无内存失败完成重建，并在Jetson平台上验证了资源受限环境的部署能力。适用时还应保留论文给出的边界：块边界可能引入渲染伪影；磁盘I/O可能成为瓶颈；回环检测的分块感知设计增加了系统复杂度。
% METHOD: 2604.18047
\paragraph{GS-STVSR（2D 高斯的连续时空视频超分）（2026.04）}
GS-STVSR~\cite{2604_18047} 是 SLAM 相邻任务。它不做相机定位和三维地图，而是用 2D Gaussian Splatting 做 continuous spatio-temporal video super-resolution。放在本章的原因是，它展示了 Gaussian kernels 如何在时空视频流中被连续运动模型驱动，而不是固定在三维静态场景里。

方法把每帧表示成 2D 高斯核集合，通过光流引导运动模块预测任意时间步的高斯位置偏移和颜色；协方差重采样对齐模块利用 covariance 的时间稳定性，从先验库采样端点协方差并拟合中间状态。调整协方差尺度即可渲染任意空间放大尺度，绕开 INR 的密集网格查询。相对 4D Gaussian Splatting SLAM，它没有相机位姿和三维地图，但它对理解“高斯作为连续时空基元”有参考价值。
\textbf{方法信息：} Huawei 参与了这篇论文。
\textbf{设计动机与边界：} 相对原始 3DGS 以独立高斯和图像重建误差为主的设定，GS-STVSR 把问题落在以下可实现的改动上：GS-STVSR的核心设计：(1) 2D高斯核时空演化：将每帧表示为2D高斯核集合，通过连续运动建模驱动参数演化，协方差矩阵的尺度分量直接对应空间放大因子，渲染任意HR尺度只需调整尺度分量；(2) 协方差重采样对齐模块（CRA）：从预定义协方差先验库采样端点协方差，利用协方差的强时间稳定性进行轻量中间拟合，防止协方差漂移；(3) 光流引导运动模块：从双向光流估计任意时间步的高斯位置偏移和颜色；(4) 自适应偏移窗口：处理大规模运动。完全绕过INR方法的密集像素级网格查询。其设计目的不是单独增加一个模块，而是用上述表示、监督或渲染约束消除原流程在目标任务中的具体失配。提出GS-STVSR，首个基于2D高斯泼溅的超高效连续时空视频超分辨率框架。通过连续运动建模驱动高斯核的时空演化，完全绕过密集网格查询。利用协方差参数的时间稳定性进行轻量中间拟合，设计光流引导运动模块和协方差重采样对齐模块。在常规时间尺度（X2-X8）保持近乎恒定的推理时间，在极端尺度X32实现3倍以上加速。适用时还应保留论文给出的边界：依赖光流估计质量；对极端运动场景可能受限；协方差先验库的容量限制。
% METHOD: 2606.17935
\paragraph{MoonSplat（Sim(3) 全局优化的在线单目重建）（2026.06）}
MoonSplat~\cite{2606_17935} 面向在线单目 3DGS 重建中的长序列鲁棒性。它继承 monocular online Gaussian mapping，但指出缺少全局优化会让尺度和位姿在长序列中逐渐漂移。方法的改变是在增量维护高斯场的同时，引入 $Sim(3)$ 全局优化来修正尺度、位姿和地图一致性。

具体流程先对新关键帧做 camera tracking，再构建 voxelized Gaussians。论文和 README 都显示其实现借鉴 Scaffold-GS：每个 voxel 先由关键帧 pointmaps 的加权平均颜色初始化 anchor，再通过共享 MLP 解码成 Gaussian primitives；color residual learning 预测相对 voxel base color 的残差，使局部颜色优化更快收敛。全局一致性部分使用带回环的 factor graph 做 $Sim(3)$ 优化，同步修正相机位姿、尺度漂移和 Gaussian anchor 位置。相对 Splat-SLAM 的全局 BA，MoonSplat 更强调单目在线重建中的尺度闭环；相对 DiskChunGS 的内存分块，它解决的是长序列几何一致性而不是核外存储。
\textbf{实现细节：} MoonSplat（2026.06） 延续 online 3DGS 在"在线单目3DGS"上的基本设定，主要改变是：MoonSplat：Monocular Online Gaussian Splatting，加入Sim(3)全局优化。3DGS 的显式高斯表示和可微渲染能力保持为基础。其设计动机是：针对在线3DGS缺少全局优化导致的相机位姿脆弱和长序列效率低问题，提出鲁棒在线重建框架。具体而言，方法在单目序列中增量维护高斯场，并用Sim(3)全局优化修正尺度、位姿和地图一致性。适用边界是：单目尺度和回环仍具挑战；动态场景可能破坏全局优化假设。
\textbf{方法信息：} 代码：\url{https://github.com/TrickyGo/MoonSplat}。
\textbf{设计动机与边界：} 相对原始 3DGS 以独立高斯和图像重建误差为主的设定，MoonSplat 把问题落在以下可实现的改动上：方法在单目序列中增量维护高斯场，并用Sim(3)全局优化修正尺度、位姿和地图一致性。其设计目的不是单独增加一个模块，而是用上述表示、监督或渲染约束消除原流程在目标任务中的具体失配。针对在线3DGS缺少全局优化导致的相机位姿脆弱和长序列效率低问题，提出鲁棒在线重建框架。适用时还应保留论文给出的边界：单目尺度和回环仍具挑战；动态场景可能破坏全局优化假设。
% METHOD: 2606.24796
\paragraph{Pocket-SLAM（面向有效渲染区域的地图剪枝）（2026.06）}
Pocket-SLAM~\cite{2606_24796} 针对大规模 3DGS-SLAM 的内存效率。普通剪枝常依据单个高斯的不透明度或局部启发式，但在 SLAM 中，一个高斯是否该保留取决于它对当前和潜在渲染区域的贡献。Pocket-SLAM 的改变是以 effective rendering region contribution 作为剪枝依据。

系统在建图过程中监测高斯对可见区域和未来可能访问区域的贡献，选择性移除低贡献高斯，从而减缓地图膨胀，同时尽量保持视图合成质量。项目实现围绕 memory-aware Gaussian map maintenance。相对 MotionGS 的运动关键帧压缩，Pocket-SLAM 更直接地作用于高斯级地图维护，适合移动设备或长时间 SLAM。
\textbf{实现细节：} Pocket-SLAM（2026.06） 是记忆感知 Gaussian 地图维护路线中的新增代表。相较于 3DGS-SLAM，主要改变是：基于渲染有效区域贡献的 Gaussian pruning 策略。3DGS 的显式高斯表示和可微渲染能力保持为基础。其设计动机是：将剪枝依据从单个高斯的不透明度等局部启发式转向对有效渲染区域的贡献，从而更适合大场景 SLAM。具体而言，在 3DGS-SLAM 流程中监测高斯对当前和潜在渲染区域的贡献，选择性移除低贡献高斯，减少地图膨胀并保持视图合成质量。它给后续方法留下的比较点是：如何针对 Gaussian map 的持续增长加入渲染区域感知剪枝。边界在于：剪枝策略可能依赖轨迹覆盖和可见性估计，对未来未观测视角的保持能力需要谨慎评估。
\textbf{实现补充：} 实现上还要核对输入、输出、参数、特征、初始化、训练损失、正则、渲染或推理接口，以及优化中的剪枝、分裂、采样、编码或解码步骤；这些抓手决定该方法到底改了 3DGS 管线的哪一环。
\textbf{设计动机与边界：} 相对原始 3DGS 以独立高斯和图像重建误差为主的设定，Pocket-SLAM 把问题落在以下可实现的改动上：在 3DGS-SLAM 流程中监测高斯对当前和潜在渲染区域的贡献，选择性移除低贡献高斯，减少地图膨胀并保持视图合成质量。为了消除原始 3DGS 在该目标任务中的具体失配，方法把上述表示、监督或渲染约束作为可学习和可执行的改动，而不是单独堆叠模块。将剪枝依据从单个高斯的不透明度等局部启发式转向对有效渲染区域的贡献，从而更适合大场景 SLAM。适用时还应保留论文给出的边界：剪枝策略可能依赖轨迹覆盖和可见性估计，对未来未观测视角的保持能力需要谨慎评估。
\textbf{方法信息：} 代码：\url{https://github.com/umn-zhaolab/pocket-slam}。

\begin{takeaway}
本章方法可以按“高斯在 SLAM 中承担什么角色”来理解。第一类把高斯作为照片级地图，前端仍由 ORB、DROID、DSO、LiDAR 或 IMU 稳定位姿，例如 Photo-SLAM、DROID-Splat、GS-LIVO 和 Gaussian-LIC。第二类让高斯直接参与跟踪和建图损失，例如 Gaussian Splatting SLAM、SplaTAM、CG-SLAM、VarSplat 和 GS3LAM。第三类改地图生命周期，包括 MotionGS 的运动过滤、RGS-SLAM 的一次性初始化、LoopSplat 的高斯子图回环、DiskChunGS 的核外分块和 Pocket-SLAM 的贡献剪枝。第四类处理真实输入的不理想性：模糊、事件、动态物体、开放词汇语义、多智能体和非视觉传感器。因而，3DGS SLAM 的技术演进不是单纯把 NeRF-SLAM 换成 Gaussian renderer，而是在初始化、tracking objective、地图增删、传感器融合和全局一致性上逐步把高斯变成 SLAM 系统的核心状态。
\end{takeaway}
